% Network Design Strategies — ICT Upper Sixth — Cours signé OnBuch+
% Official MINESEC syllabus. Compiler avec : tectonic ICT04-network-design-strategies.tex
\newcommand{\DOCMATIERE}{ICT}
\newcommand{\DOCNIVEAU}{Upper Sixth}
\newcommand{\DOCMODULE}{Upper Sixth — ICT}
\newcommand{\DOCLECON}{4}
\newcommand{\DOCTITRE}{Network Design Strategies}
% OnBuch+ — préambule « cours signés » ENRICHI (compilé avec tectonic / XeLaTeX)
% Système visuel « Pop » d'OnBuch+ : fond crème (jamais blanc), bordures encre
% épaisses, ombres « dures » décalées, titres Archivo Black en capitales.
% Version MATHÉMATIQUES : graphes (pgfplots/tikz), boîtes pédagogiques
% (définition, propriété, méthode, attention, le savais-tu, exercice résolu,
% exercice, TP, bilan), page de garde et signature OnBuch+ ancrée en bas.
%
% Macros attendues dans main.tex AVANT \input{preamble} :
%   \DOCMATIERE  \DOCNIVEAU  \DOCMODULE  \DOCLECON (numéro)  \DOCTITRE
\documentclass[11pt]{article}

\usepackage{fontspec}
\usepackage[english]{babel}
\usepackage[a4paper,margin=2cm,top=3.4cm,bottom=2.8cm,headheight=1.4cm,footskip=1.3cm]{geometry}
\usepackage{amsmath,amssymb,amsfonts}
\usepackage{mathtools} % \xmapsto, \coloneqq... souvent utilisés par les modèles
\usepackage{cancel} % simplifications barrées (bilans de forces)
% \abs{z} (module, valeur absolue) et \norm{u} (norme), utilisés par les modèles
\providecommand{\abs}{}\let\abs\relax\DeclarePairedDelimiter{\abs}{\lvert}{\rvert}
\providecommand{\norm}{}\let\norm\relax\DeclarePairedDelimiter{\norm}{\lVert}{\rVert}
\usepackage[table]{xcolor} % [table] : fournit aussi \rowcolors (alternance de lignes)
\usepackage{pagecolor}
\usepackage{tikz}
\usetikzlibrary{shapes.symbols,circuits.logic.US,circuits.logic.IEC,arrows.meta,positioning,calc,shapes.geometric,shapes.misc,shapes.arrows,decorations.pathmorphing,decorations.pathreplacing,decorations.markings,patterns,shadows,mindmap,trees,fit,backgrounds,matrix,intersections,angles,quotes}
\usepackage{pgfplots}
\usepackage[version=4]{mhchem} % équations nucléaires \ce{^{235}_{92}U}
\usepackage[european,siunitx]{circuitikz} % schémas électriques
\pgfplotsset{compat=1.18}
\usepgfplotslibrary{fillbetween,groupplots} % aires entre courbes (calcul intégral)
\usepackage{enumitem}
\usepackage{fancyhdr}
% [most] charge toutes les bibliothèques tcolorbox (listings, minted, raster,
% documentation...) et gonfle inutilement la mémoire TeX sur un document long
% (~300 boîtes) : on ne charge que ce qui est réellement utilisé.
\usepackage[skins,breakable]{tcolorbox}
\usepackage{graphicx}
\usepackage{colortbl}
\usepackage{booktabs}
\usepackage{tabularx}
\usepackage{diagbox} % case d'en-tête diagonale des tableaux à double entrée
% Colonnes extensibles centrée / gauche / droite (C, L, R), souvent utilisées par les modèles
\newcolumntype{C}{>{\centering\arraybackslash}X}
\newcolumntype{L}{>{\raggedright\arraybackslash}X}
\newcolumntype{R}{>{\raggedleft\arraybackslash}X}
\usepackage{array}
\usepackage{multicol}
\usepackage{siunitx}
% parse-numbers=false : accepte aussi \SI{2,5 \times 10^{-3}}{...}
\sisetup{locale=UK,output-decimal-marker={.},group-minimum-digits=4,parse-numbers=false}
\DeclareSIUnit\ohm{\ensuremath{\symup{\Omega}}} % U+2126 absent de la police
\DeclareSIUnit\division{div} % réglages d'oscilloscope (ms/div, V/div)
\DeclareSIUnit\year{yr}\DeclareSIUnit\annum{yr}\DeclareSIUnit\Ma{Ma}\DeclareSIUnit\tr{tr} % tours (vitesse de rotation, ex. \SI{1500}{\tr\per\minute})
\usepackage{qrcode}
% \degree : commande fréquemment utilisée par les modèles pour un angle ou une
% température en dehors de \SI{}{} (non fournie par les paquets chargés).
\providecommand{\degree}{^{\circ}}
% \qed (fin de démonstration), utilisé par les modèles sans amsthm
\providecommand{\qed}{\hfill$\square$}
% \barre : double barre des valeurs interdites dans un tableau de variations (tkz-tab)
\providecommand{\barre}{\ensuremath{\|}}
% environnement proof (amsthm non chargé) utilisé par les modèles
\ifdefined\proof\else\newenvironment{proof}[1][Proof]{\par\noindent\textit{#1.}\ }{\qed\par}\fi
\usepackage{lastpage}
\usepackage{hyperref}
\hypersetup{hidelinks}

% ============================================================
% Palette « Pop » OnBuch+ — mêmes hex que la classe P (plus_ui.dart)
% ============================================================
\definecolor{poporange}{HTML}{F59321}
\definecolor{popdark}{HTML}{E07A0C}
\definecolor{popcream}{HTML}{FFF6E8}
\definecolor{popink}{HTML}{1C1714}
\definecolor{poppurple}{HTML}{7A5AE0}
\definecolor{popgreen}{HTML}{2E9E6A}
\definecolor{poppink}{HTML}{E0457B}
\definecolor{popblue}{HTML}{2F7DE1}
\definecolor{popgold}{HTML}{C98A12}
\definecolor{popmuted}{HTML}{8A7F76}
\definecolor{popline}{HTML}{EADFD2}
\definecolor{popgray}{HTML}{8A7F76}
\colorlet{popgrey}{popgray}
% Alias tolérants (le modèle nomme parfois les couleurs différemment)
\colorlet{popwhite}{white}
\colorlet{popblack}{popink}
\colorlet{popred}{poppink}
\colorlet{popyellow}{popgold}
% Teintes claires pour fonds de tableaux / zones
\colorlet{popblueL}{popblue!10!white}
\colorlet{poporangeL}{poporange!14!white}
\colorlet{popgreenL}{popgreen!12!white}
\colorlet{poppurpleL}{poppurple!12!white}
\colorlet{poppinkL}{poppink!12!white}
\colorlet{popgoldL}{popgold!14!white}

\pagecolor{popcream}

% ============================================================
% Typographie
% ============================================================
\defaultfontfeatures{Ligatures=TeX}
\setmainfont{PlusJakartaSans-Medium.ttf}[Path=../../../fonts/,BoldFont=PlusJakartaSans-Bold.ttf]
% Maths en sans-serif (Fira Math) avec chiffres et lettres droites de Plus
% Jakarta, pour que les formules se fondent dans le texte.
\usepackage[math-style=ISO,bold-style=ISO]{unicode-math}
\setmathfont{FiraMath-Regular.otf}[Path=../../../fonts/]
\setmathfont{PlusJakartaSans-Medium.ttf}[Path=../../../fonts/,range={up/num,up/{Latin,latin},\mathup}]
% (icomma retiré : décimales avec point en anglais)
% Caractères mathématiques Unicode absents des polices de texte (Plus Jakarta,
% Archivo Black) : rendus par leur équivalent mathématique, en texte comme en
% formule — sinon ils s'affichent en carré vide (ex. « Arithmétique dans ℤ »).
\usepackage{newunicodechar}
\newunicodechar{ℝ}{\ensuremath{\mathbb{R}}}
\newunicodechar{ℤ}{\ensuremath{\mathbb{Z}}}
\newunicodechar{ℂ}{\ensuremath{\mathbb{C}}}
\newunicodechar{ℕ}{\ensuremath{\mathbb{N}}}
\newunicodechar{ℚ}{\ensuremath{\mathbb{Q}}}
\newunicodechar{𝔻}{\ensuremath{\mathbb{D}}}
\newunicodechar{α}{\ensuremath{\alpha}}
\newunicodechar{β}{\ensuremath{\beta}}
\newunicodechar{γ}{\ensuremath{\gamma}}
\newunicodechar{δ}{\ensuremath{\delta}}
\newunicodechar{ε}{\ensuremath{\varepsilon}}
\newunicodechar{θ}{\ensuremath{\theta}}
\newunicodechar{λ}{\ensuremath{\lambda}}
\newunicodechar{μ}{\ensuremath{\mu}}
\newunicodechar{σ}{\ensuremath{\sigma}}
\newunicodechar{τ}{\ensuremath{\tau}}
\newunicodechar{φ}{\ensuremath{\varphi}}
\newunicodechar{ψ}{\ensuremath{\psi}}
\newunicodechar{ω}{\ensuremath{\omega}}
\newunicodechar{Σ}{\ensuremath{\Sigma}}
% majuscules produites par \MakeUppercase dans les titres (α-aminés -> Α)
\newunicodechar{Α}{\ensuremath{\alpha}}
\newunicodechar{Β}{\ensuremath{\beta}}
\newunicodechar{Γ}{\ensuremath{\Gamma}}
\newunicodechar{Δ}{\ensuremath{\Delta}}
\newunicodechar{Ε}{\ensuremath{\varepsilon}}
\newunicodechar{Θ}{\ensuremath{\Theta}}
\newunicodechar{Λ}{\ensuremath{\Lambda}}
\newunicodechar{Μ}{\ensuremath{\mu}}
\newunicodechar{Τ}{\ensuremath{\tau}}
\newunicodechar{Φ}{\ensuremath{\Phi}}
\newunicodechar{Ψ}{\ensuremath{\Psi}}
\newunicodechar{Ω}{\ensuremath{\Omega}}
\newunicodechar{↦}{\ensuremath{\mapsto}}
\newunicodechar{∈}{\ensuremath{\in}}
\newunicodechar{∉}{\ensuremath{\notin}}
\newunicodechar{∧}{\ensuremath{\wedge}}
\newunicodechar{⇔}{\ensuremath{\Leftrightarrow}}
\newunicodechar{⟺}{\ensuremath{\Longleftrightarrow}}
\newunicodechar{⇒}{\ensuremath{\Rightarrow}}
\newunicodechar{⟹}{\ensuremath{\Longrightarrow}}
\newunicodechar{∘}{\ensuremath{\circ}}
\newunicodechar{∪}{\ensuremath{\cup}}
\newunicodechar{∩}{\ensuremath{\cap}}
\newunicodechar{≡}{\ensuremath{\equiv}}
\newunicodechar{⊂}{\ensuremath{\subset}}
\newunicodechar{⊥}{\ensuremath{\perp}}
\newunicodechar{∃}{\ensuremath{\exists}}
\newunicodechar{∀}{\ensuremath{\forall}}
\newunicodechar{∛}{\ensuremath{\sqrt[3]{\,}}}
\newunicodechar{∜}{\ensuremath{\sqrt[4]{\,}}}
\newunicodechar{⃗}{\ensuremath{{}^{\to}}}
\newunicodechar{ⁿ}{\ifmmode^{n}\else\textsuperscript{\ensuremath{n}}\fi}
\newunicodechar{ˣ}{\ifmmode^{x}\else\textsuperscript{\ensuremath{x}}\fi}
\newunicodechar{ᵇ}{\ifmmode^{b}\else\textsuperscript{\ensuremath{b}}\fi}
\newunicodechar{ʳ}{\ifmmode^{r}\else\textsuperscript{\ensuremath{r}}\fi}
\newunicodechar{ᵘ}{\ifmmode^{u}\else\textsuperscript{\ensuremath{u}}\fi}
\newunicodechar{ʸ}{\ifmmode^{y}\else\textsuperscript{\ensuremath{y}}\fi}
\newunicodechar{ᵃ}{\ifmmode^{a}\else\textsuperscript{\ensuremath{a}}\fi}
\newunicodechar{ᵉ}{\ifmmode^{e}\else\textsuperscript{\ensuremath{e}}\fi}
\newunicodechar{ᵏ}{\ifmmode^{k}\else\textsuperscript{\ensuremath{k}}\fi}
\newunicodechar{ᵖ}{\ifmmode^{p}\else\textsuperscript{\ensuremath{p}}\fi}
\newunicodechar{ᴺ}{\ifmmode^{N}\else\textsuperscript{\ensuremath{N}}\fi}
\newunicodechar{ᶜ}{\ifmmode^{c}\else\textsuperscript{\ensuremath{c}}\fi}
\newunicodechar{ʰ}{\ifmmode^{h}\else\textsuperscript{\ensuremath{h}}\fi}
\newunicodechar{ᵠ}{\ifmmode^{q}\else\textsuperscript{\ensuremath{q}}\fi}
\newunicodechar{ₙ}{\ifmmode_{n}\else\textsubscript{\ensuremath{n}}\fi}
\newunicodechar{ₐ}{\ifmmode_{a}\else\textsubscript{\ensuremath{a}}\fi}
\newunicodechar{ᵢ}{\ifmmode_{i}\else\textsubscript{\ensuremath{i}}\fi}
\newunicodechar{ᵣ}{\ifmmode_{r}\else\textsubscript{\ensuremath{r}}\fi}
\newunicodechar{ₖ}{\ifmmode_{k}\else\textsubscript{\ensuremath{k}}\fi}
\newunicodechar{ᵦ}{\ifmmode_{\beta}\else\textsubscript{\ensuremath{\beta}}\fi}
\newunicodechar{ₓ}{\ifmmode_{x}\else\textsubscript{\ensuremath{x}}\fi}
\newfontfamily\popdisplay{ArchivoBlack-Regular.ttf}[Path=../../../fonts/]
\newfontfamily\poplogo{SpaceGrotesk-Bold.ttf}[Path=../../../fonts/]
\newfontfamily\popsemi{PlusJakartaSans-SemiBold.ttf}[Path=../../../fonts/]
\newfontfamily\popxbold{PlusJakartaSans-ExtraBold.ttf}[Path=../../../fonts/]
\newfontfamily\popxboldls{PlusJakartaSans-ExtraBold.ttf}[Path=../../../fonts/,LetterSpace=3.0]

\setlist[enumerate]{leftmargin=1.7em,itemsep=5pt,topsep=4pt}
\setlist[itemize]{leftmargin=1.7em,itemsep=4pt,topsep=4pt,label={\color{poporange}\textbullet}}
\setlength{\parindent}{0pt}
\setlength{\parskip}{5pt}
\renewcommand{\arraystretch}{1.25}

% pgfplots : style Pop par défaut
\pgfplotsset{
  every axis/.append style={axis line style={popink,line width=1pt},tick style={popink},
    grid style={popline,line width=0.5pt},label style={font=\small},tick label style={font=\footnotesize}},
  popcurve/.style={poporange,line width=1.8pt,no markers,smooth},
}
% Même style disponible dans un \draw TikZ simple (hors pgfplots)
\tikzset{popcurve/.style={poporange,line width=1.8pt}}
\tikzset{popgrid/.style={popline,very thin}}
\colorlet{popcurve}{poporange} % le nom du style est parfois utilisé comme couleur

% ============================================================
% Pill / squiggle
% ============================================================
\newcommand{\poppill}[3]{%
  \tikz[baseline=-0.55ex]\node[draw=popink,line width=1.2pt,rounded corners=2.6pt,%
    fill=#2,inner xsep=6.5pt,inner ysep=3pt]{%
    \popxboldls\fontsize{7}{7}\selectfont\color{#3}\MakeUppercase{#1}};%
}
\newcommand{\popsquiggle}[1]{%
  \tikz[baseline=-0.4ex]{
    \draw[#1,line width=2.6pt,line cap=round]
      plot [smooth] coordinates {(0,0.09) (0.34,0.01) (0.68,0.21) (1.02,0.01) (1.36,0.10)};
  }%
}
% Mot-clé surligné (marqueur orange sous le texte)
\newcommand{\cle}[1]{{\popxbold\color{popdark}#1}}

% ============================================================
% En-tête / pied
% ============================================================
\pagestyle{fancy}
\fancyhf{}
\renewcommand{\headrulewidth}{0pt}
\renewcommand{\footrulewidth}{0pt}
\lhead{\raisebox{-0.15ex}{\poplogo\fontsize{13}{13}\selectfont\color{popink}OnBuch\color{poporange}+}}
\rhead{\poppill{\DOCMATIERE\ \textbullet\ \DOCNIVEAU\ \textbullet\ Chapter \DOCLECON}{white}{popink}}
\lfoot{\popxbold\fontsize{7.6}{7.6}\selectfont\color{popmuted}\MakeUppercase{Signed course OnBuch+}\ \textbullet\ {\popsemi\DOCTITRE}}
\rfoot{\tikz[baseline=-0.6ex]\node[rounded corners=8.5pt,fill=popink,minimum height=17pt,minimum width=17pt,inner xsep=5pt,inner sep=0]{\popxbold\fontsize{8}{8}\selectfont\color{popcream}\thepage\,/\,\pageref*{LastPage}};}

% ============================================================
% Titres
% ============================================================
\newcommand{\chaptitre}[2]{%
  \begin{tcolorbox}[enhanced,colback=poporange,colframe=popink,boxrule=2.6pt,arc=11pt,
      drop shadow=popink,left=15pt,right=15pt,top=12pt,bottom=11pt,before skip=3pt,after skip=18pt]
    \poppill{#2}{popink}{popcream}\par\vspace{9pt}
    {\raggedright\hyphenpenalty=10000\exhyphenpenalty=10000\popdisplay\fontsize{22}{24}\selectfont\color{popcream}\MakeUppercase{#1}\par}
  \end{tcolorbox}
}
% Section numérotée : pastille orange avec le numéro + titre Archivo
\newcounter{popsec}
\newcommand{\coursec}[1]{%
  \stepcounter{popsec}\setcounter{popsub}{0}%
  \par\vspace{16pt}\noindent
  \begin{minipage}{\linewidth}
  \tikz[baseline=-0.7ex]\node[draw=popink,line width=1.6pt,fill=poporange,rounded corners=4pt,minimum size=22pt,inner sep=0]{\popdisplay\fontsize{12}{12}\selectfont\color{popcream}\thepopsec};%
  \hspace{8pt}{\popdisplay\fontsize{15}{17}\selectfont\color{popink}\MakeUppercase{#1}}\par
  \vspace{4pt}\noindent{\color{poporange}\rule{36pt}{3.6pt}}
  \end{minipage}\par\nopagebreak\vspace{8pt}
}
\newcounter{popsub}
\newcommand{\courssub}[1]{%
  \stepcounter{popsub}%
  \par\vspace{9pt}\noindent{\popxbold\fontsize{11.5}{13}\selectfont\color{popink}{\color{poporange}\thepopsec.\thepopsub}\hspace{6pt}#1}\par\nopagebreak\vspace{4pt}
}

% ============================================================
% Boîtes pédagogiques — même recette : fond blanc, bordure épaisse colorée,
% ombre dure encre, badge plein. Titre optionnel [#1].
% ============================================================
\tcbset{popbase/.style={breakable,enhanced,colback=white,boxrule=2.3pt,arc=9pt,
  shadow={3pt}{-3pt}{0pt}{popink},left=14pt,right=14pt,top=11pt,bottom=11pt,before skip=11pt,after skip=15pt,
  fontupper=\popsemi\fontsize{10.6}{14.5}\selectfont\color{popink},
  fontlower=\popsemi\fontsize{10.6}{14.5}\selectfont\color{popink}}}
% Coche dessinée (le glyphe ✓ n'existe pas dans les polices du document)
\newcommand{\popcheck}[1]{\tikz[baseline=-0.5ex]\draw[#1,line width=1.6pt,line cap=round,line join=round](-0.13,0.0)--(-0.04,-0.09)--(0.13,0.1);}
\providecommand{\coche}{\popcheck{popgreen}}
\providecommand{\resultat}[1]{\par\smallskip\noindent\fcolorbox{popgreen}{popgreenL}{\parbox{\dimexpr\linewidth-2\fboxsep-2\fboxrule\relax}{\textbf{Result.} #1}}\par\smallskip}
\newcommand{\popboxhead}[3]{\poppill{#1}{#2}{white}\ifx\relax#3\relax\else\hspace{6pt}{\popxbold\fontsize{10.6}{12}\selectfont\color{popink}#3}\fi\par\vspace{7pt}}

\newtcolorbox{aretenir}[1][]{popbase,colframe=popgreen,before upper={\popboxhead{Key points}{popgreen}{#1}}}
\newtcolorbox{exemplebox}[1][]{popbase,colframe=poporange,before upper={\popboxhead{Exemple}{poporange}{#1}}}
\newtcolorbox{definition}[1][]{popbase,colframe=popblue,before upper={\popboxhead{Definition}{popblue}{#1}}}
\newtcolorbox{propriete}[1][]{popbase,colframe=poppurple,before upper={\popboxhead{Property}{poppurple}{#1}}}
\newtcolorbox{methode}[1][]{popbase,colframe=popgold,before upper={\popboxhead{Method}{popgold}{#1}}}
\newtcolorbox{attention}[1][]{popbase,colframe=poppink,before upper={\popboxhead{Warning}{poppink}{#1}}}
\newtcolorbox{experience}[1][]{popbase,colframe=popblue,colback=popblueL,before upper={\popboxhead{Experiment}{popblue}{#1}}}
% « Le savais-tu ? » : carte violette pleine (contexte, culture, Cameroun)
\newtcolorbox{savaistu}[1][]{popbase,colback=poppurple,colframe=popink,
  fontupper=\popsemi\fontsize{10.4}{14.2}\selectfont\color{white},
  before upper={\poppill{Did you know?}{white}{poppurple}\ifx\relax#1\relax\else\hspace{6pt}{\popxbold\color{white}#1}\fi\par\vspace{7pt}}}
% Exercice résolu : énoncé en haut, \tcblower puis la solution sur fond crème
\newcounter{popexo}\newcounter{popexr}
\newtcolorbox{exoresolu}[1][]{popbase,colframe=poporange,colbacklower=popcream,
  segmentation style={popink,line width=1.2pt,dashed},
  before upper={\stepcounter{popexr}\popboxhead{Worked example \thepopexr}{poporange}{#1}},
  before lower={\poppill{Solution}{popink}{popcream}\par\vspace{6pt}}}
% Exercice (non résolu en place) — la correction est dans la partie Corrigés
\newtcolorbox{exercice}[1][]{popbase,colframe=popink,
  before upper={\stepcounter{popexo}\popboxhead{Exercise \thepopexo}{popink}{#1}}}
% Corrigé
\newtcolorbox{corrige}[1][]{popbase,colframe=popgreen,colback=popgreenL,
  before upper={\popboxhead{Solution}{popgreen}{#1}}}
% Objectifs / prérequis / bilan
\newtcolorbox{objectifs}{popbase,colframe=popink,colback=poporangeL,
  before upper={\popboxhead{Chapter objectives}{popink}{}}}
\newtcolorbox{prerequis}{popbase,colframe=popink,colback=popblueL,
  before upper={\popboxhead{Prerequisites}{popblue}{}}}
\newtcolorbox{bilan}{popbase,colframe=popink,colback=white,boxrule=2.8pt,
  before upper={\popboxhead{Fiche bilan}{poporange}{L'essentiel à connaître pour l'examen}}}
% Figure encadrée avec légende
\newtcolorbox{popfigure}[1][]{enhanced,breakable=false,colback=white,colframe=popline,boxrule=1.4pt,arc=8pt,
  left=10pt,right=10pt,top=10pt,bottom=8pt,before skip=10pt,after skip=12pt,halign=center,
  lower separated=false,halign lower=center,
  fontlower=\popsemi\fontsize{9}{11.5}\selectfont\color{popmuted},
  #1}
\newcommand{\legende}[1]{\par\vspace{6pt}{\popsemi\fontsize{9}{11.5}\selectfont\color{popmuted}#1}}

% ============================================================
% Signature OnBuch+
% ============================================================
\newcommand{\poplogoinline}[1]{{\poplogo\fontsize{#1}{#1}\selectfont\color{popink}OnBuch\color{poporange}+}}
\newcommand{\signatureonbuch}{%
  \begin{tcolorbox}[enhanced,colback=popink,colframe=popink,boxrule=2.2pt,arc=10pt,
      drop shadow=poporange,left=14pt,right=14pt,top=10pt,bottom=10pt,before skip=0pt,after skip=0pt]
    \tikz[baseline=-0.7ex]\node[circle,fill=popgreen,draw=popcream,line width=1.3pt,minimum size=22pt,inner sep=0]{\popcheck{white}};%
    \hspace{9pt}\begin{minipage}[c]{0.62\linewidth}
      {\popxbold\fontsize{12}{13}\selectfont\color{popcream}Signed\ \textbullet\ {\poplogo OnBuch\color{poporange}+}}\par\vspace{2pt}
      {\raggedright\popsemi\fontsize{8}{10.5}\selectfont\color{popline}\MakeUppercase{OnBuch+ teaching team}\\\MakeUppercase{Follows the official MINESEC syllabus}\par}
    \end{minipage}\hfill
    {\poplogo\fontsize{22}{22}\selectfont\color{popcream}OnBuch\color{poporange}+}
  \end{tcolorbox}%
}

% ============================================================
% Page de garde
% #1 titre · #2 matière · #3 niveau · #4 module · #5 n° leçon · #6 durée ·
% #7 description / objectifs (texte ou itemize)
% ============================================================
\newcommand{\pagegarde}[7]{%
  \thispagestyle{empty}
  \newgeometry{a4paper,margin=1.7cm,top=1.6cm,bottom=1.4cm}%
  \begin{tikzpicture}[remember picture,overlay]
    \fill[poporange!18] ([xshift=-3.2cm,yshift=-3.6cm]current page.north east) circle (4.6cm);
    \fill[poppurple!14] ([xshift=2.4cm,yshift=7.6cm]current page.south west) circle (3.8cm);
  \end{tikzpicture}%
  {\poplogo\fontsize{30}{30}\selectfont\color{popink}OnBuch\color{poporange}+}\hfill
  \raisebox{0.6ex}{\poppill{Signed course \textbullet\ Premium}{popink}{popcream}}\par
  \vspace{1pt}\popsquiggle{poppurple}\par
  \vspace{8pt}
  \begin{tcolorbox}[enhanced,colback=poporange,colframe=popink,boxrule=2.8pt,arc=13pt,
      drop shadow=popink,left=18pt,right=18pt,top=12pt,bottom=13pt,after skip=10pt]
    \poppill{#2 \textbullet\ #3}{popink}{popcream}\hspace{5pt}\poppill{#4}{white}{popink}\par\vspace{8pt}
    {\popdisplay\fontsize{14}{14}\selectfont\color{popink}CHAPTER #5}\par\vspace{4pt}
    {\raggedright\hyphenpenalty=10000\exhyphenpenalty=10000\popdisplay\fontsize{25}{27}\selectfont\color{popcream}\MakeUppercase{#1}\par}
    \vspace{8pt}
    {\popxbold\fontsize{9.5}{9.5}\selectfont\color{popink}\MakeUppercase{Suggested time: #6}}
  \end{tcolorbox}
  \begin{tcolorbox}[enhanced,colback=white,colframe=popink,boxrule=2.2pt,arc=10pt,
      drop shadow=popink,left=16pt,right=16pt,top=10pt,bottom=10pt,after skip=8pt]
    \poppill{In this chapter}{poppurple}{white}\par\vspace{5pt}
    \popsemi\fontsize{9.3}{12.6}\selectfont\color{popink}#7
  \end{tcolorbox}
  \noindent\begin{minipage}[t]{0.487\linewidth}
    \begin{tcolorbox}[enhanced,colback=popgreen,colframe=popink,boxrule=2.2pt,arc=10pt,
        drop shadow=popink,left=11pt,right=11pt,top=10pt,bottom=10pt,height=3.1cm,valign=center]
      \tcbox[colback=white,colframe=popink,boxrule=1.3pt,arc=5pt,boxsep=0pt,nobeforeafter,
        left=3pt,right=3pt,top=3pt,bottom=3pt,valign=center]{\qrcode[height=1.9cm]{https://play.google.com/store/apps/details?id=com.luvvix.onbuch&hl=en}}%
      \hspace{8pt}\begin{minipage}[c]{0.5\linewidth}
        {\popdisplay\fontsize{11}{12.5}\selectfont\color{white}\MakeUppercase{Download OnBuch}}\par\vspace{4pt}
        {\raggedright\popsemi\fontsize{8.4}{11}\selectfont\color{white}Free \textbullet\ Google Play\\AI tutor, past papers, results\par}
      \end{minipage}
    \end{tcolorbox}
  \end{minipage}\hfill
  \begin{minipage}[t]{0.487\linewidth}
    \begin{tcolorbox}[enhanced,colback=poppurple,colframe=popink,boxrule=2.2pt,arc=10pt,
        drop shadow=popink,left=11pt,right=11pt,top=10pt,bottom=10pt,height=3.1cm,valign=center]
      \tikz[baseline=-0.7ex]\node[circle,fill=white,draw=popink,line width=1.3pt,minimum size=1.6cm,inner sep=0]{\popdisplay\fontsize{24}{24}\selectfont\color{poppurple}+};%
      \hspace{8pt}\begin{minipage}[c]{0.6\linewidth}
        {\popdisplay\fontsize{11}{12.5}\selectfont\color{white}\MakeUppercase{Go OnBuch+}}\par\vspace{4pt}
        {\raggedright\popsemi\fontsize{8.4}{11}\selectfont\color{white}All signed courses, unlimited AI tutor, PDF sheets — OnBuch+ tab in the app\par}
      \end{minipage}
    \end{tcolorbox}
  \end{minipage}
  \vspace{8pt}
  \popcontenu
  \vfill
  \signatureonbuch
  \restoregeometry
  \newpage
}

% Rangée « Ce cours contient » (page de garde)
\newcommand{\poptile}[3]{% #1 couleur #2 gros texte #3 légende
  \begin{tcolorbox}[enhanced,colback=#1,colframe=popink,boxrule=2pt,arc=9pt,shadow={2.5pt}{-2.5pt}{0pt}{popink},
      left=6pt,right=6pt,top=9pt,bottom=9pt,halign=center,valign=center,height=1.75cm,width=\linewidth,nobeforeafter]
    {\popdisplay\fontsize{12.5}{13}\selectfont\color{white}\MakeUppercase{#2}}\par\vspace{3pt}
    {\centering\popsemi\fontsize{7.8}{9.5}\selectfont\color{white}#3\par}
  \end{tcolorbox}}
\newcommand{\popcontenu}{%
  \par\noindent{\popxboldls\fontsize{7.5}{7.5}\selectfont\color{popmuted}THIS SIGNED COURSE CONTAINS}\par\vspace{7pt}
  \noindent\begin{minipage}[t]{0.235\linewidth}\poptile{popblue}{Course}{complete, illustrated, step by step}\end{minipage}\hfill
  \begin{minipage}[t]{0.235\linewidth}\poptile{poporange}{Methods}{and worked examples}\end{minipage}\hfill
  \begin{minipage}[t]{0.235\linewidth}\poptile{poppink}{Exercises}{exam style + solutions}\end{minipage}\hfill
  \begin{minipage}[t]{0.235\linewidth}\poptile{popgreen}{Summary sheet}{the essentials to remember}\end{minipage}\par}

% Sommaire
\newtcolorbox{sommaire}{enhanced,colback=white,colframe=popink,boxrule=2.2pt,arc=10pt,
  drop shadow=popink,left=18pt,right=18pt,top=13pt,bottom=13pt,before skip=6pt,after skip=20pt,
  before upper={\poppill{Contents}{poppurple}{white}\par\vspace{9pt}\popsemi\fontsize{10.6}{14.5}\selectfont\color{popink}}}

% Signature de fin de cours — poussée tout en bas de la dernière page
\newcommand{\findecours}{%
  \par\vfill\nopagebreak
  \begin{center}\popsquiggle{poporange}\end{center}\vspace{4pt}
  \signatureonbuch
}
\AtBeginDocument{\def\llbracket{\mathopen{[\mkern-3.5mu[}}\def\rrbracket{\mathclose{]\mkern-3.5mu]}}} % intervalles d'entiers (glyphe ⟦ absent de la police)
% Glyphes absents de Fira Math : redéfinis à partir de glyphes disponibles
\AtBeginDocument{%
  \def\setminus{\mathbin{\backslash}}\let\smallsetminus\setminus
  \def\vdots{\mathord{\vbox{\baselineskip3.2pt\lineskiplimit0pt\kern4pt\hbox{.}\hbox{.}\hbox{.}}}}%
  \def\ddots{\mathinner{\mkern1mu\raise6.5pt\hbox{.}\mkern2mu\raise3.6pt\hbox{.}\mkern2mu\raise0.7pt\hbox{.}\mkern1mu}}%
  \def\triangle{\mathord{\scalebox{0.9}{$\symup{\Delta}$}}}%
  \def\nabla{\mathord{\raisebox{\depth}{\scalebox{1}[-1]{$\symup{\Delta}$}}}}%
  \def\checkmark{\popcheck{popdark}}%
  \def\star{\mathbin{\ast}}\def\bigstar{\mathord{\ast}}%
  \def\diamond{\mathbin{\scalebox{0.6}{$\Diamond$}}}%
  \def\bigcup{\mathop{\mathchoice{\raisebox{-0.3em}{\scalebox{1.7}{$\cup$}}}{\scalebox{1.25}{$\cup$}}{\cup}{\cup}}\displaylimits}%
  \def\bigcap{\mathop{\mathchoice{\raisebox{-0.3em}{\scalebox{1.7}{$\cap$}}}{\scalebox{1.25}{$\cap$}}{\cap}{\cap}}\displaylimits}%
  \def\lesseqqgtr{\mathrel{\lessgtr}}\def\bigoplus{\mathop{\oplus}\displaylimits}%
}
\providecommand{\ssi}{if and only if} % abréviation fréquente
\AtBeginDocument{\providecommand{\wideparen}{\overparen}} % arc de cercle

% Fonctions usuelles que les modèles utilisent sans que amsmath les fournisse
\providecommand{\arsinh}{\operatorname{arsinh}}\providecommand{\arcosh}{\operatorname{arcosh}}
\providecommand{\artanh}{\operatorname{artanh}}\providecommand{\arsech}{\operatorname{arsech}}
\providecommand{\arcsch}{\operatorname{arcsch}}\providecommand{\arcoth}{\operatorname{arcoth}}
\providecommand{\arcsinh}{\operatorname{arsinh}}\providecommand{\arccosh}{\operatorname{arcosh}}
\providecommand{\arctanh}{\operatorname{artanh}}\providecommand{\sech}{\operatorname{sech}}
\providecommand{\csch}{\operatorname{csch}}\providecommand{\arcsec}{\operatorname{arcsec}}
\providecommand{\arccsc}{\operatorname{arccsc}}\providecommand{\arccot}{\operatorname{arccot}}
\providecommand{\sgn}{\operatorname{sgn}}\providecommand{\lcm}{\operatorname{lcm}}
\providecommand{\Var}{\operatorname{Var}}\providecommand{\Cov}{\operatorname{Cov}}
\providecommand{\permil}{\ensuremath{{}^{0}\!/_{00}}}\providecommand{\textperthousand}{\ensuremath{{}^{0}\!/_{00}}}

%% FIGFIX-BEGIN v5 — correctifs globaux des figures (docs/onbuch-generation/tools/figfix)
\usepackage{adjustbox}\usepackage{etoolbox}\tcbuselibrary{raster}
% toute figure TikZ/pgfplots plus large que la ligne est réduite à la largeur disponible
\BeforeBeginEnvironment{tikzpicture}{\begin{adjustbox}{max width=\linewidth}}
\AfterEndEnvironment{tikzpicture}{\end{adjustbox}}
% légendes pgfplots lisibles par-dessus les courbes
\pgfplotsset{every axis/.append style={legend style={fill=white,fill opacity=0.92,text opacity=1,draw=black!15,inner sep=3pt}}}
% étiquettes d'arêtes (syntaxe edge["..."]) avec fond blanc
\tikzset{every edge quotes/.append style={fill=white,inner sep=1.2pt}}
%% FIGFIX-END
\begin{document}

\pagegarde{Network Design Strategies}{ICT}{Upper Sixth}{Upper Sixth — ICT}{4}{1 periods}{This chapter introduces the fundamental strategies for designing computer networks, focusing on the two major architectures: peer-to-peer and client/server. Learners will compare their structures, advantages, limitations and costs, and will be able to recommend an appropriate architecture for real-life organisations in Cameroon.
\vspace{4pt}
\begin{multicols}{2}\small
\begin{itemize}
\item By the end of this chapter I can define network architecture and explain its role in network design.
\item By the end of this chapter I can describe the structure and functioning of a peer-to-peer network.
\item By the end of this chapter I can describe the structure and functioning of a client/server network.
\item By the end of this chapter I can identify the roles of servers, clients and peers in a network.
\item By the end of this chapter I can compare peer-to-peer and client/server architectures using criteria such as cost, security, scalability and administration.
\item By the end of this chapter I can state the advantages and disadvantages of each architecture.
\end{itemize}\end{multicols}}

\begin{sommaire}
\begin{multicols}{2}
\begin{enumerate}
\item Introduction to Network Architecture
\item Peer-to-Peer Architecture
\item Client/Server Architecture
\item Comparing Architectures and Choosing a Design Strategy
\item Integration activity
\item Methods and worked examples
\item Exercises
\item Solutions to the exercises
\item Summary sheet
\end{enumerate}
\end{multicols}
\end{sommaire}

\begin{objectifs}
\begin{itemize}
\item By the end of this chapter I can define network architecture and explain its role in network design.
\item By the end of this chapter I can describe the structure and functioning of a peer-to-peer network.
\item By the end of this chapter I can describe the structure and functioning of a client/server network.
\item By the end of this chapter I can identify the roles of servers, clients and peers in a network.
\item By the end of this chapter I can compare peer-to-peer and client/server architectures using criteria such as cost, security, scalability and administration.
\item By the end of this chapter I can state the advantages and disadvantages of each architecture.
\item By the end of this chapter I can recommend a suitable architecture for a given organisation (school, cybercafé, bank, ministry).
\item By the end of this chapter I can draw and label simple diagrams of both architectures.
\item By the end of this chapter I can answer GCE-style structured and essay questions on network architectures.
\end{itemize}
\end{objectifs}

\begin{prerequis}
\begin{itemize}
\item Definition and advantages of a computer network (Lower Sixth)
\item Network hardware: computers, cables, switches, hubs, NICs (Lower Sixth)
\item Network topologies: bus, star, ring (Lower Sixth)
\item Types of networks: LAN, MAN, WAN (Lower Sixth)
\item Basic notions of resource sharing (files, printers, internet)
\end{itemize}
\end{prerequis}

\newpage

\coursec{Introduction to Network Architecture}

At a busy cybercafé in Bamenda, ten computers share a single printer and one internet connection, yet the owner complains that everything slows down whenever customers stream videos. Meanwhile, the nearby microfinance bank runs its own network where one powerful machine controls all the accounts data, and tellers cannot work when it fails. Why do two organisations choose such different ways of organising their computers? The answer lies in \cle{network architecture} --- the strategic design decision that determines cost, speed, security and reliability. In this chapter, you will learn to analyse, compare and choose between peer-to-peer and client/server architectures like a real network engineer.

\courssub{What is Network Architecture?}

\textbf{From cables to strategy.}\par
When students first meet computer networks, they usually think of the visible things: cables, switches, Wi-Fi routers and computers on desks. But two networks can use exactly the same cables and the same switch, yet behave completely differently. In one, every computer is equal and shares its own files freely; in the other, one powerful machine is in charge and all the others depend on it. The difference is not in the hardware --- it is in the \emph{organisation of roles}. That organisation is what we call the architecture.

Think of a secondary school. The \emph{buildings and corridors} are like the cables and devices. But the way the school is \emph{run} --- a principal who takes decisions, teachers who deliver lessons, students who receive them --- is the organisation. Two schools may have identical buildings but very different organisations. In the same way, two networks may have identical hardware but very different architectures.

\begin{definition}[Network Architecture]
The \cle{network architecture} of a computer network is the overall design and organisation of the computers and resources in the network: it specifies the \textbf{role} played by each machine (equal partner, dedicated server, or client), \textbf{who controls} the shared resources (files, printers, user accounts, internet access), and \textbf{how data and services flow} between machines. It answers the question: \emph{``Who is in charge, and who depends on whom?''}
\end{definition}

\begin{exemplebox}[Two networks, same hardware, different architectures]
A small printing business in Douala owns four computers and one printer, all connected to the same switch.
\begin{itemize}
\item \textbf{Choice A:} every computer stores its own documents and shares its own folders with the others; the printer is attached to one of the computers, which shares it. No machine is more important than another.
\item \textbf{Choice B:} one powerful computer is reserved to store all documents centrally, manage user passwords and host the printer; the other three must log in through it to work.
\end{itemize}
Same cables, same switch, same computers --- but two completely different \emph{architectures}. Choice A is a \cle{peer-to-peer} architecture; Choice B is a \cle{client/server} architecture.
\end{exemplebox}

\courssub{Physical Topology versus Logical Architecture}

One of the most frequent confusions at the GCE Advanced Level is between \emph{topology} and \emph{architecture}. They describe two different ``layers'' of a network's design.

\begin{definition}[Physical topology and logical architecture]
\begin{itemize}
\item The \cle{physical topology} is the \textbf{geometric layout} of the network: how the devices are physically (or wirelessly) interconnected. The classic topologies are \textbf{bus}, \textbf{star}, \textbf{ring} and \textbf{mesh}. It answers: \emph{``How are the machines wired together?''}
\item The \cle{logical architecture} (or simply \emph{architecture}) is the \textbf{organisation of roles and control} on the network: which machines provide services and which consume them. The two main architectures are \textbf{peer-to-peer} and \textbf{client/server}. It answers: \emph{``Who serves, who is served, and who controls the resources?''}
\end{itemize}
The two notions are \textbf{independent}: a star topology can carry either a peer-to-peer or a client/server architecture.
\end{definition}

\begin{attention}[Topology is NOT architecture]
A very common exam mistake is to answer ``star'' or ``bus'' when the question asks about \emph{architecture}, or to answer ``client/server'' when asked about \emph{topology}. Remember:
\begin{itemize}
\item \textbf{Topology} (bus, star, ring, mesh) = the \emph{shape} of the cabling --- what you see on a floor plan.
\item \textbf{Architecture} (peer-to-peer, client/server) = the \emph{distribution of roles} --- what you see when you ask who controls the data.
\end{itemize}
A school computer lab wired as a \emph{star} may run a \emph{client/server} architecture (all students log on to a central server). The star describes the cables; client/server describes the roles. Never mix the two vocabularies in an exam answer.
\end{attention}

The figure below shows the same physical star topology carrying two different logical architectures.

\begin{popfigure}
\begin{tikzpicture}[
  pc/.style={draw=popink, fill=popblueL, rounded corners=2pt, minimum width=1.15cm, minimum height=0.7cm, font=\scriptsize, align=center},
  srv/.style={draw=popdark, fill=popgold, rounded corners=2pt, minimum width=1.3cm, minimum height=0.8cm, font=\scriptsize, align=center},
  sw/.style={draw=popink, fill=poppurple, text=white, circle, inner sep=2pt, minimum size=0.85cm, font=\scriptsize},
  lbl/.style={font=\small\bfseries, text=popdark},
  note/.style={font=\scriptsize\itshape, text=popmuted, align=center}
]
% ===== LEFT : star topology, peer-to-peer roles =====
\node[lbl] at (0,3.1) {Same physical star topology};
\node[sw] (S1) at (0,0.6) {Switch};
\node[pc, align=center] (A1) at (-2.4,2.1){PC 1\\(equal)};
\node[pc, align=center] (A2) at (2.4,2.1){PC 2\\(equal)};
\node[pc, align=center] (A3) at (-2.4,-0.9){PC 3\\(equal)};
\node[pc, align=center] (A4) at (2.4,-0.9){PC 4\\(equal)};
\draw[popline, thick] (S1) -- (A1);
\draw[popline, thick] (S1) -- (A2);
\draw[popline, thick] (S1) -- (A3);
\draw[popline, thick] (S1) -- (A4);
\node[note, align=center] at (0,-1.9){Logical view: PEER-TO-PEER\\every PC shares its own files};
\draw[poporange, thick, <->] (A1.south) .. controls (-1.2,0.9) .. (A3.north);
\draw[poporange, thick, <->] (A2.south) .. controls (1.2,0.9) .. (A4.north);

% ===== RIGHT : same star, client/server roles =====
\begin{scope}[xshift=8.6cm]
\node[lbl] at (0,3.1) {Same physical star topology};
\node[sw] (S2) at (0,0.6) {Switch};
\node[srv, align=center] (B0) at (0,2.2){SERVER\\(controls all)};
\node[pc] (B1) at (-2.4,0.6) {Client 1};
\node[pc] (B2) at (2.4,0.6) {Client 2};
\node[pc] (B3) at (-2.4,-1.0) {Client 3};
\node[pc] (B4) at (2.4,-1.0) {Client 4};
\draw[popline, thick] (S2) -- (B0);
\draw[popline, thick] (S2) -- (B1);
\draw[popline, thick] (S2) -- (B2);
\draw[popline, thick] (S2) -- (B3);
\draw[popline, thick] (S2) -- (B4);
\node[note, align=center] at (0,-1.9){Logical view: CLIENT/SERVER\\clients depend on the server};
\draw[popgreen, thick, ->] (B1.north) -- (S2);
\draw[popgreen, thick, ->] (B2.north) -- (S2);
\end{scope}
\end{tikzpicture}
\legende{The physical star topology (cabling through a central switch) is identical on both sides, but the logical architectures differ: equals sharing with equals (left) versus clients depending on a central server (right).}
\end{popfigure}

\courssub{The Two Main Architectures: an Overview}

Network design offers many variations, but at GCE Advanced Level we focus on the two fundamental architectures. They will be studied in depth in the following lessons; here we only introduce them.

\textbf{Peer-to-peer (P2P).}\par
In a \cle{peer-to-peer} architecture, all computers are \textbf{equals} (``peers''). Each machine can act both as a provider and as a consumer of resources: it shares its own folders, and it uses folders shared by the others. There is no central controller, no dedicated server, and usually no centralised security. This is the natural architecture of a home network or a small cybercafé: cheap, simple to set up, but difficult to secure and to manage as it grows.

\textbf{Client/server.}\par
In a \cle{client/server} architecture, roles are \textbf{specialised}. One or more powerful machines, the \textbf{servers}, are dedicated to providing services: storing files, managing user accounts and passwords, hosting databases, sharing printers, delivering web pages or e-mail. The other machines, the \textbf{clients}, request and consume these services. Control and security are centralised on the server. This is the architecture of banks, ministries, universities and large companies: more expensive and requiring skilled administration, but secure, manageable and able to grow.

\begin{center}
\begin{tabularx}{\linewidth}{|l|X|X|}
\hline
\rowcolor{popblueL} \textbf{Aspect} & \textbf{Peer-to-peer} & \textbf{Client/server} \\
\hline
Roles of machines & All equal (peers) & Specialised: servers and clients \\
\hline
Control of resources & Decentralised, each user & Centralised on the server \\
\hline
Typical size & Few machines (about 2--10) & Tens to thousands of machines \\
\hline
Typical cost & Low & High (server hardware, licences, administrator) \\
\hline
Typical examples & Home network, small cybercafé & Bank, ministry, school lab with server \\
\hline
\end{tabularx}
\end{center}

\courssub{The Six Design Criteria}

How does an engineer decide which architecture to adopt? By judging each candidate design against a small set of standard criteria. These criteria are the vocabulary of every network design discussion, and examiners expect you to use them precisely.

\begin{definition}[Design criteria for a network architecture]
A \cle{design criterion} is a measurable quality used to evaluate and compare network designs. The six key criteria are:
\begin{enumerate}
\item \textbf{Cost} --- the total expense of the network: hardware (servers cost far more than ordinary PCs), software licences, cabling, and the salary of qualified administrators.
\item \textbf{Security} --- the ability to protect data and resources: user authentication, access rights, protection against unauthorised access and data loss.
\item \textbf{Scalability} --- the ability of the network to \emph{grow} (more users, more machines, more services) without a complete redesign or a collapse of performance.
\item \textbf{Performance} --- the speed and responsiveness experienced by users: how quickly files open, pages load and transactions complete, especially under heavy load.
\item \textbf{Ease of administration} --- how simply the network can be managed day to day: creating accounts, installing software, applying updates, backing up data.
\item \textbf{Reliability} --- the ability to keep working despite failures: what happens if one machine breaks down? Is there a single point of failure?
\end{enumerate}
\end{definition}

\begin{methode}[The six-criteria checklist for evaluating any network design]
Whenever you are asked to analyse, compare or choose a network architecture (in an exam or in real life), proceed systematically:
\begin{enumerate}
\item \textbf{Identify the context:} number of users, type of data (sensitive or not), budget available, expected growth.
\item \textbf{Cost:} can the organisation afford servers, licences and an administrator?
\item \textbf{Security:} does the data (bank accounts, exam results, patient files) require centralised, strict access control?
\item \textbf{Scalability:} will the network grow beyond a handful of machines in the coming years?
\item \textbf{Performance:} must the network stay fast under heavy simultaneous use?
\item \textbf{Ease of administration:} is a skilled technician available to manage a server?
\item \textbf{Reliability:} is a breakdown acceptable, or must the service run continuously?
\item \textbf{Conclude and justify:} choose the architecture whose strengths match the most critical criteria, and \emph{justify your choice in terms of the criteria}, never just by naming the architecture.
\end{enumerate}
\end{methode}

\begin{popfigure}
\begin{tikzpicture}[
  crit/.style={draw=popink, fill=popblueL, rounded corners=3pt, minimum width=2.5cm, minimum height=0.75cm, font=\scriptsize\bfseries, align=center},
  arch/.style={draw=popdark, rounded corners=4pt, minimum width=3.4cm, minimum height=1.0cm, font=\small\bfseries, align=center, text=white},
  lab/.style={font=\scriptsize, text=popmuted, align=center, midway, sloped}
]
\node[arch, fill=poporange] (p2p) at (-4.6,0) {PEER-TO-PEER};
\node[arch, fill=popgreen] (cs) at (4.6,0) {CLIENT/SERVER};
\node[crit] (c1) at (0,3.0) {Cost};
\node[crit] (c2) at (0,1.8) {Security};
\node[crit] (c3) at (0,0.6) {Scalability};
\node[crit] (c4) at (0,-0.6) {Performance};
\node[crit] (c5) at (0,-1.8) {Administration};
\node[crit] (c6) at (0,-3.0) {Reliability};
\draw[poporange, thick, ->] (c1.west) -- (p2p.north east) node[lab, above]{low cost};
\draw[popgreen, thick, ->] (c1.east) -- (cs.north west) node[lab, above]{high cost};
\draw[popgreen, thick, ->] (c2.east) -- (cs.north west) node[lab, above]{strong};
\draw[poporange, thick, ->] (c2.west) -- (p2p.east) node[lab, above]{weak};
\draw[popgreen, thick, ->] (c3.east) -- (cs.west) node[lab, above]{excellent};
\draw[poporange, thick, ->] (c3.west) -- (p2p.east) node[lab, above]{poor};
\draw[popgreen, thick, ->] (c4.east) -- (cs.west) node[lab, above]{managed};
\draw[poporange, thick, ->] (c4.west) -- (p2p.east) node[lab, below]{degrades};
\draw[popgreen, thick, ->] (c5.east) -- (cs.south west) node[lab, above]{centralised};
\draw[poporange, thick, ->] (c5.west) -- (p2p.south east) node[lab, below]{per-machine};
\draw[popgreen, thick, ->] (c6.east) -- (cs.south west) node[lab, below]{server = critical};
\draw[poporange, thick, ->] (c6.west) -- (p2p.south east) node[lab, below]{no single failure};
\end{tikzpicture}
\legende{Concept map: the six design criteria connect the two architectures. Each architecture has characteristic strengths (green = client/server, orange = peer-to-peer) on each criterion.}
\end{popfigure}

\courssub{Network Architectures Around Us in Cameroon}

The two architectures are not abstract theory --- you meet them every day.

\begin{exemplebox}[Cameroonian examples of each architecture]
\begin{itemize}
\item \textbf{Home network (peer-to-peer):} in a family home in Yaoundé, a laptop, two smartphones and a smart TV share one Wi-Fi router. Each device shares its own photos or files directly; no device is ``in charge''.
\item \textbf{Small cybercafé (peer-to-peer):} the Bamenda cybercafé of our opening story: ten PCs, one shared printer, one internet connection, no central server. Cheap to run, but one customer streaming videos slows everyone, and there is no central control of what users do.
\item \textbf{School computer lab (often client/server):} in many well-equipped secondary schools, students log on with individual accounts stored on a central server; the teacher can control, install software and save students' work centrally.
\item \textbf{Banks and microfinance institutions (client/server):} a microfinance bank centralises all account data on a secure server; tellers' machines are clients. Security and data integrity are excellent --- but if the server fails, all work stops (a reliability issue: the \emph{single point of failure}).
\item \textbf{Ministries and e-government services (client/server):} centralised servers host citizen databases and online services, with strict access control, exactly because the data is sensitive and the number of users is huge.
\end{itemize}
\end{exemplebox}

\begin{savaistu}[From standalone PCs to organised networks]
In the early days of computing, each personal computer worked \emph{alone}: files moved between machines on floppy disks (the famous ``sneakernet'' --- walking the disk across the room!). As organisations acquired more computers, the need to share expensive printers, centralise data and communicate led to the first local networks, and engineers had to decide \emph{how to organise} them. Early small office networks were naturally peer-to-peer, while large institutions adopted centralised client/server models inspired by even older mainframe systems, where one giant computer served many simple terminals. Today's cloud services are, in a sense, a gigantic client/server architecture: your phone is the client, and huge data centres are the servers. This historical note is enrichment: it will not be examined directly, but it helps you understand \emph{why} the two architectures exist.
\end{savaistu}

\courssub{Worked Example and Consolidation}

\begin{exoresolu}[Choosing an architecture for a growing business]
\textbf{Statement.} A small accounting firm in Buea starts with 3 employees sharing files and one printer. The owner expects to hire 12 more employees within two years and will then handle confidential financial records of many clients. Using the six design criteria, advise the owner on the architecture to plan for, and justify your answer.
\tcblower
\textbf{Solution.} We apply the six-criteria checklist.
\begin{enumerate}
\item \textbf{Context:} currently 3 users, growing to about 15; confidential financial data; moderate budget.
\item \textbf{Cost:} a peer-to-peer network is cheaper now, but the firm can plan an investment given the growth.
\item \textbf{Security:} confidential client records \emph{require} centralised accounts and access rights --- this favours client/server strongly.
\item \textbf{Scalability:} growth from 3 to 15 machines is painful in peer-to-peer (each machine manages its own sharing) but natural in client/server.
\item \textbf{Performance:} with 15 simultaneous users, a dedicated server keeps file access fast and stable.
\item \textbf{Ease of administration:} creating accounts and backing up data centrally is far simpler with a server; a part-time technician can be hired.
\item \textbf{Reliability:} a server introduces a single point of failure, so the firm must plan regular backups (and ideally a backup server).
\end{enumerate}
\textbf{Conclusion:} the firm may start today with a simple peer-to-peer network for its 3 machines, but it should \emph{plan and budget} for a client/server architecture, because security, scalability, performance and centralised administration --- the critical criteria for its future activity --- all point that way. The justification is expressed in terms of the criteria, as required.
\end{exoresolu}

\begin{exercice}[Analysing the cybercafé and the bank]
For each of the two organisations in the opening story (the Bamenda cybercafé and the microfinance bank):
\begin{enumerate}
\item State the architecture it uses and justify your answer from the description.
\item Identify, for each, the two design criteria that are most problematic.
\item Suggest one realistic improvement for each organisation.
\end{enumerate}
\end{exercice}

\begin{corrige}[Exercise 1]
\begin{enumerate}
\item \textbf{Cybercafé:} peer-to-peer --- ten equal PCs share a printer and a connection with no central controller. \textbf{Bank:} client/server --- one powerful machine controls all account data and the tellers' machines depend on it.
\item \textbf{Cybercafé:} \emph{performance} (streaming by one customer slows everyone, no central bandwidth control) and \emph{security/administration} (no central control of users). \textbf{Bank:} \emph{reliability} (the server is a single point of failure: when it fails, all tellers stop) and \emph{cost} (powerful server, licences, administrator).
\item \textbf{Cybercafé:} introduce a central machine (or a router with management features) acting as a server/gateway to control bandwidth per user and centralise printing --- a move towards client/server. \textbf{Bank:} add a backup (redundant) server and a strict backup plan so that tellers can keep working if the main server fails.
\end{enumerate}
\end{corrige}

\begin{aretenir}[Key points of this section]
\begin{itemize}
\item The \cle{network architecture} is the overall organisation of roles and control in a network: who provides services, who consumes them, who controls the resources.
\item \textbf{Physical topology} (bus, star, ring, mesh) describes the \emph{layout of the cabling}; \textbf{logical architecture} (peer-to-peer, client/server) describes the \emph{organisation of roles}. The two are independent --- never confuse them in an exam.
\item \textbf{Peer-to-peer:} all machines are equal; cheap and simple, but weak in security, scalability and administration. \textbf{Client/server:} specialised servers serve clients; secure, scalable and centrally managed, but costly and the server is a single point of failure.
\item Judge every design with the \textbf{six criteria}: cost, security, scalability, performance, ease of administration, reliability --- and always justify a choice in terms of these criteria.
\end{itemize}
\end{aretenir}

\coursec{Peer-to-Peer Architecture}

In the previous section we saw that a network architecture describes \emph{who does what} on a network: which machines request services and which machines provide them. We now study the first and simplest of the two great architectures: the \cle{peer-to-peer} model, in which \emph{every} machine plays both roles.

\courssub{Intuition: a group of friends sharing equally}

Imagine five students in Buea who each own a laptop. Manka has music files, Tabi has a printer, Enow has course notes in PDF, Ali has films, and Sirri has past GCE questions. They connect their laptops to the same small router (or the same hotspot) and agree: ``whatever I have, you may take; whatever you have, I may take.'' Nobody is the boss. Each student decides \emph{personally} what to share and with whom.

That is exactly the spirit of a peer-to-peer network. There is no special, powerful machine whose only job is to serve the others. Every computer is simultaneously a \emph{client} (when it fetches a file from a neighbour) and a \emph{server} (when it gives a file to a neighbour).

\begin{definition}[Peer-to-peer network]
A \cle{peer-to-peer network} (P2P network) is a network architecture in which all computers have \textbf{equal status}: each machine, called a \emph{peer}, can act \textbf{both as a client and as a server}. Resources (files, folders, printers) are shared directly between peers, without any dedicated central server.
\end{definition}

\begin{definition}[Peer and workgroup]
A \cle{peer} is any computer in a peer-to-peer network; the word stresses that all machines are ``equals''. A \cle{workgroup} is the Microsoft Windows implementation of a peer-to-peer network: a named logical group of computers (for example \texttt{WORKGROUP}) on the same local network, in which each machine manages its \emph{own} user accounts and its \emph{own} shared resources. There is no central administrator and no central database of users.
\end{definition}

\begin{savaistu}[Did you know?]
When you send a photo from your phone to a friend's laptop through Bluetooth, or share a document between two laptops connected to the same phone hotspot in a Douala student hostel, you are using peer-to-peer logic: two equal devices exchanging data directly, with no server in between. Large-scale systems such as BitTorrent generalise the same idea to millions of peers on the Internet.
\end{savaistu}

\courssub{How resources are shared in a P2P network}

The key idea is \emph{decentralised control}. In a peer-to-peer network:

\begin{itemize}
\item \textbf{Each user controls his or her own sharing.} Tabi alone decides that her folder \texttt{PastQuestions} is shared, and she alone sets the permissions (read-only or read/write) and the password.
\item \textbf{Sharing is configured machine by machine.} If five computers must each share a folder, the sharing must be configured five times, once on each machine.
\item \textbf{Accounts are local.} A user account created on Enow's laptop exists only on that laptop. To let Ali connect, Enow must create an account (or enable a sharing password) \emph{on his own machine}.
\item \textbf{Any peer can offer any service.} A printer plugged into one PC can be shared so that all other peers print through it; that PC then acts as a print server \emph{in addition to} being somebody's ordinary workstation.
\end{itemize}

\begin{popfigure}
\begin{tikzpicture}[
  pc/.style={rectangle, draw=popblue, thick, fill=popblueL, rounded corners=2pt, minimum width=1.7cm, minimum height=0.8cm, align=center, font=\small},
  res/.style={font=\scriptsize, text=popdark, align=center},
  sw/.style={rectangle, draw=poporange, very thick, fill=white, rounded corners=2pt, minimum width=2.2cm, minimum height=0.8cm, align=center, font=\small\bfseries}
]
\node[sw] (hub) at (0,0) {Switch / Router};
\node[pc] (p1) at (90:3.1) {PC 1 (Manka)};
\node[pc] (p2) at (162:3.1) {PC 2 (Tabi)};
\node[pc] (p3) at (234:3.1) {PC 3 (Enow)};
\node[pc] (p4) at (306:3.1) {PC 4 (Ali)};
\node[pc] (p5) at (18:3.1) {PC 5 (Sirri)};
\draw[thick, popline] (hub) -- (p1);
\draw[thick, popline] (hub) -- (p2);
\draw[thick, popline] (hub) -- (p3);
\draw[thick, popline] (hub) -- (p4);
\draw[thick, popline] (hub) -- (p5);
\node[res, anchor=south] at (90:3.9) {shares folder \texttt{Music}};
\node[res, anchor=east, align=right] at (162:4.1) {shares\\printer};
\node[res, anchor=north east, align=right] at (234:3.9) {shares folder \texttt{Notes}};
\node[res, anchor=north west, align=left] at (306:3.9) {shares folder \texttt{Films}};
\node[res, anchor=west, align=left] at (18:4.1) {shares folder \texttt{GCE}};
\end{tikzpicture}
\legende{A five-computer peer-to-peer network in a star layout. Every PC is equal: each one shares its own folder or printer and can use the resources of the others. There is no dedicated server.}
\end{popfigure}

\begin{attention}[No dedicated server does not mean no switch]
Do not confuse the \emph{physical} topology with the \emph{logical} architecture. The five PCs above are physically wired in a star around a switch, but the switch only forwards frames; it provides no file service. ``No dedicated server'' means no machine is reserved for serving the others --- not that there is no network equipment at all.
\end{attention}

\courssub{Characteristics of a peer-to-peer network}

\begin{propriete}[Characteristics of a peer-to-peer network]
A typical peer-to-peer network has the following characteristics:
\begin{itemize}
\item \textbf{Small size:} usually \textbf{fewer than 10 computers}; beyond that, management becomes chaotic.
\item \textbf{Low cost:} no dedicated server machine and no expensive server operating system to buy.
\item \textbf{Simple setup:} ordinary desktop operating systems (Windows Home editions, Linux desktops) are sufficient.
\item \textbf{Decentralised administration:} each user administers his or her own machine; there is no central administrator.
\item \textbf{No single point of failure for the whole network:} if one peer is switched off, only \emph{its} resources disappear; the rest of the network keeps working.
\end{itemize}
These characteristics are directly examinable: learn them with a justification for each.
\end{propriete}

\courssub{Setting up a Windows workgroup}

\begin{methode}[Steps to set up a simple Windows workgroup for file and printer sharing]
\begin{enumerate}
\item \textbf{Connect the machines.} Plug all PCs into the same switch or connect them to the same Wi-Fi router/hotspot, and check that each has a valid IP address on the same network (e.g.\ via \texttt{ipconfig}).
\item \textbf{Put all PCs in the same workgroup.} Open \emph{System Properties} $\rightarrow$ \emph{Computer name} and set the same workgroup name (e.g.\ \texttt{CYBERCAFE}) on every machine; give each PC a distinct computer name.
\item \textbf{Turn on sharing.} In the \emph{Network and Sharing Center}, enable \emph{Network discovery} and \emph{File and printer sharing}; set the network profile to \emph{Private}.
\item \textbf{Share the resources.} On each PC, right-click the folder or printer $\rightarrow$ \emph{Properties} $\rightarrow$ \emph{Sharing}, choose the users or groups allowed, and set their permission level (Read or Read/Write).
\item \textbf{Create local user accounts} on each sharing machine for the people who must connect (or configure password-protected sharing).
\item \textbf{Test from another peer.} From a second PC, open the network view (or type \texttt{\textbackslash\textbackslash PCNAME}), authenticate, and open a shared file; then print a test page on the shared printer.
\end{enumerate}
\end{methode}

\begin{exoresolu}[Designing a small P2P network]
A small cybercafé in Bamenda has 6 identical PCs and one printer. The owner wants customers' files and the printer to be shared, but has a very small budget and no trained administrator. Propose an architecture and justify it; state two precautions the owner should take.
\tcblower
\textbf{Solution.} With only 6 machines (fewer than 10), a tiny budget and no administrator, the \textbf{peer-to-peer architecture} is appropriate: no dedicated server must be bought, and an ordinary Windows workgroup is enough.
\begin{itemize}
\item Connect the 6 PCs to one switch; put them all in the same workgroup, e.g.\ \texttt{CAFE}.
\item Plug the printer into one PC and share it; the other five print through that peer.
\item On each PC, share only the intended folder, with \emph{Read} permission where modification is not needed.
\end{itemize}
\textbf{Precautions:} (1) set a sharing password / local accounts on every PC so that only staff connect, since there is no central security control; (2) agree on a backup routine, because data is scattered across all six machines and a disk failure on one PC destroys the files stored there.
\end{exoresolu}

\courssub{Advantages and disadvantages}

\begin{popfigure}
\begin{tikzpicture}
\node[rectangle, draw=popgreen, very thick, fill=popgreenL, rounded corners=3pt, text width=6.6cm, align=left, anchor=north west] (adv) at (0,0)
{\textbf{\textcolor{popdark}{Advantages}}\\[2pt]
\textbullet\ Cheap: no server to buy\\
\textbullet\ Easy and quick to install\\
\textbullet\ No dedicated server needed\\
\textbullet\ No dependence on a single machine\\
\textbullet\ Each user keeps control of own files};
\node[rectangle, draw=poppink, very thick, fill=white, rounded corners=3pt, text width=6.6cm, align=left, anchor=north east] (dis) at (\linewidth,0)
{\textbf{\textcolor{poppink}{Disadvantages}}\\[2pt]
\textbullet\ Weak, scattered security\\
\textbullet\ No centralised administration\\
\textbullet\ Performance drops as network grows\\
\textbullet\ Data scattered: backups are hard\\
\textbullet\ A switched-off peer removes its resources};
\end{tikzpicture}
\legende{Summary of the strengths and weaknesses of the peer-to-peer architecture.}
\end{popfigure}

\textbf{Advantages.} The architecture is \textbf{cheap}: all machines are ordinary PCs, and no server licence is required. It is \textbf{easy to install and configure}: creating a workgroup takes minutes. It needs \textbf{no dedicated server} and therefore no specialist administrator. Finally, there is \textbf{no dependence on a single machine}: unlike a client/server network, the failure of one peer does not paralyse the whole network.

\textbf{Disadvantages.} Security is \textbf{weak}: permissions are set separately on each machine, so the overall policy is inconsistent and easy to misconfigure. Administration is \textbf{not centralised}: creating an account for a new employee may mean touching every PC. Performance \textbf{degrades as the network grows}, because each PC shares its processor time between its user's work and the requests of others. And because \textbf{data is scattered} across all machines, backup, updating and locating files become difficult.

\begin{attention}[Common mistake: ``P2P has no security at all'']
This is false and costs marks. Each peer \emph{can} set permissions: shared folders can be password-protected and limited to read-only, and local user accounts control who connects. What P2P lacks is \textbf{centralised} security: there is no single place where an administrator defines and enforces one policy for the whole network. Say precisely this in the examination.
\end{attention}

\courssub{Where is P2P suitable?}

Peer-to-peer is the right design strategy whenever the network is \emph{small}, the budget is \emph{tight} and administration skills are \emph{limited}:

\begin{itemize}
\item \textbf{Homes:} a family sharing photos and one printer between three or four devices.
\item \textbf{Small offices:} a two-to-eight-person business (an accounting firm, a nursery school office) sharing documents.
\item \textbf{Small cybercafés} with only a few machines, as in the worked example.
\item \textbf{Ad-hoc exchanges:} two laptops sharing files through Bluetooth or a hotspot --- everyday P2P logic, even without a permanent network (enrichment).
\end{itemize}

As soon as the number of users grows, security requirements rise or data must be centralised, the \emph{client/server} architecture studied in the next section becomes the better choice.

\begin{exercice}[Checking your understanding]
\begin{enumerate}
\item Define a peer-to-peer network and explain the double role of each peer.
\item A shop in Yaoundé has 4 computers and one printer. Justify the choice of a P2P architecture and describe how the printer can be shared.
\item Give three disadvantages of P2P that would push a growing company of 60 employees to abandon it.
\item True or false (justify): ``In a peer-to-peer network, it is impossible to protect a shared folder with a password.''
\end{enumerate}
\end{exercice}

\begin{corrige}[Exercise 1]
\begin{enumerate}
\item A P2P network is an architecture in which all computers have equal status and share resources directly, with no dedicated server. Each peer acts as a \emph{client} when it requests a resource from another peer, and as a \emph{server} when it provides one of its own resources (folder, printer) to others.
\item With only 4 machines, a small budget and no administrator, P2P is suitable: connect the PCs to one switch, place them in the same Windows workgroup, plug the printer into one PC, enable \emph{File and printer sharing} and share the printer; the other three PCs then print through that peer.
\item (Any three) Weak, non-centralised security; administration must be repeated on each machine; performance degrades as the number of peers grows; data scattered across machines makes backup difficult. With 60 employees these problems become unmanageable, so a client/server architecture with centralised accounts and security is preferable.
\item \textbf{False.} Each peer can protect its shares with passwords, local accounts and permission levels (Read, Read/Write). What is missing is \emph{centralised} control of security for the whole network, not security itself.
\end{enumerate}
\end{corrige}

\coursec{Client/Server Architecture}

In the previous section we explored peer-to-peer networks, where every computer acts as both a client and a server, sharing resources directly with one another. That model works well for very small groups, but it quickly becomes difficult to manage as the number of users grows: passwords are scattered across many machines, backups are inconsistent, and there is no single point to enforce security policies. To overcome these limitations, organisations — from the Ministry of Finance in Yaoundé to a large secondary school in Buea — adopt a fundamentally different architecture: the \cle{client/server} model. In this section we define the model, identify the roles of servers and clients, list the main types of servers, explain how centralised administration works, and weigh the advantages and disadvantages so that you can decide when this architecture is the right choice.

\courssub{Fundamental Concepts: Server, Client and Dedicated Server}

\begin{definition}[Client/Server Network]
A \textbf{client/server network} is a network architecture in which one or more \emph{dedicated servers} provide services and resources to \emph{client} computers. The server waits passively for requests; the client initiates the communication by asking for a specific service (file access, printing, authentication, web page, database query, etc.). The server processes the request and sends back a response.
\end{definition}

\begin{definition}[Server, Client and Dedicated Server]
\begin{itemize}
\item \textbf{Server:} A computer (or a virtual machine) that runs a \emph{server operating system} (e.g., Windows Server 2022, Ubuntu Server, Red Hat Enterprise Linux) and server applications. Its primary purpose is to \emph{serve} other computers. It typically has more RAM, multiple CPU cores, redundant power supplies and RAID storage arrays.
\item \textbf{Client:} A workstation (desktop, laptop, thin client, smartphone) that runs a \emph{client operating system} (Windows 10/11 Pro, Ubuntu Desktop, macOS) and client applications. It \emph{requests} services from the server.
\item \textbf{Dedicated server:} A server that is \emph{not} used as a regular user workstation. It does not run user productivity software (word processors, games, web browsers for daily surfing). Its sole job is to provide network services reliably and securely.
\end{itemize}
\end{definition}

\begin{attention}[Common Mistake: Powerful Hardware $\neq$ Server]
A common error in examinations is to state that ``the most powerful computer in the room is the server''. \textbf{Wrong.} A server is defined by its \emph{role} and the \emph{software} it runs, not only by its hardware specifications. You can install Windows Server on a modest machine and it will act as a server (though performance may suffer). Conversely, a high-end gaming PC running Windows 11 Home is a \emph{client} if it only requests services. In the GCE, always justify the server/client distinction by referring to the \emph{operating system} and the \emph{services provided}.
\end{attention}

\courssub{Main Types of Servers and Their Services}

A single physical server can host several server roles simultaneously (e.g., a domain controller that also runs DNS and DHCP). In large organisations, each role may be placed on a separate dedicated machine for performance and fault isolation. The table below summarises the server types required by the syllabus.

\begin{propriete}[Main Types of Servers and Services Provided]
\begin{center}
\begin{tabularx}{\linewidth}{|l|X|X|}
\hline
\rowcolor{popblueL}
\textbf{Server Type} & \textbf{Service Provided} & \textbf{Typical Software Examples} \\
\hline
File Server & Centralised storage, shared folders, access control (NTFS permissions, Samba shares) & Windows File Services, Samba (Linux), FreeNAS/TrueNAS \\
Print Server & Manages print queues, distributes jobs to network printers, driver repository & Windows Print Services, CUPS (Linux) \\
Mail Server & Sends, receives, stores and forwards e-mail (SMTP, IMAP, POP3) & Microsoft Exchange, Postfix + Dovecot, Zimbra \\
Web Server & Hosts websites and web applications (HTTP/HTTPS) & Apache HTTP Server, Nginx, Microsoft IIS \\
Database Server & Stores structured data, processes SQL queries, enforces integrity & Microsoft SQL Server, PostgreSQL, MySQL/MariaDB, Oracle Database \\
Application Server & Runs business logic for client applications (ERP, CRM, custom apps) & WildFly/JBoss, GlassFish, Tomcat (Java EE), .NET Core on Linux/Windows \\
\hline
\end{tabularx}
\end{center}
\end{propriete}

\begin{savaistu}[Cameroon Context: Servers in Public Administration]
Many Cameroonian ministries and parastatals (e.g., \emph{MINFI}, \emph{MINESEC}, \emph{CAMTEL}, \emph{ANTIC}) run their critical databases and e-government portals on centralised database and application servers housed in the \emph{National Data Centre} (or in secure server rooms). These servers are often virtualised on clusters to ensure high availability. When you register for the GCE online or check results on the \emph{GCE Board} portal, your browser (client) talks to a web server, which in turn queries a database server — a classic multi-tier client/server interaction.
\end{savaistu}

\courssub{Visualising the Architecture}

The diagram below shows a typical client/server LAN: a central server connected to a switch, with six client workstations. The server is labelled with the multiple roles it may simultaneously fulfil.

\begin{popfigure}
\begin{tikzpicture}[
  node distance=1.5cm and 2cm,
  server/.style={draw, thick, rounded corners, fill=popblue!10, minimum width=2.5cm, minimum height=1.8cm, align=center},
  client/.style={draw, thick, rounded corners, fill=popgreen!10, minimum width=1.8cm, minimum height=1.2cm, align=center},
  switch/.style={draw, thick, circle, fill=poporange!20, minimum size=1.2cm},
  link/.style={thick, popdark}
]
% Switch at center
\node[switch] (sw) {Switch};

% Server above switch
\node[server, above=of sw, align=center] (srv){
  \textbf{DEDICATED SERVER}\\
  \small Windows Server 2022 / Ubuntu Server\\
  \textit{Roles:}\\
  File Server \textbullet\ Print Server\\
  Domain Controller (AD) \textbullet\ DNS/DHCP\\
  Database Server \textbullet\ Web Server
};

% Clients around switch
\foreach \angle/\lbl in {150/Client 1, 110/Client 2, 70/Client 3, 30/Client 4, -30/Client 5, -70/Client 6} {
  \node[client] (c\angle) at ($(sw)+(\angle:3cm)$) {\lbl\\ \small Windows 11 Pro / Ubuntu Desktop};
  \draw[link] (sw) -- (c\angle);
}

% Server to switch
\draw[link, very thick] (srv) -- (sw) node[midway, left, align=center] {\small 1 Gbps / 10 Gbps\\ uplink};

% Legend
\node[below=1.2cm of sw, align=center, text width=12cm, font=\small] {
  \textbf{Figure 3.1:} Client/Server LAN topology. The dedicated server (top) provides multiple services to six clients via a central switch. Each client authenticates against the Domain Controller and accesses shared files, printers, databases and intranet web sites hosted on the server.
};
\end{tikzpicture}
\legende{Client/Server LAN with a multi-role dedicated server and six client workstations.}
\end{popfigure}

\courssub{The Client Request -- Server Response Cycle}

Every interaction in a client/server network follows a predictable pattern: the client initiates a request, the server processes it, and the server returns a response. Understanding this cycle is essential for troubleshooting and for answering GCE questions on network protocols.

\begin{experience}[Tracing a Login and File Access]
Imagine a teacher at \emph{Government Bilingual High School Molyko} (Buea) logs on to a domain-joined laptop and opens a shared marksheet stored on the file server. The sequence below shows the key steps (simplified for clarity).
\end{experience}

\begin{popfigure}
\begin{tikzpicture}[
  node distance=1.2cm and 1.5cm,
  proc/.style={draw, thick, rectangle, rounded corners, fill=popblue!10, minimum width=3cm, minimum height=1cm, align=center},
  arrow/.style={->, thick, popdark},
  label/.style={font=\small, align=center}
]
% Client side
\node[proc, align=center] (c1){Client: User types\\ username/password};
\node[proc, below=of c1, align=center] (c2){Client: Sends Kerberos\\ AS-REQ to DC (UDP 88)};
\node[proc, below=of c2, align=center] (c3){Client: Receives TGT,\\ requests service ticket (TGS-REQ)};
\node[proc, below=of c3, align=center] (c4){Client: Presents service ticket\\ to File Server (SMB port 445)};
\node[proc, below=of c4, align=center] (c5){Client: Opens marksheet.xlsx\\ (read/write)};

% Server side (Domain Controller)
\node[proc, right=3.5cm of c1, align=center] (s1){Domain Controller (AD):\\ Verifies credentials,\\ issues TGT};
\node[proc, right=3.5cm of c2, align=center] (s2){Domain Controller (AD):\\ Issues service ticket\\ for File Server (CIFS)};
\node[proc, right=3.5cm of c4, align=center] (s3){File Server: Validates ticket,\\ checks NTFS ACLs,\\ grants access};

% Arrows
\draw[arrow] (c1) -- node[label, midway, above] {1. Login} (s1);
\draw[arrow] (c2) -- node[label, midway, above] {2. AS-REQ} (s1);
\draw[arrow] (s1) -- node[label, midway, above] {3. AS-REP (TGT)} (c2);
\draw[arrow] (c3) -- node[label, midway, above] {4. TGS-REQ} (s2);
\draw[arrow] (s2) -- node[label, midway, above] {5. TGS-REP (ST)} (c3);
\draw[arrow] (c4) -- node[label, midway, above, align=center]{6. SMB Session Setup\\ with Service Ticket} (s3);
\draw[arrow] (s3) -- node[label, midway, above] {7. Access Granted / Data} (c4);
\draw[arrow] (c4) -- (c5);

% Dashed line separating client/server
\draw[dashed, popmuted, thick] ($(c1.west)+(-0.5,0.5)$) -- ($(c5.west)+(-0.5,-0.5)$) node[midway, left, rotate=90, font=\small] {Network};

% Labels
\node[above=0.2cm of c1, font=\bfseries\small, popblue] {CLIENT SIDE};
\node[above=0.2cm of s1, font=\bfseries\small, poporange] {SERVER SIDE (DC)};
\node[above=0.2cm of s3, font=\bfseries\small, popgreen] {SERVER SIDE (FILE)};
\end{tikzpicture}
\legende{Client request -- server response cycle: Kerberos authentication (steps 1--5) followed by SMB file access (steps 6--7).}
\end{popfigure}

\begin{methode}[How Centralised Authentication Works (Step-by-Step)]
\begin{enumerate}
\item \textbf{Account creation:} The network administrator creates a user account \emph{once} on the Domain Controller (Active Directory or OpenLDAP). The account stores the username, hashed password, group memberships and password policies.
\item \textbf{Domain join:} Each client computer is \emph{joined to the domain}. This creates a trust relationship and a computer account in the directory.
\item \textbf{User logon:} When a user starts a client workstation, they enter their credentials. The client contacts the Domain Controller (via Kerberos on UDP/TCP 88) to obtain a \emph{Ticket Granting Ticket (TGT)}.
\item \textbf{Resource access:} When the user tries to open a shared folder, the client requests a \emph{service ticket} for the file server (CIFS/SMB). The Domain Controller issues the ticket encrypted with the file server's secret key.
\item \textbf{Access decision:} The client presents the service ticket to the file server. The file server decrypts it, checks the user's Security Identifier (SID) against the folder's Access Control List (ACL), and grants or denies access.
\item \textbf{Single sign-on:} Because the TGT is cached, the user can access multiple servers (mail, web, database) without re-entering their password during the session.
\end{enumerate}
\end{methode}

\courssub{Advantages and Disadvantages of Client/Server Architecture}

\begin{aretenir}[Advantages of Client/Server Networks]
\begin{itemize}
\item \textbf{Strong centralised security:} User accounts, passwords, group policies and permissions are managed in one place (the directory service). Revoking access for a departed employee takes seconds.
\item \textbf{Easy, reliable backup:} Critical data resides on the server's RAID arrays. A single scheduled backup job protects the entire organisation's data. No need to chase users' local hard drives.
\item \textbf{Scalability:} Adding users only requires creating accounts on the server. The server hardware can be upgraded (more RAM, CPUs, storage) or additional servers can be added (load balancing, clustering) to handle growth from 50 to 5,000 users.
\item \textbf{Better performance for large networks:} Dedicated server hardware (multi-core CPUs, ECC RAM, 10 GbE NICs, NVMe RAID) handles heavy concurrent workloads (database transactions, large file transfers) far better than peer-to-peer workstations.
\item \textbf{Centralised software deployment and updates:} Group Policy, SCCM/Intune, Ansible or SaltStack can push applications, patches and configuration changes to all clients automatically.
\item \textbf{Audit and compliance:} Centralised logging (Windows Event Forwarding, syslog, SIEM) makes it possible to meet regulatory requirements (e.g., financial audit trails for banks in Douala).
\end{itemize}
\end{aretenir}

\begin{aretenir}[Disadvantages of Client/Server Networks]
\begin{itemize}
\item \textbf{High initial and recurring cost:} Server-grade hardware (redundant PSUs, RAID controllers, ECC RAM), server OS licences (Windows Server CALs, RHEL subscriptions), and skilled system administrators are expensive. A small cybercafé in Bamenda cannot justify this.
\item \textbf{Single point of failure:} If the only domain controller or file server goes down, \emph{no one} can log on or access files. Mitigation requires redundancy (second DC, failover clustering, VM replication) which adds cost and complexity.
\item \textbf{Complex setup and maintenance:} Designing Active Directory sites, configuring DNS/DHCP, managing Group Policy Objects (GPOs), hardening servers, patching monthly — all require specialised knowledge (MCSE, RHCE, LPIC).
\item \textbf{Network dependency:} Clients \emph{must} have network connectivity to the server for authentication and resource access. Offline work requires cached credentials and offline files configuration.
\item \textbf{Overkill for tiny networks:} For 2--5 computers in a home or very small office, the administrative burden outweighs the benefits.
\end{itemize}
\end{aretenir}

\courssub{Suitable Contexts and Hybrid Situations}

\begin{exemplebox}[Where Client/Server is the Clear Choice]
\begin{itemize}
\item \textbf{Commercial banks (Afriland First Bank, UBA Cameroon, Ecobank):} Core banking application runs on clustered application/database servers; teller workstations are thin clients. Zero tolerance for data inconsistency.
\item \textbf{Universities (University of Yaoundé I, University of Buea):} Thousands of students/staff; centralised identity management (LDAP/AD), library systems, e-learning platforms (Moodle on web/app servers), research databases.
\item \textbf{Ministries and Government Agencies (MINSANTE, MINCOMMERCE, ANTIC):} Secure document management, HR/payroll systems, e-government portals. Compliance with national cybersecurity directives.
\item \textbf{Large secondary schools (e.g., GBHS Molyko, Lycée Leclerc Yaoundé):} 50+ workstations in labs, admin offices, library. Centralised student records, printing, internet filtering (proxy), automated backups of exam materials.
\item \textbf{Medium/Large enterprises (CDC, SONARA, Brasseries du Cameroun):} ERP systems (SAP, Oracle, Sage) on application servers; barcode scanners and PCs as clients.
\end{itemize}
\end{exemplebox}

\begin{savaistu}[Enrichment: Hybrid Architectures in Practice]
Real-world networks are rarely pure. A \textbf{hybrid} design is common:
\begin{itemize}
\item A school may run a domain controller for staff and admin offices (client/server) but allow the computer lab to operate in a workgroup (peer-to-peer) for simplicity during practical sessions.
\item A small business with 15 employees might use a NAS (Synology/QNAP) as a file/print server (client/server for storage) while workstations share a directly connected USB printer peer-to-peer.
\item \textbf{GCE Discussion Tip:} If asked ``Can a network use both architectures simultaneously?'', answer \textbf{Yes}. Explain that different \emph{segments} or \emph{services} may follow different models. The defining characteristic is \emph{per service}: if a dedicated machine provides a service to others, that service follows client/server; if two workstations share a folder directly, that share follows peer-to-peer.
\end{itemize}
\end{savaistu}

\courssub{Worked Example: Designing a Small Office Network}

\begin{exoresolu}[Scenario: Law Firm in Douala]
\textbf{Statement:} A law firm in Douala has 12 employees: 8 lawyers, 2 paralegals, 1 accountant, 1 receptionist. They need:
\begin{itemize}
\item Centralised storage of case files (confidential, must be backed up daily).
\item Shared network printer (heavy duty).
\item Internal e-mail for formal communication (not Gmail/Yahoo).
\item A simple intranet site for templates and precedents.
\item Each user must log on with their own credentials; access to case files is restricted by practice area (Corporate, Litigation, Property).
\end{itemize}
The managing partner asks you to recommend a network architecture and justify your choice.

\textbf{Detailed Solution:}

\textbf{1. Recommended Architecture: Client/Server.}
\textbf{Justification:}
\begin{itemize}
\item \textbf{Confidentiality and access control:} Case files require strict permissions (Corporate lawyers $\neq$ Litigation lawyers). Only a domain-based file server with NTFS ACLs / Samba + LDAP can enforce this centrally.
\item \textbf{Backup requirement:} Daily backup of all case files is trivial with a single file server (Veeam, Windows Server Backup, rsync to off-site). Peer-to-peer would require 12 separate backup configurations.
\item \textbf{User management:} 12 users, each with unique credentials, password expiry, account lockout policy. Active Directory (or OpenLDAP + Samba) handles this natively.
\item \textbf{Services needed:} File server, Print server, Mail server (Postfix/Dovecot or Exchange), Web server (Apache/Nginx for intranet). All can run on one or two physical servers (virtualised).
\item \textbf{Scalability:} If the firm grows to 30 staff, the domain scales without redesign.
\end{itemize}

\textbf{2. Server Specification (Example):}
\begin{itemize}
\item \textbf{Physical Server 1 (Virtualisation Host):} Dell PowerEdge R450 / HP ProLiant DL360 Gen10, $2\times$ Intel Xeon Silver 4310, 128 GB ECC RAM, $2\times$ 480 GB SSD (RAID 1 for hypervisor), $4\times$ 2 TB NVMe (RAID 10 for VM storage), Dual 10 GbE NICs, Redundant PSU.
\item \textbf{VM 1 -- Domain Controller + DNS + DHCP + Print Server:} Windows Server 2022 Standard (2 vCPU, 8 GB RAM, 80 GB disk).
\item \textbf{VM 2 -- File Server + Intranet Web Server:} Ubuntu Server 22.04 LTS (4 vCPU, 16 GB RAM, 2 TB disk for shares, Samba + Apache).
\item \textbf{VM 3 -- Mail Server:} Ubuntu Server 22.04 LTS (2 vCPU, 8 GB RAM, 200 GB disk, Postfix + Dovecot + Roundcube webmail).
\item \textbf{Backup:} Veeam Backup \& Replication (free edition) to a separate NAS (Synology DS923+, $4\times$ 4 TB HDD RAID 5) in a different room. Daily incremental, weekly full, monthly off-site copy.
\end{itemize}

\textbf{3. Client Specification:}
12 business-grade laptops (HP ProBook / Dell Latitude), Windows 11 Pro, joined to domain. Group Policy pushes: mapped drives (G: for Corporate, H: for Litigation, I: for Property), printer, intranet homepage, Windows Update from WSUS.

\textbf{4. Estimated Cost Context (Cameroon, indicative):}
Server hardware $\approx$ 3.5--4.5 million FCFA; Windows Server 2022 Standard (16 cores) + 12 User CALs $\approx$ 1.2--1.5 million FCFA; NAS + disks $\approx$ 800k FCFA; Laptops $\approx$ 600k FCFA each. Total $\approx$ 12--15 million FCFA. Justifiable for a professional firm handling high-value cases.

\textbf{5. Why Not Peer-to-Peer?}
\begin{itemize}
\item No centralised permission enforcement $\rightarrow$ risk of data leakage.
\item 12 separate backup jobs $\rightarrow$ high probability of missed backups.
\item Password management nightmare (12 local accounts $\times$ password changes).
\item No single point to deploy security patches or antivirus policies.
\end{itemize}
\end{exoresolu}

\courssub{Summary and Examination Focus}

\begin{aretenir}[Key Points for GCE Revision]
\begin{itemize}
\item \textbf{Definition:} Client/server = dedicated server(s) provide services to clients. Server defined by \emph{role + server OS}, not just hardware.
\item \textbf{Server types:} File, Print, Mail, Web, Database, Application. Know one typical software for each (e.g., Samba, CUPS, Postfix, Apache, PostgreSQL, Tomcat).
\item \textbf{Centralised authentication:} Account created once on Domain Controller $\rightarrow$ user logs on from any domain-joined client $\rightarrow$ Kerberos tickets grant access to resources $\rightarrow$ single sign-on.
\item \textbf{Advantages:} Security, backup, scalability, performance, centralised management, audit.
\item \textbf{Disadvantages:} Cost, single point of failure, complexity, network dependency, overkill for tiny nets.
\item \textbf{Suitable for:} Banks, universities, ministries, large schools, enterprises ($>$ 10--20 users, sensitive data, regulatory needs).
\item \textbf{Hybrid:} Real networks often mix both; justify per service/segment.
\end{itemize}
\end{aretenir}

\begin{attention}[Examination Traps to Avoid]
\begin{itemize}
\item \textbf{Do not confuse} ``server'' (role) with ``mainframe'' or ``supercomputer''. A Raspberry Pi running Samba is a file server.
\item \textbf{Do not write} ``peer-to-peer has no server''. It has \emph{no dedicated server}; every node can act as a server.
\item \textbf{When asked for ``disadvantages'',} do not list ``virus spreads easily'' — that is a peer-to-peer trait. In client/server, a compromised server is catastrophic, but clients are more controlled.
\item \textbf{In diagram questions:} Draw a central server connected to a switch, clients connected to the same switch. Label the server with at least two roles (e.g., ``File Server + Domain Controller''). Show the request/response arrows if asked for the cycle.
\end{itemize}
\end{attention}

\begin{exercice}[Practice Questions]
\begin{enumerate}
\item Define the term ``dedicated server'' and explain why a high-specification gaming PC running Windows 11 Home is \emph{not} a dedicated server in a client/server network.
\item List \textbf{four} distinct types of servers and, for each, state the network service it provides and the standard port(s) it typically uses.
\item Describe the sequence of events when a domain user logs on to a client workstation and then opens a file on a shared folder, mentioning Kerberos, TGT, service ticket and ACL.
\item A small NGO in Bertoua has 6 laptops and a limited budget. They need shared files and a shared printer. Advise whether they should implement a client/server or peer-to-peer architecture. Justify your answer with \textbf{three} reasons.
\item ``A client/server network eliminates the risk of data loss.'' State whether this statement is true or false and explain your reasoning.
\item (Discussion) Explain how a secondary school with 40 computers in two labs and an admin block could use a \textbf{hybrid} architecture. Identify which segments would use client/server and which might use peer-to-peer, and why.
\end{enumerate}
\end{exercice}

\begin{corrige}[Exercise 1]
A dedicated server is a computer that runs a server operating system and server applications whose sole purpose is to provide network services to clients; it is not used as a regular user workstation. A gaming PC running Windows 11 Home is a client OS, runs user applications (games, browsers), and does not provide managed network services — therefore it is a client, not a server, regardless of its hardware power.
\end{corrige}

\begin{corrige}[Exercise 2]
\begin{enumerate}
\item File Server — provides centralised file storage and sharing — TCP 445 (SMB/CIFS), TCP 2049 (NFS).
\item Print Server — manages print queues and distributes jobs — TCP 9100 (RAW), TCP 631 (IPP/CUPS).
\item Mail Server — sends/receives/stores e-mail — TCP 25 (SMTP), TCP 110 (POP3), TCP 143 (IMAP), TCP 587 (Submission), TCP 993 (IMAPS), TCP 995 (POP3S).
\item Web Server — hosts websites — TCP 80 (HTTP), TCP 443 (HTTPS).
\item Database Server — processes SQL queries — TCP 1433 (MS SQL), TCP 3306 (MySQL/MariaDB), TCP 5432 (PostgreSQL), TCP 1521 (Oracle).
\item Application Server — runs business logic — TCP 8080/8443 (Tomcat), TCP 9990 (WildFly management), varies by platform.
\end{enumerate}
(Any four with correct service and port earn full marks.)
\end{corrige}

\begin{corrige}[Exercise 3]
1) User enters credentials at client logon screen. 2) Client sends AS-REQ (Kerberos) to Domain Controller (KDC) on UDP/TCP 88. 3) KDC verifies password hash, issues TGT (Ticket Granting Ticket) encrypted with user's key — AS-REP. 4) Client decrypts TGT, caches it. 5) User accesses shared folder $\rightarrow$ client sends TGS-REQ to KDC requesting service ticket for File Server (CIFS). 6) KDC issues service ticket (ST) encrypted with File Server's secret key — TGS-REP. 7) Client presents ST to File Server via SMB (port 445). 8) File Server decrypts ST, extracts user SID/group SIDs, checks folder's NTFS ACL. 9) If ACE matches $\rightarrow$ access granted; file data returned.
\end{corrige}

\begin{corrige}[Exercise 4]
Recommendation: \textbf{Peer-to-peer} (workgroup) is more appropriate.
Reasons: 1) Only 6 users — administrative overhead of a domain (server hardware, OS licence, CALs, sysadmin) is disproportionate. 2) Budget limited — a simple NAS (Synology/QNAP) or one laptop sharing a folder and printer meets needs at low cost. 3) No complex security requirements (no strict compartmentalisation, no regulatory audit) — local accounts and shared folder permissions suffice.
\end{corrige}

\begin{corrige}[Exercise 5]
\textbf{False.} Client/server \emph{facilitates} backup (centralised data) but does not \emph{eliminate} risk. Risks remain: server hardware failure (mitigated by RAID/cluster), ransomware encrypting server shares, administrator error, fire/theft in server room, backup media failure, untested restores. A robust backup strategy (3-2-1 rule, off-site, regular test restores) is still essential.
\end{corrige}

\begin{corrige}[Exercise 6]
Hybrid proposal: \textbf{Admin block (10 PCs) + Library/Staff room (5 PCs) = Client/Server.} Join to domain (Windows Server / Samba+LDAP). Centralised accounts for teachers/admin, shared admin files (exam papers, payroll), network printers, internet proxy/filtering, automated backups. \textbf{Computer Labs (2 $\times$ 12 PCs = 24 PCs) = Peer-to-Peer (Workgroup).} Students log on with generic accounts (``lab1user'', ``lab2user'') or local accounts. No need for individual domain accounts. File sharing between student PCs for group work can be ad-hoc. Teacher's station in each lab could be a domain member to push screen monitoring software. Justification: Admin block needs security, audit, individual accountability $\rightarrow$ client/server. Labs need simplicity, quick re-imaging, minimal login friction $\rightarrow$ peer-to-peer reduces domain traffic and login delays for 24 simultaneous student logons.
\end{corrige}

\coursec{Comparing Architectures and Choosing a Design Strategy}

In the previous sections we studied the two fundamental network architectures in detail: the \cle{peer-to-peer} (P2P) architecture, in which all machines are equal and share resources directly, and the \cle{client/server} architecture, in which powerful central servers provide services to client workstations. Knowing each architecture separately is not enough for the GCE examination or for real professional life. The real question an ICT specialist faces is: \emph{given a particular organisation, which architecture should be chosen, and why?} This final section gives you the tools to answer that question rigorously: a side-by-side comparison, decision criteria, a case-analysis method, worked examples drawn from Cameroonian situations, and the exam technique for writing a top-mark comparison essay.

\courssub{Side-by-Side Comparison of the Two Architectures}

Imagine you must choose between hiring ten motorbike taxis (\emph{bendskins}) or one bus to move people around a town. The motorbikes are cheap to start with, need no central garage, and each one works independently; but coordinating ten riders is chaotic, and if one rider carries the only copy of a parcel, losing that rider loses the parcel. The bus is expensive, needs a professional driver and a depot, but is organised, secure and can carry many more passengers. Neither is ``better'' in absolute terms: the correct choice depends on \emph{how many passengers, how much money, and how much organisation you have}. Network architectures behave in exactly the same way.

\begin{propriete}[Comparison of Peer-to-Peer and Client/Server Architectures]
The two architectures can be compared on seven standard criteria. This table is a summary of examinable facts; memorising it (with justifications) is essential for the GCE.
\begin{center}
\begin{tabularx}{\linewidth}{|l|X|X|}
\hline
\rowcolor{popblueL} \textbf{Criterion} & \textbf{Peer-to-Peer (P2P)} & \textbf{Client/Server (C/S)} \\
\hline
\textbf{Cost} & Low: no dedicated server, no server licence, no specialist staff. Uses ordinary PCs and a cheap switch. & High: dedicated server hardware, server operating system licences, possibly a paid administrator. \\
\hline
\textbf{Security} & Weak: each user manages security on his own machine; no central control of accounts or access rights. & Strong: centralised user accounts, passwords and access rights managed on the server. \\
\hline
\textbf{Scalability} & Poor: performance and manageability collapse beyond about 10 machines. & Excellent: hundreds or thousands of clients can be added; servers can be upgraded or multiplied. \\
\hline
\textbf{Administration} & Decentralised: every machine is configured individually; no central backup. & Centralised: one administrator manages users, data, backups and updates from the server. \\
\hline
\textbf{Performance} & Each machine shares its own resources, so a busy PC slows down its own user; degrades as the network grows. & Powerful dedicated server handles requests; stable performance even with many users. \\
\hline
\textbf{Reliability} & If a machine holding shared files is switched off, those files are unavailable; no single point of failure but many weak points. & Server is a single point of failure, but servers are built with redundancy (backups, RAID, UPS) and are rarely off. \\
\hline
\textbf{Size} & Small networks: typically 2 to 10 computers. & Medium to large networks: from about 10 computers to thousands. \\
\hline
\end{tabularx}
\end{center}
\end{propriete}

\begin{attention}[The ``10 Computers'' Rule Is a Guideline, Not a Law]
The limit of ``about 10 computers'' for peer-to-peer is a \emph{practical rule of thumb}, not a technical barrier. A P2P network of 12 machines will still function; it simply becomes increasingly difficult to manage (each user account must be created on every machine, backups are done machine by machine, and performance degrades). In the exam, treat 10 as the boundary beyond which client/server becomes strongly recommended, and always justify the boundary by \emph{manageability}, not by an imaginary technical impossibility.
\end{attention}

The figure below presents the same comparison visually, as a matrix. Green cells indicate the architecture that ``wins'' on that criterion; orange cells indicate the weaker side.

\begin{popfigure}
\begin{tikzpicture}[font=\small]
% Column headers
\node[fill=popink, text=white, minimum width=2.6cm, minimum height=0.8cm, align=center] at (0,6.4) {\textbf{Criterion}};
\node[fill=popgreen, text=white, minimum width=4.4cm, minimum height=0.8cm, align=center] at (3.6,6.4) {\textbf{Peer-to-Peer}};
\node[fill=popblue, text=white, minimum width=4.4cm, minimum height=0.8cm, align=center] at (8.1,6.4) {\textbf{Client/Server}};
% Row 1 cost
\node[fill=popblueL, minimum width=2.6cm, minimum height=0.7cm, align=center] at (0,5.6) {Cost};
\node[fill=popgreenL, minimum width=4.4cm, minimum height=0.7cm, align=center] at (3.6,5.6) {Low};
\node[fill=poporange!25, minimum width=4.4cm, minimum height=0.7cm, align=center] at (8.1,5.6) {High};
% Row 2 security
\node[fill=popblueL, minimum width=2.6cm, minimum height=0.7cm, align=center] at (0,4.8) {Security};
\node[fill=poporange!25, minimum width=4.4cm, minimum height=0.7cm, align=center] at (3.6,4.8) {Weak, per-machine};
\node[fill=popgreenL, minimum width=4.4cm, minimum height=0.7cm, align=center] at (8.1,4.8) {Strong, centralised};
% Row 3 scalability
\node[fill=popblueL, minimum width=2.6cm, minimum height=0.7cm, align=center] at (0,4.0) {Scalability};
\node[fill=poporange!25, minimum width=4.4cm, minimum height=0.7cm, align=center] at (3.6,4.0) {Poor (max $\sim$10 PCs)};
\node[fill=popgreenL, minimum width=4.4cm, minimum height=0.7cm, align=center] at (8.1,4.0) {Excellent (1000s)};
% Row 4 administration
\node[fill=popblueL, minimum width=2.6cm, minimum height=0.7cm, align=center] at (0,3.2) {Administration};
\node[fill=poporange!25, minimum width=4.4cm, minimum height=0.7cm, align=center] at (3.6,3.2) {Decentralised};
\node[fill=popgreenL, minimum width=4.4cm, minimum height=0.7cm, align=center] at (8.1,3.2) {Centralised};
% Row 5 performance
\node[fill=popblueL, minimum width=2.6cm, minimum height=0.7cm, align=center] at (0,2.4) {Performance};
\node[fill=poporange!25, minimum width=4.4cm, minimum height=0.7cm, align=center] at (3.6,2.4) {Degrades with size};
\node[fill=popgreenL, minimum width=4.4cm, minimum height=0.7cm, align=center] at (8.1,2.4) {Stable, dedicated};
% Row 6 reliability
\node[fill=popblueL, minimum width=2.6cm, minimum height=0.7cm, align=center] at (0,1.6) {Reliability};
\node[fill=poporange!25, minimum width=4.4cm, minimum height=0.7cm, align=center] at (3.6,1.6) {Files lost if PC off};
\node[fill=popgreenL, minimum width=4.4cm, minimum height=0.7cm, align=center] at (8.1,1.6) {High, with backups};
% Row 7 size
\node[fill=popblueL, minimum width=2.6cm, minimum height=0.7cm, align=center] at (0,0.8) {Typical size};
\node[fill=popgreenL, minimum width=4.4cm, minimum height=0.7cm, align=center] at (3.6,0.8) {2 to 10 computers};
\node[fill=popgreenL, minimum width=4.4cm, minimum height=0.7cm, align=center] at (8.1,0.8) {10 to 1000s};
\end{tikzpicture}
\legende{Comparison matrix of peer-to-peer and client/server architectures across the seven criteria. Green cells show the stronger architecture on each criterion.}
\end{popfigure}

\begin{aretenir}[Key Points of the Comparison]
\begin{itemize}
\item P2P wins only on \cle{cost} and \cle{simplicity}; client/server wins on \cle{security}, \cle{scalability}, \cle{administration}, \cle{performance} and \cle{reliability}.
\item The deciding factor is therefore rarely ``which is better'' but ``which fits the \textbf{size, budget and security needs} of this organisation''.
\item In an exam answer, every advantage quoted must be \emph{linked to the scenario} (e.g.\ ``the bank handles customers' money, so the strong centralised security of client/server is required'').
\end{itemize}
\end{aretenir}

\courssub{Decision Criteria: Which Questions to Ask}

Choosing an architecture is a \emph{decision under constraints}. A good network designer does not start from the technology; he starts from the organisation. Five questions, in order of importance, guide the choice:

\begin{enumerate}
\item \textbf{Number of users (size).} Fewer than about 10 computers $\Rightarrow$ P2P is feasible. More than 10, or several sites $\Rightarrow$ client/server is almost compulsory.
\item \textbf{Budget.} Can the organisation afford a server machine, server software licences and possibly a salaried administrator? A small business with a tight budget may be forced towards P2P.
\item \textbf{Security needs.} Does the organisation handle sensitive data (money, medical records, exam results, customer accounts)? If yes, centralised security is required $\Rightarrow$ client/server.
\item \textbf{Availability of a trained administrator.} Client/server networks need someone who can manage user accounts, backups and updates. If no such person exists (and none can be hired), P2P is safer.
\item \textbf{Growth plans.} An organisation that expects to grow (more staff, more branches) should invest in client/server from the start, because migrating later is costly and disruptive.
\end{enumerate}

Notice that the five questions can pull in opposite directions: a school may have 30 computers (pointing to client/server) but no budget and no administrator (pointing to P2P). In such conflicts, \textbf{security needs and size normally outweigh cost}, because a network that cannot protect its data or cannot be managed is useless, however cheap it was. The decision tree below organises the questions in the order in which they should be asked.

\begin{popfigure}
\begin{tikzpicture}[font=\small,
decision/.style={diamond, draw=popink, fill=popblueL, align=center, inner sep=1pt, aspect=2.2},
result/.style={rectangle, rounded corners, draw=popink, align=center, minimum height=0.8cm},
lab/.style={font=\footnotesize\itshape}]
\node[decision, align=center] (q1) at (0,0){More than 10\\computers?};
\node[decision, align=center] (q2) at (-4.2,-2.4){Sensitive data or\\high security need?};
\node[decision, align=center] (q3) at (4.2,-2.4){Budget for server\\and administrator?};
\node[result, fill=popgreenL, align=center] (p2p1) at (-6.4,-4.8){\textbf{Peer-to-Peer}\\small, simple needs};
\node[decision, align=center] (q4) at (-1.8,-4.8){Trained admin\\available?};
\node[result, fill=popgreenL, align=center] (cs1) at (6.4,-4.8){\textbf{Client/Server}\\large network};
\node[result, fill=popgreenL, align=center] (cs2) at (-1.8,-7.2){\textbf{Client/Server}\\security is priority};
\node[result, fill=popgreenL, align=center] (p2p2) at (-6.4,-7.2){\textbf{Peer-to-Peer}\\accept weaker security};
\draw[-latex, thick, popink] (q1) -- node[lab, above left]{No} (q2);
\draw[-latex, thick, popink] (q1) -- node[lab, above right]{Yes} (q3);
\draw[-latex, thick, popink] (q2) -- node[lab, above left]{No} (p2p1);
\draw[-latex, thick, popink] (q2) -- node[lab, above right]{Yes} (q4);
\draw[-latex, thick, popink] (q3) -- node[lab, right]{Yes} (cs1);
\draw[-latex, thick, popink] (q3.west) -- node[lab, above]{No} (q4.east);
\draw[-latex, thick, popink] (q4) -- node[lab, right]{Yes} (cs2);
\draw[-latex, thick, popink] (q4) -- node[lab, above left]{No} (p2p2);
\end{tikzpicture}
\legende{Decision tree for choosing a network architecture. Follow the questions from the top; the leaf reached gives the recommended architecture.}
\end{popfigure}

\begin{methode}[Four-Step Strategy for Choosing an Architecture]
When faced with a scenario (in an exam or in real life), proceed as follows.
\begin{enumerate}
\item \textbf{Read the scenario carefully} and extract the facts: number of computers, number of sites, type of activity, budget mentioned, staff skills, growth plans.
\item \textbf{Identify the needs} behind those facts: is security critical? Is the budget tight? Will the network grow? Is there an administrator?
\item \textbf{Match each need to the architecture} that satisfies it, using the comparison table (e.g.\ need for centralised security $\Rightarrow$ client/server; tiny budget and 5 PCs $\Rightarrow$ P2P).
\item \textbf{State and justify the choice}: name the recommended architecture and explain, criterion by criterion, \emph{why it fits this particular organisation}, and why the other one does not.
\end{enumerate}
\end{methode}

\courssub{Worked Case Studies}

\begin{exoresolu}[Case Study 1: A Cybercafé in Buea]
\textbf{Statement.} Mrs.\ Ngum opens a small cybercafé in Molyko, Buea, with 6 computers for customers and 1 computer at the reception desk. Her budget is limited, she has no employed network specialist, and the computers are used mainly for browsing, printing and typing documents. Advise her on the network architecture to adopt, justifying your answer.
\tcblower
\textbf{Solution (applying the 4-step method).}

\emph{Step 1 — Facts.} 7 computers in total, one room (single site), limited budget, no network specialist, simple uses (Internet, printing, file sharing).

\emph{Step 2 — Needs.} Low cost; simplicity of management; no sensitive data stored centrally; no planned expansion beyond a few machines.

\emph{Step 3 — Matching.} The network has fewer than 10 machines, so P2P is technically sufficient. The budget is small and there is no administrator, which rules out the expense of a dedicated server. Security needs are modest (each customer PC is independent; the reception PC can be password-protected individually).

\emph{Step 4 — Justified choice.} We recommend a \textbf{peer-to-peer architecture}. Justification: (i) with only 7 computers, P2P works well since it suits networks of up to about 10 machines; (ii) it is cheap — no server hardware, no server licence, no paid administrator — which respects the limited budget; (iii) administration is simple enough for the owner to handle (sharing the printer and folders on each PC); (iv) no sensitive centralised data is kept, so the weaker security of P2P is not a serious risk here. A client/server network would be an unjustified expense for this scenario.
\end{exoresolu}

\begin{exoresolu}[Case Study 2: A Microfinance Bank in Douala]
\textbf{Statement.} A microfinance bank has its head office in Akwa, Douala, and two branches in Bonabéri and Deïdo, with about 40 computers in total. The computers handle customer accounts, savings and loan records. The bank plans to open two more branches within three years and employs an IT officer. Recommend a network architecture, justifying your answer.
\tcblower
\textbf{Solution.}

\emph{Step 1 — Facts.} About 40 computers across 3 sites; financial (sensitive) data; employed IT officer; clear growth plans.

\emph{Step 2 — Needs.} Strong centralised security (customers' money and account data); centralised data storage so that a customer can be served at any branch; reliable backups; ability to grow; professional administration.

\emph{Step 3 — Matching.} 40 computers is far beyond the practical limit of P2P ($\sim$10 machines). Financial data demands the centralised accounts, access rights and backups that only a server provides. The presence of an IT officer satisfies the administration requirement, and the growth plan demands scalability.

\emph{Step 4 — Justified choice.} We recommend a \textbf{client/server architecture}. Justification: (i) \emph{size} — 40 computers on 3 sites cannot be managed peer-to-peer; (ii) \emph{security} — customer financial records must be protected by centralised user accounts and access rights on a server; (iii) \emph{administration} — the IT officer can manage all accounts, data and backups centrally; (iv) \emph{reliability} — the server can be protected with backups and a UPS so account data is never lost; (v) \emph{scalability} — new branches and workstations can be added easily as the bank expands. The high cost of client/server is justified here because the bank's activity depends entirely on the security and availability of its data.
\end{exoresolu}

\begin{attention}[Common Mistake: Always Recommending Client/Server]
Many candidates believe client/server is ``the modern, correct answer'' and recommend it for every scenario. This loses marks. For a small organisation (a 5-PC school office, a small cybercafé, a home network), client/server is \emph{wrong} because its cost and complexity are unjustified: there is no budget for a server and no administrator to run it. Cost and size very often make \textbf{peer-to-peer the correct answer}. Always let the scenario's constraints decide, never a personal preference.
\end{attention}

\begin{savaistu}[Did You Know? — Mixed Networks Exist]
Real organisations sometimes combine the two architectures in a \cle{hybrid} (mixed) network: a central server holds the important shared data and user accounts (client/server part), while workstations also share folders or a printer directly among themselves (peer-to-peer part). This gives small organisations some centralised security without paying for a fully managed server environment. For the GCE, however, you are expected to master the two pure architectures and to choose between them; mention hybrid networks only as an enriching remark, never as a way to avoid making a justified choice.
\end{savaistu}

\courssub{Exam Technique: Writing the GCE Comparison Essay}

GCE Advanced Level questions on this topic typically ask you to ``compare peer-to-peer and client/server networks'' or to ``advise an organisation, justifying your choice''. A well-structured answer has three parts:

\begin{enumerate}
\item \textbf{Introduction (2--3 lines).} Define both architectures briefly and announce that the choice depends on the organisation's needs.
\item \textbf{Body: the comparison.} Either present a comparison table (criterion by criterion, as in the Property box above) or develop each criterion in a short paragraph using the \cle{PEEL} method.
\item \textbf{Conclusion: justified recommendation.} Name the chosen architecture and justify it by linking at least three criteria \emph{to the facts of the scenario}.
\end{enumerate}

\begin{methode}[The PEEL Method for a Comparison Essay]
Use PEEL for each paragraph of the body and for the conclusion:
\begin{enumerate}
\item \textbf{P — Point:} state one clear idea (``Client/server offers stronger security than peer-to-peer'').
\item \textbf{E — Evidence:} give the fact (``user accounts, passwords and access rights are managed centrally on the server'').
\item \textbf{E — Explanation:} explain why it matters (``so an administrator controls exactly who can read or modify each file, which is impossible when every PC manages its own security'').
\item \textbf{L — Link:} connect back to the question or scenario (``therefore, for a bank handling customers' savings, client/server is the appropriate choice'').
\end{enumerate}
A paragraph without the Link is an unfinished paragraph: it describes, but does not \emph{answer the question}.
\end{methode}

\begin{attention}[Exam Pitfall: Describing Without Justifying]
The most frequent reason for losing marks is writing a beautiful description of client/server (or of both architectures) without ever explaining \emph{why it fits the given case}. Examiners reward \textbf{application}, not recitation. The sentence ``client/server has centralised security'' earns little; the sentence ``\emph{because} the microfinance bank stores customers' account records, the centralised security of client/server is required to control access to those records'' earns full marks. Every point must carry the words ``because'' or ``therefore'' tied to the scenario.
\end{attention}

\begin{savaistu}[Did You Know?]
In Cameroon, the choice between architectures is visible everywhere. Small cybercafés and many primary-school computer rooms run peer-to-peer networks because of cost, while banks, mobile-money operators such as MTN Mobile Money and Orange Money, and state bodies like ANTIC rely on large client/server infrastructures with centralised data centres — because when millions of transactions are involved, centralised security and reliability are not optional.
\end{savaistu}

\begin{exercice}[Applying the Method]
For each organisation below, recommend an architecture and justify your choice on at least three criteria.
\begin{enumerate}
\item A private nursery school in Bamenda with 4 computers used by the secretary and teachers for records and printing; no IT staff; very small budget.
\item A regional hospital in Yaoundé with 120 computers across 5 buildings, storing patients' medical records, with a team of 3 IT technicians and plans to computerise the pharmacy next year.
\item A family home in Limbe with 3 computers and one printer shared between parents and children.
\end{enumerate}
\end{exercice}

\begin{corrige}[Exercise 1]
\begin{enumerate}
\item \textbf{Peer-to-peer.} Size: 4 computers is well within the P2P limit of about 10. Cost: no money for a server or licences, and P2P needs only the existing PCs and a switch. Administration: with no IT staff, the simple per-machine management of P2P is the only realistic option. Security needs are modest, so P2P's weak security is acceptable.
\item \textbf{Client/server.} Size: 120 computers on 5 buildings is far beyond P2P. Security: medical records are highly sensitive and require centralised accounts and access rights. Administration: the 3 technicians can run the servers, backups and updates centrally. Scalability: the planned pharmacy computerisation is easily added to a client/server network.
\item \textbf{Peer-to-peer.} Size: 3 computers; cost: a home cannot justify a server; administration: family members can share the printer and folders directly; no sensitive centralised data, so strong security is unnecessary.
\end{enumerate}
\end{corrige}

\begin{aretenir}[To Conclude the Chapter]
\begin{itemize}
\item There is no universally ``best'' architecture: the best choice is the one that \textbf{fits the size, budget, security needs, staffing and growth plans} of the organisation.
\item P2P: cheap and simple, but limited to about 10 machines, with weak, decentralised security. Client/server: costly and needing an administrator, but secure, scalable, centralised, performant and reliable.
\item In the exam, always \emph{read the scenario, identify needs, match needs to criteria, and justify} — using PEEL paragraphs and ending with a recommendation explicitly linked to the case.
\end{itemize}
\end{aretenir}

\newpage

\coursec{Integration activity}

\coursec{Network Design Strategies}

\courssub{Lesson 13: Network Architecture — Integration Activity}

\courssub{Aims of the Activity}

By the end of this integration activity, you should be able to:

\begin{itemize}
\item analyse a real-life situation and identify the \cle{needs} of different users of a network;
\item distinguish the contexts in which a \cle{peer-to-peer} architecture or a \cle{client/server} architecture is appropriate;
\item draw clear, labelled \cle{network diagrams} for a proposed design;
\item justify a design decision using the criteria of \cle{cost}, \cle{security}, \cle{scalability} and \cle{ease of administration};
\item communicate a technical recommendation in simple language to a non-specialist decision-maker (a school principal);
\item defend a technical choice orally and respond to counter-arguments, as required in the GCE Advanced Level problem-solving approach.
\end{itemize}

This activity mobilises all the competences developed in this chapter: description of architectures, comparison of their strengths and weaknesses, and selection of a network design strategy adapted to constraints.

\courssub{Situation: Designing the Network for ``Lyc\'ee de Molyko''}

The Lyc\'ee de Molyko, Buea, has just received funding from an education support project to equip two rooms with computers and to connect them into local area networks (LANs). The principal, who is not a computer specialist, has asked the school's ICT club (your class) to advise her before any money is spent.

\textbf{Room A --- The Student Computer Lab.}
\begin{itemize}
\item \textbf{Equipment:} 10 identical desktop computers, one black-and-white laser printer, one switch, one router with an Internet connection from a local ISP (e.g. CAMTEL or MTN).
\item \textbf{Users:} Form Five and Lower Sixth students, working in pairs during practical ICT lessons (word processing, spreadsheets, introduction to programming).
\item \textbf{Supervision:} one teacher supervises the room; there is \textbf{no full-time network administrator}.
\item \textbf{Budget:} very limited --- the funding covers the 10 computers and little else. Buying an expensive dedicated server machine and a server operating system licence is almost impossible.
\item \textbf{Needs:} students must share the printer, share files (e.g.\ a worksheet prepared by the teacher), and access the Internet. Students do not store sensitive data; work is usually printed or copied to flash drives at the end of the lesson.
\end{itemize}

\textbf{Room B --- The Administration Office.}
\begin{itemize}
\item \textbf{Equipment:} 15 staff computers used by the bursar, the principal's secretary, the vice-principals and the discipline masters.
\item \textbf{Data handled:} student records (marks, report cards), fee payments, staff salaries and official correspondence. These data are \textbf{confidential}: loss, unauthorised modification or leakage would have serious consequences for the school.
\item \textbf{Needs:} data must be stored \textbf{centrally}, protected by \textbf{user accounts and passwords} with different access rights (e.g.\ only the bursar modifies fee records), and \textbf{backed up automatically every night}. The office expects to grow: two new secretaries will be hired next year, and the school plans to connect the two rooms together later.
\item \textbf{Budget:} the project can finance one powerful machine to act as a server and a trained staff member can be designated to manage it part-time.
\end{itemize}

The principal's exact words were: \emph{``I do not want to buy the most expensive solution everywhere, but I cannot accept that our students' records are lost or read by just anyone. Explain to me what we should install in each room, and why.''}

\courssub{Tasks}

Work in groups of four. Produce your answers on paper (or in a word processor) and be ready to defend them in the class debate.

\begin{enumerate}
\item \textbf{Needs analysis.} For each room, list at least four concrete needs (hardware sharing, data sharing, security, backup, growth, budget). Present your answer as a two-column table.

\item \textbf{Choice of architecture.} For each room, state which architecture you recommend (peer-to-peer or client/server) and give \textbf{two} reasons per room, directly linked to the needs you listed in question 1.

\item \textbf{Diagrams.} Draw a labelled diagram of each proposed network. Show the computers, the printer(s), the switch, the router/Internet link, and (where applicable) the server. Indicate on each diagram where the data are stored and who controls access.

\item \textbf{Comparison and justification.} Copy and complete the comparison table below, filling the last two columns with short, precise arguments for Room A and Room B.

\begin{center}
\begin{tabularx}{\linewidth}{|l|X|X|}
\hline
\rowcolor{popblueL} \textbf{Criterion} & \textbf{Room A: Lab (your choice)} & \textbf{Room B: Administration (your choice)} \\
\hline
Cost of installation & & \\
\hline
Security of data & & \\
\hline
Scalability (adding machines) & & \\
\hline
Ease of administration & & \\
\hline
Main risk accepted & & \\
\hline
\end{tabularx}
\end{center}

\item \textbf{Recommendation letter.} Write a one-page recommendation to the principal. It must: (a) state clearly what to install in each room; (b) justify each choice in at most three sentences per room, in non-technical language; (c) mention one future improvement (e.g.\ interconnecting the two LANs, adding a firewall or an uninterruptible power supply).

\item \textbf{Class debate.} Each group presents its recommendation in three minutes. Be ready to answer challenges such as: \emph{``Why not put a server in the lab too?''}, \emph{``What happens in Room B if the server breaks down?''}, \emph{``Is peer-to-peer really safe enough for students?''}
\end{enumerate}

\begin{attention}[Common Traps to Avoid]
\begin{itemize}
\item Do not choose client/server ``because it is better'' in general: it is better \emph{only where its advantages are needed}. A design must match needs and budget.
\item Do not confuse the \textbf{Internet connection} with the \textbf{network architecture}: both rooms can have Internet access whatever architecture is chosen.
\item In a peer-to-peer network, every machine can act as both client and server; do not draw a central server in Room A if you chose peer-to-peer.
\item Justify with \emph{criteria} (cost, security, scalability, administration), not with vague words like ``modern'' or ``fast''.
\end{itemize}
\end{attention}

\courssub{Model Answers}

\begin{exemplebox}[1. Needs Analysis]
\begin{center}
\begin{tabularx}{\linewidth}{|X|X|}
\hline
\rowcolor{popblueL} \textbf{Room A --- Student Lab} & \textbf{Room B --- Administration} \\
\hline
Share one printer among 10 PCs & Central, unique storage of student records and fees \\
\hline
Share teacher's worksheets and small files & Strong security: accounts, passwords, access rights per user \\
\hline
Internet access for research lessons & Automatic nightly backup of all data \\
\hline
Very low cost; no dedicated server affordable & Possibility to add machines (2 new secretaries soon) \\
\hline
Simple management by one teacher, no administrator & Part-time administrator available; server budgeted \\
\hline
No sensitive data stored permanently & Confidentiality is critical (fees, marks, salaries) \\
\hline
\end{tabularx}
\end{center}
\end{exemplebox}

\begin{exemplebox}[2. Choice of Architecture]
\begin{itemize}
\item \textbf{Room A: Peer-to-peer.}
      \begin{enumerate}
      \item The budget cannot cover a dedicated server and its licence; peer-to-peer uses the existing 10 PCs and a cheap switch.
      \item The needs are limited to sharing a printer, small files and Internet access among a small, fixed number of machines (10), with no sensitive data --- exactly the situation peer-to-peer handles well.
      \end{enumerate}
\item \textbf{Room B: Client/Server.}
      \begin{enumerate}
      \item Confidential records require centralised user accounts, passwords and access rights, which only a server administers reliably.
      \item Automatic central backup and easy addition of new machines (scalability) are natural in client/server, and the budget provides a server and a part-time administrator.
      \end{enumerate}
\end{itemize}
\end{exemplebox}

\begin{exemplebox}[3. Network Diagrams]
\begin{popfigure}
\begin{center}
\begin{tikzpicture}[scale=0.85, every node/.style={font=\small}]
% --- Room A: Peer-to-Peer ---
\node[font=\bfseries] at (2.5,4.0) {Room A: Peer-to-Peer Lab};
% Switch
\node[draw, fill=popblueL, rounded corners, minimum width=1.8cm, minimum height=0.9cm] (swA) at (2.5,2.0) {Switch};
% PCs (5 shown, representing 10)
\foreach \x/\n in {-1.8/PC1, -0.4/PC2, 1.0/PC3, 2.4/PC4, 3.8/PC5, 5.2/PC6, 6.6/PC7} {
  \node[draw, fill=white, rounded corners, minimum width=0.85cm, minimum height=0.55cm] (\n) at (\x,3.6) {\n};
  \draw[popline] (\n) -- (swA);
}
\node[font=\footnotesize] at (7.5,3.6) {$\dots$ PC10};
% Printer
\node[draw, fill=popgreenL, rounded corners, minimum width=1.4cm, minimum height=0.6cm] (prA) at (-1.8,0.4) {Printer};
\draw[popline] (prA) -- (swA);
% Router / Internet
\node[draw, fill=popgreenL, rounded corners, minimum width=1.4cm, minimum height=0.6cm] (rtA) at (6.8,0.4) {Router};
\draw[popline] (rtA) -- (swA);
\node[draw, ellipse, fill=popblue!10, minimum width=1.6cm, minimum height=0.9cm] (net) at (6.8,-1.0) {Internet};
\draw[popline] (rtA) -- (net);
% Annotation
\node[align=center, font=\footnotesize, text width=3.0cm, text=popdark] at (-3.2,2.0) {Each PC shares its own files; data stored locally on each PC};
% --- Room B: Client/Server ---
\node[font=\bfseries] at (12.5,4.0) {Room B: Client/Server Office};
% Switch
\node[draw, fill=popblueL, rounded corners, minimum width=1.8cm, minimum height=0.9cm] (swB) at (12.5,2.0) {Switch};
% Server
\node[draw, fill=poporange!40, rounded corners, minimum width=2.4cm, minimum height=1.0cm] (srv) at (12.5,3.8) {Server};
\node[align=center, font=\footnotesize, text width=2.5cm, text=popdark] at (15.0,3.8) {User accounts,\\ access rights,\\ central data,\\ nightly backup};
\draw[popline, thick] (srv) -- (swB);
% Clients (4 shown, representing 15)
\foreach \x/\n in {9.0/C1, 10.4/C2, 11.8/C3, 13.2/C4, 14.6/C5, 16.0/C6} {
  \node[draw, fill=white, rounded corners, minimum width=0.8cm, minimum height=0.55cm] (\n) at (\x,0.2) {\n};
  \draw[popline] (\n) -- (swB);
}
\node[font=\footnotesize] at (17.0,0.2) {$\dots$ C15};
% Printer (optional, shared via server)
\node[draw, fill=popgreenL, rounded corners, minimum width=1.2cm, minimum height=0.55cm] (prB) at (12.5,-1.0) {Printer};
\draw[popline] (prB) -- (swB);
\end{tikzpicture}
\end{center}
\legende{Proposed designs: peer-to-peer for the lab (left), client/server for the administration (right).}
\end{popfigure}
\end{exemplebox}

\begin{exemplebox}[4. Completed Comparison Table]
\begin{center}
\begin{tabularx}{\linewidth}{|l|X|X|}
\hline
\rowcolor{popblueL} \textbf{Criterion} & \textbf{Room A: Peer-to-peer} & \textbf{Room B: Client/server} \\
\hline
Cost of installation & Very low: no server, no server OS licence; uses existing 10 PCs and one switch & Higher: one powerful server, server OS licence (e.g. Windows Server / Linux), backup device (NAS/tape) \\
\hline
Security of data & Low: each user manages own PC; no central access control. Acceptable because no sensitive data stored permanently & High: centralised user accounts, passwords, granular access rights (read/write/execute) per service/folder \\
\hline
Scalability (adding machines) & Limited: beyond $\approx$10 machines, manual configuration on each PC becomes confusing and error-prone & Good: new clients simply join domain/workgroup; accounts and policies applied centrally from server \\
\hline
Ease of administration & Simple for one teacher (no admin skills needed), but settings (IP, sharing) repeated on every PC & Centralised: one part-time administrator manages all users, shares, backups, updates from the server console \\
\hline
Main risk accepted & A student may accidentally delete another's shared file; no central backup. Risk accepted due to low data criticality & Server hardware failure stops the whole office; mitigated by nightly backups, RAID disks, and a UPS \\
\hline
\end{tabularx}
\end{center}
\end{exemplebox}

\begin{exemplebox}[5. Recommendation Letter (Model)]
\textbf{To: The Principal, Lyc\'ee de Molyko}\\
\textbf{From: ICT Club, Upper Sixth}\\
\textbf{Date: \today}\\
\textbf{Subject: Network Architecture Recommendation for Computer Lab and Administration Office}

Dear Madam Principal,

Following our analysis of the needs and constraints of both rooms, we recommend two different network architectures, each adapted to its specific context.

\textbf{Room A (Student Computer Lab): Peer-to-Peer Network.}
The ten computers will be connected to a single switch, which also links the printer and the router for Internet access. In this architecture, every computer is equal: each can share its own folders and use the printer directly. This solution requires no extra server hardware and no expensive server operating system licence, keeping the cost to the absolute minimum. Since students do not store confidential data and files are saved on flash drives at the end of each lesson, the lower security and lack of central backup are acceptable risks. The supervising teacher can manage the network without specialised training.

\textbf{Room B (Administration Office): Client/Server Network.}
One powerful computer will be designated as the server, running a server operating system. All 15 staff computers (clients) connect to the switch and access data \emph{only} through the server. The server enforces user accounts, passwords, and precise access rights (e.g. the bursar alone can modify fee records). It performs automatic backups every night, guaranteeing that no record is lost even if a client machine fails or a user makes a mistake. Although this solution costs more (server hardware, licence, backup media), it is the \emph{only} architecture that provides the confidentiality, integrity and availability your administrative data demand.

\textbf{Future Improvement.}
When the school is ready, the two LANs can be interconnected via the existing router. We strongly recommend adding a \cle{firewall} between the administration server and the student lab (or the Internet) to block unauthorised access attempts. The administration server should also be protected by an \cle{Uninterruptible Power Supply (UPS)} to prevent data corruption during Buea's frequent power fluctuations.

We are available to present this proposal and answer any questions at your convenience.

Respectfully,
The ICT Club (Upper Sixth)
\end{exemplebox}

\begin{exemplebox}[6. Expected Debate Points]
\begin{itemize}
\item \textbf{``Why not put a server in the lab too?''} A server would cost 3--5 times more (hardware + OS licence + CALs) for benefits (central backup, strict access control) that are not needed because the lab handles no confidential data and has no budget for an administrator.
\item \textbf{``What happens in Room B if the server breaks down?''} Work pauses because clients cannot access central data. However, \emph{no data is lost} thanks to the nightly backup. Downtime is minimised by: (a) a UPS protecting against power cuts, (b) RAID disks on the server (survives one disk failure), (c) a maintenance contract with a local ICT firm (e.g. CAMTEL Business or a certified partner) for 4-hour on-site response.
\item \textbf{``Is peer-to-peer really safe enough for students?''} Yes, because the \emph{threat model} is different. In the lab, the risk is accidental file deletion or a student reading a neighbour's homework --- inconvenient but not disastrous. In the administration, the risk is leakage of salary sheets or alteration of exam marks --- legally and reputationally disastrous. Security investment must be proportional to the value of the asset protected.
\end{itemize}
\end{exemplebox}

\begin{savaistu}[Network Architecture in Cameroon]
In Cameroon, most cybercaf\'es and small school labs historically used peer-to-peer (workgroup) networks because they require no server licence and can be set up by anyone with basic Windows knowledge. However, with the rise of \emph{e-government} services (e.g. \texttt{www.minesec.gov.cm} for exam registration, \texttt{www.antic.cm} for cybersecurity reporting) and mobile money integration (MTN MoMo, Orange Money APIs), even small organisations are moving to client/server or cloud-based solutions to meet data protection requirements (Law No. 2010/012 on Cybersecurity and Cybercrime). The \cle{Agence Nationale des Technologies de l'Information et de la Communication (ANTIC)} recommends centralised logging and access control for any system handling personal data --- a strong argument for client/server in administration offices.
\end{savaistu}

\newpage

\coursec{Methods and worked examples}

\coursec{Network Design Strategies}

\courssub{Introduction to Network Architecture}

\begin{definition}[Network Architecture]
Network architecture refers to the \cle{structural and logical design} of a computer network, defining how devices (nodes) are interconnected, how tasks are distributed among them, and how resources (files, printers, applications, Internet access) are shared and managed.
\end{definition}

\textbf{Key Elements:}\par
\begin{itemize}
    \item \textbf{Node:} Any device (PC, server, printer, IP phone, access point) connected to the network.
    \item \textbf{Link:} The physical (copper UTP/STP, fibre optic, radio) or logical connection between nodes.
    \item \textbf{Resource:} Any asset shared on the network (data files, databases, hardware peripherals, software applications).
    \item \textbf{Role:} The function a node plays in resource interaction: \cle{Client} (requests services), \cle{Server} (provides services), or \cle{Peer} (both requests and provides).
\end{itemize}

\begin{propriete}[Fundamental Design Decision]
The choice between \cle{Peer-to-Peer (P2P)} and \cle{Client/Server (C/S)} architecture is the primary strategic decision in LAN design. It dictates hardware procurement, software licensing, administrative overhead, security posture, scalability, and compliance with regulations (e.g., Cameroon Law No. 2010/012 on Cybersecurity).
\end{propriete}

\begin{savaistu}[Cameroon Context]
In Cameroon, most cybercafés (\textit{cyber}) and very small businesses (e.g., a boutique in Mokolo market with 3 PCs sharing a printer) operate de facto P2P networks. However, any structure handling sensitive data (banks, ministries, hospitals, schools with $>10$ PCs) \emph{must} adopt Client/Server to comply with Law No. 2010/012 on Cybersecurity and ANTIC guidelines on access control, audit trails, and data integrity.
\end{savaistu}

\courssub{Peer-to-Peer Architecture}

\begin{definition}[Peer-to-Peer (P2P) Architecture]
A network model where all nodes (\cle{peers}) have equivalent capabilities and responsibilities. Each workstation acts as both a \cle{client} (requesting resources) and a \cle{server} (providing resources like shared folders or printers). There is \textbf{no dedicated server hardware} and \textbf{no centralised authentication database}.
\end{definition}

\begin{propriete}[Characteristics of P2P]
\begin{itemize}
    \item \textbf{Symmetric Roles:} Every computer can share and access resources; no functional distinction between nodes.
    \item \textbf{Decentralised Administration:} Each user manages their own machine (local user accounts, sharing permissions, OS updates, antivirus).
    \item \textbf{Local Security:} Authentication is local (SAM database on each Windows PC). No Single Sign-On (SSO); a user needs an account on every machine they access.
    \item \textbf{Workstation OS:} Runs standard desktop OS (Windows 10/11 Pro, Ubuntu Desktop), not a Server OS.
    \item \textbf{Scale Limit:} Practical limit \textbf{$\approx 10$ nodes} (GCE syllabus guideline). Performance and management complexity grow exponentially beyond this.
\end{itemize}
\end{propriete}

\begin{exemplebox}[Typical P2P Scenario]
A graphic design startup in Yaoundé (Ngousso) with 6 PCs connected to a switch. The printer is plugged into PC\#1 via USB. Project files are in a shared folder on PC\#3. Each PC has its own local user accounts. This is a classic P2P setup.
\end{exemplebox}

\begin{attention}[P2P Traps in Exams]
\begin{itemize}
    \item \textbf{Do not confuse} ``Workgroup'' (Windows term for P2P logical grouping) with ``Domain'' (Client/Server logical grouping).
    \item \textbf{Do not write} ``P2P has no security''. Write ``P2P has \emph{decentralised, inconsistent} security difficult to enforce at scale''.
    \item \textbf{Never recommend} P2P for $>10$ nodes, or where centralised backup, Group Policy, or regulatory compliance (ANTIC) are required.
\end{itemize}
\end{attention}

\courssub{Client/Server Architecture}

\begin{definition}[Client/Server (C/S) Architecture]
A network model with \cle{asymmetric roles}: dedicated \cle{servers} provide resources and services to multiple \cle{clients} (workstations). Servers run a \cle{Server Operating System} (Windows Server, Linux Server) and host centralised services (Active Directory, File Storage, Databases, Printing, DHCP, DNS). Authentication is centralised via a \cle{Domain Controller}.
\end{definition}

\begin{propriete}[Key Server Roles in a Modern LAN]
\begin{center}
\begin{tabularx}{\linewidth}{|l|X|}
\hline
\rowcolor{popblueL}
\textbf{Role} & \textbf{Function} \\
\hline
\textbf{Domain Controller (DC)} & Active Directory Domain Services (AD DS), DNS, DHCP, Kerberos/KDC, Group Policy, Certificate Services. \textbf{Heart of C/S}. \\
\hline
\textbf{File Server} & Centralised storage (SMB/NFS), DFS Namespaces, FSRM (Quotas, File Screens), RAID redundancy, Shadow Copies (Previous Versions). \\
\hline
\textbf{Print Server} & Central spooling, driver repository, printer pooling, auditing, driver isolation. \\
\hline
\textbf{Application/Database Server} & Hosts ERP, CRM, School Mgt System (SQL Server, PostgreSQL, MySQL). Requires high RAM/IOPS (NVMe SSD, RAID 10). \\
\hline
\textbf{Web/Proxy Server} & IIS/Apache/Nginx, Reverse Proxy, Forward Proxy (Cache, Content Filter), SSL Offloading/Inspection. \\
\hline
\textbf{NPS / RADIUS} & Network Policy Server for 802.1X (Wired/WiFi), VPN, MAB (MAC Auth Bypass). Centralised AAA (Authentication, Authorisation, Accounting). \\
\hline
\textbf{WDS / MDT} & Windows Deployment Services / Microsoft Deployment Toolkit for automated OS imaging and deployment. \\
\hline
\end{tabularx}
\end{center}
\end{propriete}

\begin{exemplebox}[Client/Server in a Cameroonian School]
GBHS Mfoundi: 65 PCs across 3 blocks. Domain Controller manages all accounts (SSO). File Server stores exam papers (backed up nightly to NAS). Print Server manages 5 heavy-duty printers. NPS handles 802.1X for Staff/Student WiFi. Proxy (FortiGate) filters Internet. All managed by 1 sysadmin via GPOs from a single console.
\end{exemplebox}

\begin{attention}[Single Point of Failure (SPOF) Mitigation]
In basic C/S, the Server is a SPOF. \textbf{Exam Answer:} ``Mitigate with RAID (disk failure), Dual Hot-Plug PSU (power failure), Clustering/HA (service failure), UPS/Generator (mains failure), Offsite/Cloud Backup (site disaster).''
\end{attention}

\courssub{Comparing Architectures and Choosing a Design Strategy}

\textbf{Methods and Worked Examples}\par

This section presents the essential methodological skills for analysing, comparing, and selecting network architectures. Each method is structured as a step-by-step procedure followed by a GCE-style worked examination question with a detailed solution. Master these patterns; they are the backbone of Paper 2 (Essay) and Paper 3 (Practical) questions on network design.

\begin{methode}[Identifying the Architecture Type from a Scenario]
\textbf{Objective:} Determine whether a described network is Peer-to-Peer (P2P) or Client/Server (C/S) based on functional clues.

\begin{enumerate}
    \item \textbf{Locate the resource control:} Ask ``Where are the shared files, printers, database, or user accounts stored?''
    \begin{itemize}
        \item If stored on \emph{one dedicated machine} that \emph{only serves} others $\rightarrow$ \cle{Client/Server}.
        \item If stored on \emph{multiple ordinary workstations} that \emph{also act as users} $\rightarrow$ \cle{Peer-to-Peer}.
    \end{itemize}
    \item \textbf{Check for a dedicated Server OS:} Mention of Windows Server, Linux Server (Ubuntu Server, CentOS/Rocky/Alma), Active Directory, Domain Controller $\rightarrow$ \cle{Client/Server}.
    \item \textbf{Check for centralised authentication:} ``Users log in to the domain'', ``Single sign-on'', ``Centralised password policy'' $\rightarrow$ \cle{Client/Server}. ``Each PC has its own local accounts'' $\rightarrow$ \cle{Peer-to-Peer}.
    \item \textbf{Count the nodes:} Typically $< 10$ nodes $\rightarrow$ P2P feasible. $> 10$ nodes $\rightarrow$ C/S strongly recommended (syllabus guideline).
    \item \textbf{Identify the role symmetry:} In P2P, every node is \emph{both} client and server (symmetric roles). In C/S, roles are \emph{asymmetric} (dedicated server, dedicated clients).
\end{enumerate}
\textbf{Reflex:} Draw a quick sketch. A star with a distinct central node labelled ``Server'' = C/S. A mesh or star with a switch but all PCs labelled ``Workstation'' = P2P.
\end{methode}

\begin{exoresolu}[Architecture Identification at \textit{Ngousso Enterprises}]
\textbf{Statement:}
Ngousso Enterprises, a small graphic design startup in Yaoundé, has 6 desktop computers connected via a switch. They share a high-end colour printer connected to PC\#1 via USB. Project files are stored in a shared folder on PC\#3. Each computer manages its own local user accounts and passwords. There is no dedicated server hardware. The owner wants to know the architecture type and one major limitation of this setup as the business grows to 15 staff.

\tcblower
\textbf{Detailed Solution:}

\noindent\textbf{Step 1: Identify resource location.}
\begin{itemize}
    \item Printer $\rightarrow$ attached to PC\#1 (a workstation).
    \item Project files $\rightarrow$ shared folder on PC\#3 (a workstation).
    \item User accounts $\rightarrow$ local on each PC.
\end{itemize}
No dedicated machine acts solely as a provider.

\noindent\textbf{Step 2: Check roles.}
All 6 computers are used by designers for daily work (client role) \emph{and} provide resources (server role). Roles are symmetric.

\noindent\textbf{Step 3: Conclusion.}
This is a \textbf{Peer-to-Peer (P2P) Architecture}.

\noindent\textbf{Step 4: Major limitation for growth to 15 staff.}
\textbf{Centralised Administration \& Security.} With 15 PCs in P2P:
\begin{itemize}
    \item \textbf{Account Management Nightmare:} Creating/deleting a user requires visiting 15 machines. No Group Policy.
    \item \textbf{Inconsistent Security:} Password policies, firewall rules, antivirus updates cannot be enforced centrally.
    \item \textbf{Backup Complexity:} Critical data scattered across 15 ``C:\'' drives. No single backup target.
    \item \textbf{Performance Bottleneck:} PC\#1 and PC\#3 (resource hosts) will slow down under heavy simultaneous access from 14 other users because they are not optimised server hardware (RAM, CPU, RAID disks, NIC teaming).
\end{itemize}
\textbf{Recommendation:} Migrate to Client/Server with a dedicated file/print server and Active Directory domain controller.
\end{exoresolu}

\begin{methode}[Comparing Architectures using a Criteria Grid (Exam Table Technique)]
\textbf{Objective:} Produce a structured comparison table guaranteed to earn full marks in a ``Compare P2P and Client/Server'' essay question (often 6--8 marks).

\begin{enumerate}
    \item \textbf{Draw a table with 3 columns:} \textbf{Criterion}, \textbf{Peer-to-Peer}, \textbf{Client/Server}.
    \item \textbf{Use exactly these 6--7 standard criteria (syllabus core):}
    \begin{enumerate}
        \item \textbf{Centralised Administration / Management}
        \item \textbf{Security / Access Control / Authentication}
        \item \textbf{Scalability / Number of Nodes Supported}
        \item \textbf{Cost (Hardware, Software, Maintenance)}
        \item \textbf{Performance / Resource Optimisation}
        \item \textbf{Reliability / Fault Tolerance / Backup}
        \item \textbf{Technical Expertise Required}
    \end{enumerate}
    \item \textbf{Fill cells with contrasting keywords (not essays).} Use ``Decentralised'' vs ``Centralised'', ``Low'' vs ``High'', ``Local accounts'' vs ``Domain/Active Directory'', ``Workstation OS'' vs ``Server OS''.
    \item \textbf{Add a concluding ``Best Fit'' row or sentence:} ``P2P: Small offices ($<10$), low budget, low security needs. C/S: Medium/Large, high security, growth planned.''
\end{enumerate}
\textbf{Examiner Tip:} Never write ``P2P is cheap, C/S is expensive'' without qualification. Say ``P2P: Low initial cost (no server HW/OS licences). C/S: High initial cost (dedicated server, Server OS CALs) but lower TCO at scale.''
\end{methode}

\begin{exoresolu}[Structured Comparison for a School Network Proposal]
\textbf{Statement:}
The ICT department of Government Bilingual High School (GBHS) Mfoundi is designing a new network for 50 computers in the computer lab, 10 in the admin block, and 5 in the library. They must choose between Peer-to-Peer and Client/Server.
\begin{enumerate}
    \item Draw a comparison table using \textbf{six} relevant criteria.
    \item State the recommended architecture with \textbf{two} justifications linked to the school context.
\end{enumerate}

\tcblower
\textbf{Detailed Solution:}

\noindent\textbf{(i) Comparison Table}

\begin{center}
\begin{tabularx}{\linewidth}{|l|X|X|}
\hline
\rowcolor{popblueL}
\textbf{Criterion} & \textbf{Peer-to-Peer} & \textbf{Client/Server} \\
\hline
\textbf{1. Administration} & Decentralised; each PC managed locally. No Group Policy. & Centralised via Domain Controller (Active Directory); Group Policies push settings to all 65 PCs instantly. \\
\hline
\textbf{2. Security / Authentication} & Local user accounts per PC; weak password enforcement; no central audit trail. & Single sign-on (domain accounts); central password policy (complexity, expiry); Kerberos authentication; audit logs on server. \\
\hline
\textbf{3. Scalability} & Practical limit $\approx 10$ nodes. Performance degrades sharply beyond this. & Designed for hundreds/thousands. Adding 65 nodes is trivial. \\
\hline
\textbf{4. Resource Sharing \& Backup} & Files/printers scattered on workstations. Backup requires agent on each PC or manual copy. & Centralised file server (RAID), print server. Single backup job (tape/cloud) covers all user data. \\
\hline
\textbf{5. Cost (Capital \& Operational)} & Low CAPEX (no server HW, Workstation OS only). High OPEX at scale (technician time per PC). & Higher CAPEX (Server HW, Windows Server licences, CALs). Lower OPEX at scale (1 admin manages all). \\
\hline
\textbf{6. Reliability / Fault Tolerance} & If PC hosting printer/files fails, resource unavailable. No RAID typically. & Server redundancy (RAID, redundant PSU, UPS). Failover clustering possible. 99.9\% uptime target. \\
\hline
\textbf{7. Expertise Required} & Basic networking skills sufficient. & Requires trained sysadmin (MCSE/MCSA, Linux admin) for AD, DNS, DHCP, GPO. \\
\hline
\end{tabularx}
\end{center}

\noindent\textbf{(ii) Recommendation: \textbf{Client/Server Architecture}.}

\noindent\textbf{Justification 1 (Scale \& Management):} With \textbf{65 nodes} spread across three buildings (Lab, Admin, Library), decentralised management is impossible. A Domain Controller allows the single ICT teacher/admin to enforce policies (disable USB, restrict web, map drives) via Group Policy Objects (GPOs) from one console.

\noindent\textbf{Justification 2 (Security \& Data Integrity):} Admin block handles sensitive student records (grades, fees via \textit{MTN Mobile Money} or \textit{Orange Money} integration) and exam data. Centralised authentication (Active Directory) with audit logs and server-side RAID backup is mandatory for data protection compliance (Law No. 2010/012 on Cybersecurity in Cameroon). P2P offers none of this.
\end{exoresolu}

\begin{methode}[Designing a Client/Server Network Topology Diagram (Paper 3 Practical Skill)]
\textbf{Objective:} Draw a correct, labelled logical topology diagram for a Client/Server LAN, showing segmentation, server roles, and connectivity. This is a frequent Paper 3 (Practical) or Paper 2 (Diagram) task.

\begin{enumerate}
    \item \textbf{Identify the Core Backbone:} Draw a central \textbf{Core Switch} (Layer 2/3). Label it ``Core Switch (L3)''.
    \item \textbf{Place the Server Farm:} Draw a block labelled ``Server Room / Rack'' connected to Core Switch with \textbf{redundant links} (2 cables) or a single high-speed link (10 Gbps). Inside, place icons for:
    \begin{itemize}
        \item Domain Controller (AD DS, DNS, DHCP)
        \item File Server (Data, RAID)
        \item Print Server (or integrated)
        \item Application/Database Server (e.g., School Management System)
        \item Proxy / Firewall Appliance (e.g., FortiGate, pfSense) $\rightarrow$ connects to Internet Router (CAMTEL/MTN Business).
    \end{itemize}
    \item \textbf{Add Access Layer per Segment:} For each physical area (Computer Lab, Admin Block, Library), draw an \textbf{Access Switch} (24/48 port PoE+). Connect each Access Switch to Core Switch via \textbf{Uplink} (Fiber SFP+ 10G or aggregated 1G).
    \item \textbf{Connect End Devices:} Show 2--3 sample PCs/Printers per Access Switch. Label ``Cat 6a UTP''.
    \item \textbf{Wireless:} Add a \textbf{Wireless LAN Controller (WLC)} connected to Core Switch, with \textbf{Access Points (APs)} hanging off Access Switches (PoE). Label SSID: ``GBHS-Staff'', ``GBHS-Students'' (VLANs).
    \item \textbf{Label VLANs / IP Subnets:} Annotate each link/segment with VLAN ID and Subnet (e.g., VLAN 10: 192.168.10.0/24 -- Admin; VLAN 20: 192.168.20.0/24 -- Lab).
    \item \textbf{Legend:} Include a small legend for icons (Server, Switch, Router, AP, Firewall, PC).
\end{enumerate}
\textbf{Mark Scheme Checklist:} Core Switch central? Redundancy shown? Server roles distinct? VLANs/IPs labelled? WLAN controller? Internet edge? Legend?
\end{methode}

\begin{exoresolu}[Logical Topology Diagram for GBHS Mfoundi]
\textbf{Statement:}
Using the scenario from the previous example (65 nodes, 3 blocks, Internet via CAMTEL fibre), draw a labelled logical network diagram showing the Client/Server architecture. Include: Core Switch, Server Farm (list 4 servers), Access Switches per block, Wireless Controller + APs, Firewall, Router, VLANs/Subnets. Use standard icons.

\tcblower
\textbf{Detailed Solution:}

\noindent\textbf{Diagram Description (Textual representation for marking; in exam you draw neatly with ruler/pencil or Visio/draw.io in practical):}

\begin{popfigure}
\centering
\begin{tikzpicture}[
    node distance=1.8cm and 2.5cm,
    server/.style={draw, rounded corners, fill=popblueL, minimum width=2.2cm, minimum height=1cm, align=center, font=\small},
    switch/.style={draw, rectangle, fill=popgreenL, minimum width=1.8cm, minimum height=1cm, align=center, font=\small},
    router/.style={draw, circle, fill=poporange, minimum width=1.2cm, minimum height=1.2cm, align=center, font=\small},
    fw/.style={draw, diamond, fill=poppink, minimum width=1.5cm, minimum height=1.5cm, align=center, font=\small, aspect=2},
    ap/.style={draw, ellipse, fill=popgold, minimum width=1.5cm, minimum height=0.8cm, align=center, font=\small},
    pc/.style={draw, rectangle, fill=white, minimum width=1.2cm, minimum height=0.7cm, align=center, font=\tiny},
    link/.style={thick, popdark},
    uplink/.style={thick, popblue, dashed},
    wlink/.style={thick, poppurple, dotted}
]
% Internet Cloud
\node[draw, rounded corners, fill=popmuted, minimum width=2.5cm, minimum height=1.5cm, align=center] (internet) at (0,5) {CAMTEL\\Fibre};

% Edge
\node[router, below=1.2cm of internet, align=center] (router){Edge Router\\CAMTEL CPE};
\node[fw, below=1.2cm of router, align=center] (firewall){Firewall\\FortiGate / pfSense};

% Core
\node[switch, below=1.5cm of firewall, fill=popblue!20, align=center] (core){Core Switch L3\\10G Uplinks};

% Server Farm (Left of Core)
\node[server, left=2.5cm of core, yshift=1.5cm, align=center] (dc){Domain Controller\\AD DS, DNS, DHCP\\192.168.1.10};
\node[server, left=2.5cm of core, yshift=0cm, align=center] (fs){File Server\\Data + RAID 5\\192.168.1.11};
\node[server, left=2.5cm of core, yshift=-1.5cm, align=center] (apps){App / DB Server\\School Mgt System\\192.168.1.12};
\node[server, left=2.5cm of core, yshift=-3cm, align=center] (ps){Print Server\\+ Proxy Cache\\192.168.1.13};

% WLC
\node[server, right=2.5cm of core, yshift=1.5cm, align=center] (wlc){WLAN Controller\\Cisco / Ubiquiti\\192.168.1.20};

% Access Switches (Right of Core / Below WLC)
\node[switch, right=2.5cm of core, yshift=-1.5cm, align=center] (sw_admin){Access Switch\\Admin Block\\PoE+ 48pt};
\node[switch, right=2.5cm of core, yshift=-3.5cm, align=center] (sw_lab){Access Switch\\Computer Lab\\PoE+ 48pt};
\node[switch, right=2.5cm of core, yshift=-5.5cm, align=center] (sw_lib){Access Switch\\Library\\PoE+ 24pt};

% End Devices (Samples)
\node[pc, below=0.8cm of sw_admin, align=center] (pc_a1){Admin PC\\VLAN 10};
\node[pc, right=1.5cm of pc_a1, align=center] (pc_a2){Admin PC\\VLAN 10};
\node[pc, below=0.8cm of sw_lab, align=center] (pc_l1){Lab PC\\VLAN 20};
\node[pc, right=1.5cm of pc_l1, align=center] (pc_l2){Lab PC\\VLAN 20};
\node[pc, below=0.8cm of sw_lib, align=center] (pc_b1){Lib PC\\VLAN 30};

% APs
\node[ap, above=0.8cm of sw_admin, align=center] (ap_a){AP Admin\\VLAN 10/40};
\node[ap, above=0.8cm of sw_lab, align=center] (ap_l){AP Lab\\VLAN 20/40};
\node[ap, above=0.8cm of sw_lib, align=center] (ap_b){AP Lib\\VLAN 30/40};

% Connections
\draw[link] (internet) -- (router);
\draw[link] (router) -- (firewall);
\draw[link, very thick] (firewall) -- (core);

% Server Farm Links
\foreach \s in {dc,fs,apps,ps} \draw[uplink] (\s) -- (core);

% WLC Link
\draw[uplink] (wlc) -- (core);

% Access Switch Uplinks
\draw[uplink] (core) -- (sw_admin);
\draw[uplink] (core) -- (sw_lab);
\draw[uplink] (core) -- (sw_lib);

% Access to End Devices
\draw[link] (sw_admin) -- (pc_a1);
\draw[link] (sw_admin) -- (pc_a2);
\draw[link] (sw_lab) -- (pc_l1);
\draw[link] (sw_lab) -- (pc_l2);
\draw[link] (sw_lib) -- (pc_b1);

% Access to APs (PoE)
\draw[wlink] (sw_admin) -- (ap_a);
\draw[wlink] (sw_lab) -- (ap_l);
\draw[wlink] (sw_lib) -- (ap_b);

% WLC to APs (Management Tunnel - logical)
\draw[poppurple, thick, decorate, decoration={snake, amplitude=1pt, segment length=4pt}] (wlc) -- (ap_a);
\draw[poppurple, thick, decorate, decoration={snake, amplitude=1pt, segment length=4pt}] (wlc) -- (ap_l);
\draw[poppurple, thick, decorate, decoration={snake, amplitude=1pt, segment length=4pt}] (wlc) -- (ap_b);

% VLAN Labels near uplinks
\node[font=\tiny, popblue, anchor=west] at ($(core.east)+(0.2,1.0)$) {VLAN 1 (Mgmt): 192.168.1.0/24};
\node[font=\tiny, popblue, anchor=west] at ($(core.east)+(0.2,0.3)$) {VLAN 10 (Admin): 192.168.10.0/24};
\node[font=\tiny, popblue, anchor=west] at ($(core.east)+(0.2,-0.4)$) {VLAN 20 (Lab): 192.168.20.0/24};
\node[font=\tiny, popblue, anchor=west] at ($(core.east)+(0.2,-1.1)$) {VLAN 30 (Lib): 192.168.30.0/24};
\node[font=\tiny, popblue, anchor=west] at ($(core.east)+(0.2,-1.8)$) {VLAN 40 (WiFi): 192.168.40.0/24};

\end{tikzpicture}
\legende{Logical Topology: Client/Server Architecture for GBHS Mfoundi. Dashed=10G Uplinks, Dotted=PoE AP links, Snake=WLC CAPWAP Tunnel.}
\end{popfigure}

\noindent\textbf{Key Marks Checklist (Self-Assessment):}
\begin{itemize}
    \item [\checkmark] Hierarchical 3-layer (Core, Distribution/Server, Access) or Collapsed Core.
    \item [\checkmark] Distinct Server Farm with named roles (DC, File, App, Print).
    \item [\checkmark] Security Edge (Router + Firewall) before Core.
    \item [\checkmark] Wireless Controller architecture (not standalone APs).
    \item [\checkmark] VLAN segmentation per functional area (Admin, Lab, Library, WiFi, Mgmt).
    \item [\checkmark] IP Subnet notation (CIDR /24).
    \item [\checkmark] Redundancy hint (dual links to servers/core).
    \item [\checkmark] Legend / Colour coding.
\end{itemize}
\end{exoresolu}

\begin{methode}[Selecting Server Hardware Specifications for a Given Load]
\textbf{Objective:} Justify server hardware choices (CPU, RAM, Storage, NIC) based on user count, services, and performance requirements. Typical Paper 2 question: ``Specify the minimum hardware for a File/Print Server for 50 users.''

\begin{enumerate}
    \item \textbf{Define the Role(s):} Single role (File only) vs Multi-role (DC + File + Print + DHCP). Multi-role needs more resources.
    \item \textbf{Estimate Concurrent Users:} Not total staff, but \emph{simultaneous} connections. Rule of thumb: 20--30\% of total staff for file/print; 100\% for DC authentication storms (morning logon).
    \item \textbf{CPU (Processor):}
    \begin{itemize}
        \item File/Print/DC (Small $<50$): 1$\times$ Xeon Silver / AMD EPYC (8--12 cores, 2.0--2.5 GHz).
        \item Database/App Server: 2$\times$ Xeon Gold (16--24 cores total), high clock speed ($>3.0$ GHz), large L3 cache.
        \item Virtualisation Host (Hyper-V/ESXi): Max cores/RAM budget allows. Core count $>$ Clock speed for VM density.
    \end{itemize}
    \item \textbf{RAM (Memory):} \textbf{Critical for File Server (Cache) \& Virtualisation.}
    \begin{itemize}
        \item Base OS: 4--8 GB.
        \item File Server Cache: Add 1--2 GB per TB of active data (working set).
        \item Per VM (if virtualised): 8--16 GB minimum.
        \item \textbf{Rule:} Fill all memory channels. Use ECC RDIMM. Target $< 80\%$ utilisation.
    \end{itemize}
    \item \textbf{Storage (The Bottleneck):}
    \begin{itemize}
        \item \textbf{OS Boot:} 2$\times$ 240/480 GB SSD (RAID 1 / Mirror).
        \item \textbf{Data:} \textbf{RAID is mandatory.}
        \begin{itemize}
            \item RAID 1 (Mirror): 2 disks. Safety, no capacity gain. Good for OS.
            \item RAID 5 (Stripe + Parity): Min 3 disks. Good read, slow write rebuild. \textbf{Avoid for heavy DB writes.}
            \item \textbf{RAID 6 (Dual Parity):} Min 4 disks. Standard for File Server capacity + safety (survives 2 disk failures).
            \item \textbf{RAID 10 (Mirrored Stripes):} Min 4 disks (even). \textbf{Best Performance + Safety.} Use for Database, Hyper-V host, high-write File Server. 50\% capacity overhead.
        \end{itemize}
        \item Disk Type: NL-SAS / Enterprise SATA (Capacity) vs SAS / NVMe SSD (Performance). \textbf{NVMe U.2/U.3 for DB/Logs.}
    \end{itemize}
    \item \textbf{Network (NIC):} Minimum 2$\times$ 1 Gbps (Teaming/LACP). Standard now: 2$\times$ 10 Gbps SFP+ (File/DB) or 25 Gbps (Virtualisation/Storage). Dedicated iDRAC/iLO management port (OOB).
    \item \textbf{Power \& Cooling:} Dual Hot-Plug PSU (Redundant 1+1). Rackmount (1U/2U) for server room. Tower only for micro-business $<10$ users.
\end{enumerate}
\textbf{Key Formula:} \textbf{Usable Space (RAID 6)} = $(N - 2) \times \text{Disk Size}$. \textbf{Usable Space (RAID 10)} = $(N / 2) \times \text{Disk Size}$.
\end{methode}

\begin{exoresolu}[Server Specification for a Growing SME in Douala]
\textbf{Statement:}
\textit{KmerTech Ltd} (Douala) is migrating from P2P to Client/Server. They have 40 staff now, projecting 80 in 3 years. They need a single physical server running Hyper-V to host 3 VMs:
\begin{itemize}
    \item VM1: Domain Controller + DNS + DHCP + Print Server.
    \item VM2: File Server (Current data 2 TB, growing 500 GB/year). Heavy Excel/CAD access.
    \item VM3: SQL Server for ERP (Moderate OLTP, 20 concurrent users).
\end{itemize}
Budget is tight but reliability is paramount. Specify the \textbf{minimum} hardware configuration (CPU, RAM, Storage Disks + RAID, NIC, Form Factor, PSU) with justifications.

\tcblower
\textbf{Detailed Solution:}

\noindent\textbf{1. CPU: 1 $\times$ Intel Xeon Silver 4310 (12 Cores, 2.1 GHz, 18 MB Cache) or AMD EPYC 7313P (16 Cores, 2.7 GHz).}
\textbf{Justification:} 3 VMs + Hyper-V Host overhead. 12--16 cores allow dedicating 4 vCPU to SQL VM, 4 vCPU to File VM, 2 vCPU to DC VM, leaving 2--6 cores for Host/peak. Single socket saves licence cost (Windows Server Standard covers 16 cores; 12/16 cores fit perfectly). Clock speed adequate for ERP single-thread performance.

\noindent\textbf{2. RAM: 128 GB (8 $\times$ 16 GB DDR4-3200 ECC RDIMM).}
\textbf{Justification:}
\begin{itemize}
    \item Host OS (Hyper-V): 8 GB.
    \item VM1 (DC/DNS/DHCP/Print): 8 GB (light).
    \item VM2 (File Server): 32 GB (Large cache for 2--3.5 TB active data; SMB performance scales with RAM).
    \item VM3 (SQL Server): 64 GB (SQL Server \emph{loves} RAM for Buffer Pool; 20 users OLTP needs 32--64 GB for low latency).
    \item Total Allocated: 112 GB. Leaves 16 GB buffer for Host management / dynamic memory ballooning. 8 DIMMs populate all 8 channels (max bandwidth).
\end{itemize}

\noindent\textbf{3. Storage:}
\begin{itemize}
    \item \textbf{BOOT (RAID 1):} 2 $\times$ 480 GB Enterprise Read-Intensive SSD (SATA/SAS). Mirror for OS/Hyper-V partition. Fast boot, redundancy.
    \item \textbf{DATA (RAID 10):} 6 $\times$ 1.92 TB Enterprise Mixed-Use NVMe U.2 SSDs (or 8 $\times$ 1.2 TB 10K SAS HDDs if budget forces HDD).
    \textbf{Why RAID 10?} SQL Server (VM3) demands high IOPS / low latency for transaction logs and tempdb. File Server (VM2) benefits from high throughput. RAID 10 provides best write performance and fastest rebuild (survives 1 disk failure per mirror pair, sometimes 2). Usable Capacity = $6/2 \times 1.92 = 5.76$ TB. Covers current 2 TB + 3 years growth (3.5 TB) with $> 60\%$ headroom.
    \item \textit{Alternative if HDD:} 8 $\times$ 2.4 TB 10K SAS in RAID 10 = 9.6 TB usable. Slower IOPS, needs more spindles.
\end{itemize}

\noindent\textbf{4. Network: 2 $\times$ 10 Gbps SFP+ (Intel X710/XL710 or Mellanox ConnectX) + 1 $\times$ 1 Gbps Dedicated Management (iDRAC/iLO).}
\textbf{Justification:} 1 Gbps is a bottleneck for File Server (SMB3 multichannel) and SQL replication/backup. 10 Gbps uplinks to Core Switch (LACP team) future-proofs for 80 users + backup traffic. OOB management essential for remote admin (Douala to Yaoundé HQ).

\noindent\textbf{5. Form Factor: 2U Rackmount (e.g., Dell PowerEdge R750 / HPE ProLiant DL360 Gen10+).}
\textbf{Justification:} 2U accommodates 6--8+ drive bays (NVMe/U.2), better airflow for high CPU/RAM density, dual PSU standard. Tower server lacks drive bays/cooling for this spec.

\noindent\textbf{6. Power: 2 $\times$ 800W Hot-Plug Redundant PSU (Platinum/Titanium efficiency).}
\textbf{Justification:} N+1 redundancy. If one PSU fails or CAMTEL power fluctuates (common in Douala), server stays up. Supports full load + 30\% headroom.

\noindent\textbf{7. Warranty/Support: 3-Year 4-Hour On-Site (ProSupport / Foundation Care).}
\textbf{Justification:} Business continuity. ``Best effort'' next-business-day is unacceptable for ERP/File Server.
\end{exoresolu}

\begin{methode}[Planning IP Addressing \& VLAN Scheme (Subnetting for Design)]
\textbf{Objective:} Design a scalable IPv4 addressing plan using VLSM (Variable Length Subnet Mask) for a Client/Server network with multiple segments (VLANs). This tests both Network Architecture and IP Addressing (ICT02/ICT03 overlap).

\begin{enumerate}
    \item \textbf{Get the Base Block:} Usually a Private Class A ($10.0.0.0/8$) or Class B ($172.16.0.0/12$) allocated by architect. Assume $10.0.0.0/8$ or $172.16.0.0/16$ for LAN.
    \item \textbf{List Segments (VLANs) with Host Counts:} Include: Servers, Management, Admin, Lab, Library, WiFi-Staff, WiFi-Students, VoIP Phones, CCTV, IoT, DMZ, Guest. Estimate max hosts $\times 2$ for growth.
    \item \textbf{Sort Descending by Host Count:} Largest subnet first (efficient VLSM).
    \item \textbf{Allocate Subnets (VLSM Algorithm):}
    \begin{enumerate}
        \item Start with Base Network Address.
        \item For each segment, find smallest $/n$ where $2^{(32-n)} - 2 \ge \text{Required Hosts}$.
        \item Assign Network ID, Broadcast, Usable Range, Gateway (usually .1 or .254).
        \item Next segment starts at \textbf{Next IP after Broadcast} of previous.
    \end{enumerate}
    \item \textbf{Document in a Table:} Columns: VLAN ID, Name, Subnet (CIDR), Mask, Network ID, Gateway, Usable Range, Broadcast, Description.
    \item \textbf{Reserve:} First 10--20 IPs per subnet for Static (Servers, Printers, APs, Switches). DHCP Scope = Remainder.
    \item \textbf{Inter-VLAN Routing:} Layer 3 Switch (Core) has SVIs (Switch Virtual Interfaces) = Gateways.
\end{enumerate}
\textbf{Trap:} Forgetting the $-2$ (Network + Broadcast). Overlapping subnets. Using /24 for everything (wastes space, limits VLANs). Not leaving growth room.
\end{methode}

\begin{exoresolu}[VLSM Addressing Plan for GBHS Mfoundi Extension]
\textbf{Statement:}
The network architect decides to use the private block \textbf{172.16.0.0/16} for the entire school LAN. Design a VLSM addressing scheme for the following VLANs. Show the calculation for the two largest subnets.
\begin{center}
\begin{tabular}{|l|c|}
\hline
\textbf{VLAN / Segment} & \textbf{Max Hosts Required (with growth)} \\
\hline
VLAN 10 -- Admin Block & 60 \\
VLAN 20 -- Computer Lab & 120 \\
VLAN 30 -- Library & 30 \\
VLAN 40 -- WiFi Staff & 80 \\
VLAN 50 -- WiFi Students & 200 \\
VLAN 99 -- Server Farm (Mgmt) & 20 \\
VLAN 100 -- Network Devices (Switches/APs) & 30 \\
VLAN 200 -- VoIP Phones & 100 \\
VLAN 300 -- CCTV / IoT & 50 \\
\hline
\end{tabular}
\end{center}

\tcblower
\textbf{Detailed Solution:}

\noindent\textbf{Step 1: Sort Descending by Host Count.}
\begin{enumerate}
    \item VLAN 50 (WiFi Students): 200 hosts
    \item VLAN 20 (Lab): 120 hosts
    \item VLAN 200 (VoIP): 100 hosts
    \item VLAN 40 (WiFi Staff): 80 hosts
    \item VLAN 10 (Admin): 60 hosts
    \item VLAN 300 (CCTV): 50 hosts
    \item VLAN 30 (Library): 30 hosts
    \item VLAN 100 (Net Devices): 30 hosts
    \item VLAN 99 (Server Mgmt): 20 hosts
\end{enumerate}

\noindent\textbf{Step 2: Calculate Prefix Length ($/n$) for each.}
Formula: $2^{(32-n)} - 2 \ge H \implies 32-n \ge \log_2(H+2) \implies n \le 32 - \lceil \log_2(H+2) \rceil$.

\begin{itemize}
    \item \textbf{VLAN 50 (200 hosts):} $H+2=202$. $\log_2(202) \approx 7.66 \rightarrow 8$ host bits. $n = 32-8 = \mathbf{/24}$. Block size = 256.
    \item \textbf{VLAN 20 (120 hosts):} $H+2=122$. $\log_2(122) \approx 6.93 \rightarrow 7$ host bits. $n = 32-7 = \mathbf{/25}$. Block size = 128.
    \item \textbf{VLAN 200 (100 hosts):} $H+2=102$. $\log_2(102) \approx 6.67 \rightarrow 7$ host bits. $n = \mathbf{/25}$. Block size = 128.
    \item \textbf{VLAN 40 (80 hosts):} $H+2=82$. $\log_2(82) \approx 6.36 \rightarrow 7$ host bits. $n = \mathbf{/25}$. Block size = 128.
    \item \textbf{VLAN 10 (60 hosts):} $H+2=62$. $\log_2(62) \approx 5.95 \rightarrow 6$ host bits. $n = \mathbf{/26}$. Block size = 64.
    \item \textbf{VLAN 300 (50 hosts):} $H+2=52$. $\log_2(52) \approx 5.7 \rightarrow 6$ host bits. $n = \mathbf{/26}$. Block size = 64.
    \item \textbf{VLAN 30 (30 hosts):} $H+2=32$. $\log_2(32) = 5 \rightarrow 5$ host bits. $n = \mathbf{/27}$. Block size = 32.
    \item \textbf{VLAN 100 (30 hosts):} Same $\rightarrow \mathbf{/27}$. Block size = 32.
    \item \textbf{VLAN 99 (20 hosts):} $H+2=22$. $\log_2(22) \approx 4.46 \rightarrow 5$ host bits. $n = \mathbf{/27}$. Block size = 32.
\end{itemize}

\noindent\textbf{Step 3: Allocate sequentially from 172.16.0.0.}

\begin{center}
\begin{tabularx}{\linewidth}{|c|l|c|c|c|c|X|}
\hline
\rowcolor{popblueL}
\textbf{VLAN} & \textbf{Name} & \textbf{CIDR} & \textbf{Mask} & \textbf{Network ID} & \textbf{Gateway} & \textbf{Usable Range / Notes} \\
\hline
50 & WiFi Students & /24 & 255.255.255.0 & 172.16.0.0 & 172.16.0.1 & 172.16.0.2 -- 172.16.0.254 (Bcast .255) \\
\hline
20 & Computer Lab & /25 & 255.255.255.128 & 172.16.1.0 & 172.16.1.1 & 172.16.1.2 -- 172.16.1.126 (Bcast .127) \\
\hline
200 & VoIP Phones & /25 & 255.255.255.128 & 172.16.1.128 & 172.16.1.129 & 172.16.1.130 -- 172.16.1.254 (Bcast .255) \\
\hline
40 & WiFi Staff & /25 & 255.255.255.128 & 172.16.2.0 & 172.16.2.1 & 172.16.2.2 -- 172.16.2.126 \\
\hline
10 & Admin Block & /26 & 255.255.255.192 & 172.16.2.128 & 172.16.2.129 & 172.16.2.130 -- 172.16.2.190 (Bcast .191) \\
\hline
300 & CCTV/IoT & /26 & 255.255.255.192 & 172.16.2.192 & 172.16.2.193 & 172.16.2.194 -- 172.16.2.254 (Bcast .255) \\
\hline
30 & Library & /27 & 255.255.255.224 & 172.16.3.0 & 172.16.3.1 & 172.16.3.2 -- 172.16.3.30 (Bcast .31) \\
\hline
100 & Net Devices & /27 & 255.255.255.224 & 172.16.3.32 & 172.16.3.33 & 172.16.3.34 -- 172.16.3.62 (Bcast .63) \\
\hline
99 & Server Mgmt & /27 & 255.255.255.224 & 172.16.3.64 & 172.16.3.65 & 172.16.3.66 -- 172.16.3.94 (Bcast .95) \\
\hline
\end{tabularx}
\end{center}

\noindent\textbf{Verification:} Last used IP = 172.16.3.95. Next free = 172.16.3.96. Total used $< 172.16.255.255$. \textbf{Plenty of space for future VLANs.}

\noindent\textbf{Step 4: DHCP Scope Configuration Example (VLAN 20 Lab).}
\begin{itemize}
    \item Scope: 172.16.1.0/25
    \item Exclusions (Static Reservations): 172.16.1.1 (Gateway), 172.16.1.2--172.16.1.10 (Printers, Teacher PC, Switch Mgmt, AP Mgmt).
    \item DHCP Pool: 172.16.1.11 -- 172.16.1.126.
    \item Lease Duration: 8 hours (school day) or 24 hours.
    \item Options: 003 Router (172.16.1.1), 006 DNS (DC IP: 172.16.3.66), 015 Domain Name (gbhs.local).
\end{itemize}
\end{exoresolu}

\begin{methode}[Writing a Network Design Proposal (Paper 2 Essay Structure)]
\textbf{Objective:} Structure a high-scoring essay answer for ``Propose a network design for [Scenario]'' (10--15 marks).

\begin{enumerate}
    \item \textbf{Requirements Analysis (2 marks):} Bullet-point the \textbf{Functional} (File share, Internet, ERP, WiFi, Printing, Centralised Auth) and \textbf{Non-Functional} (Security, Scalability 3yrs, Availability 99.9\%, Budget, Manageability, Cameroon Law compliance) requirements extracted from the scenario.
    \item \textbf{Architecture Choice \& Justification (2 marks):} State ``Client/Server''. Justify with 2--3 points from Method 2 (Scale $>10$, Centralised Security, Scalability).
    \item \textbf{Logical Topology Diagram (3 marks):} Describe or sketch (Method 3). Mention: Hierarchical (Core-Distribution-Access or Collapsed Core), Server Farm, VLAN Segmentation, DMZ, Wireless Controller, Redundancy.
    \item \textbf{IP Addressing Plan (2 marks):} State Base Block, VLAN strategy, VLSM summary, DHCP/DNS roles (Method 5).
    \item \textbf{Server Roles \& Specifications (2 marks):} List VMs/Physical Roles (DC, File, App, Print, NPS/RADIUS, WDS/MDT). Reference specs (Method 4).
    \item \textbf{Security Implementation (2 marks):} Firewall (Edge), NPS/RADIUS (802.1X for WiFi/Wired), GPO (USB disable, Password Policy), WSUS/Endpoint Protection, Backup Strategy (3-2-1 Rule), ANTIC/Cybersecurity Law compliance.
    \item \textbf{Implementation Phases \& Testing (1--2 marks):} Phase 1: Core/Server Room. Phase 2: Admin Block (Pilot). Phase 3: Lab/Library. Phase 4: WiFi. Testing: Ping, DHCP lease, Domain Join, GPO apply, Internet, Backup Restore Test.
    \item \textbf{Bill of Materials (BoM) Summary / Budget Note (1 mark):} Key items: Core Switch (L3), Access Switches (PoE), Server (HW+OS+CALs), Firewall, APs, Controller, Cabling (Cat6a/Fibre), Rack/UPS. Mention ``Procurement via MINMAP/COPIL procedures''.
\end{enumerate}
\textbf{Golden Rule:} \textbf{Contextualise.} Don't just list generic tech. Say ``Windows Server 2022 Standard (16-core licence) for DC/File/Print VMs'' not ``Server Software''. Say ``FortiGate 60F for UTM'' not ``Firewall''.
\end{methode}

\begin{exoresolu}[Full Design Proposal: \textit{Ministry of Youth Affairs Regional Delegation}]
\textbf{Statement:}
The Ministry of Youth Affairs (MINJEC) wants to network its Regional Delegation in Bafoussam. Requirements:
\begin{itemize}
    \item 30 staff PCs (Admin), 10 PCs in a Youth Training Lab.
    \item Centralised user accounts, file storage (5 TB), network printing.
    \item Internet for all (MTN Business Fibre 100 Mbps).
    \item Secure WiFi for staff and guests (separate).
    \item CCTV (8 IP cameras) and VoIP phones (30 handsets).
    \item Must comply with Cameroon Cybersecurity Law (2010/012) and ANTIC guidelines.
    \item Budget: 15 Million FCFA (Hardware + Licences). 3-year horizon.
\end{itemize}
Write a structured design proposal covering: Architecture Choice, Logical Topology, IP Plan (VLANs), Server Roles, Security, and Key Hardware BoM.

\tcblower
\textbf{Detailed Solution:}

\noindent\textbf{1. Requirements Analysis}
\textbf{Functional:} Centralised Auth (AD), File Server (5 TB + growth), Print Server, Internet Gateway, WiFi (Corporate + Guest), VoIP, CCTV, Centralised AV/Updates.
\textbf{Non-Functional:} Security (Law 2010/012, ANTIC), Availability (Business hours 7:30--15:30, critical), Scalability (3 yrs, +50\% users), Manageability (1 SysAdmin), Budget (15M FCFA $\approx$ 23k EUR).

\noindent\textbf{2. Architecture Choice: Client/Server.}
\textbf{Justification:} 40+ wired nodes + WiFi + VoIP + CCTV $> 10$ node P2P limit. Legal requirement for centralised audit logs/user management (Law 2010/012 Art 12--15) mandates Domain Controller. Centralised backup of 5 TB+ impossible in P2P.

\noindent\textbf{3. Logical Topology (Collapsed Core)}
\begin{itemize}
    \item \textbf{Edge:} MTN ONT $\rightarrow$ FortiGate 60F (UTM: Firewall, IPS, AV, Web Filter, SSL Inspect, SD-WAN failover to 4G).
    \item \textbf{Core:} 1 $\times$ HPE Aruba 6300M 48p PoE+ (JL659A) -- Layer 3 (Routing, SVIs, DHCP Relay). Stackable for future redundancy.
    \item \textbf{Server Room (Rack 42U):} 1 $\times$ Dell PowerEdge R350 (1U) running Hyper-V.
    \item \textbf{Access:} Core switch \emph{is} the access switch (Collapsed Core) -- 48 ports cover Admin (30) + Lab (10) + VoIP (30 via PC pass-through) + APs (4) + CCTV (8) + Spares. \textit{Note: If distance $>100$m to Lab, add 1 $\times$ 24p PoE+ Access Switch in Lab linked via Fibre (SFP+ 10G). Assume single building here.}
    \item \textbf{Wireless:} 4 $\times$ Aruba AP-535 (WiFi 6) powered by Core PoE+. Managed by \textbf{Aruba Central (Cloud)} -- No on-prem controller hardware cost. SSIDs: \texttt{MINJEC-Staff} (802.1X $\rightarrow$ NPS), \texttt{MINJEC-Guest} (Captive Portal, VLAN 99, Rate Limited 5 Mbps).
\end{itemize}

\noindent\textbf{4. IP Addressing Plan (Base: 10.10.0.0/16)}
\begin{center}
\begin{tabular}{|c|l|c|c|}
\hline
\rowcolor{popblueL}
\textbf{VLAN} & \textbf{Name} & \textbf{Subnet} & \textbf{Gateway / Notes} \\
\hline
10 & Mgmt (Switches, Server, iDRAC) & 10.10.0.0/24 & 10.10.0.1 (Core SVI) \\
20 & Admin Staff (Wired) & 10.10.10.0/24 & 10.10.10.1 (DHCP) \\
30 & Training Lab (Wired) & 10.10.20.0/24 & 10.10.20.1 (DHCP) \\
40 & VoIP Phones & 10.10.30.0/24 & 10.10.30.1 (DHCP Opt 66/150) \\
50 & WiFi Staff (802.1X) & 10.10.40.0/23 & 10.10.40.1 (510 hosts) \\
60 & WiFi Guest (Captive) & 10.10.50.0/24 & 10.10.50.1 (Isolated ACL) \\
70 & CCTV / IoT & 10.10.60.0/24 & 10.10.60.1 (No Internet ACL) \\
80 & Server Farm (Data) & 10.10.100.0/24 & 10.10.100.1 (Static IPs) \\
\hline
\end{tabular}
\end{center}
\textit{Inter-VLAN Routing on Core Switch SVIs. ACLs: Guest $\nrightarrow$ Internal; CCTV $\nrightarrow$ Internet; Staff $\rightarrow$ Internet (Filtered).}

\noindent\textbf{5. Server Roles (Virtualised on Single Host)}
\begin{itemize}
    \item \textbf{VM1 (DC01):} AD DS, DNS, DHCP (Scopes for VLANs 20,30,40,50,60,70), NPS (RADIUS for 802.1X), WDS/MDT (Imaging). \textit{Spec: 4 vCPU, 16 GB RAM, 120 GB OS Disk.}
    \item \textbf{VM2 (FILE01):} File Server (DFS Namespaces, FSRM Quotas/Screens), Print Server. \textit{Spec: 4 vCPU, 32 GB RAM, 8 TB RAID 10 Data Disk (4 $\times$ 2.4 TB SAS HDD or 4 $\times$ 1.92 TB SSD).}
    \item \textbf{VM3 (UTIL01):} WSUS (Updates), SCCM/MECM Client (or PDQ Deploy), Veeam Backup \& Replication Server. \textit{Spec: 4 vCPU, 16 GB RAM, 2 TB Disk.}
\end{itemize}
\textbf{Host Spec (R350):} 1$\times$ Xeon Silver 4310 (12C), 128 GB RAM, 2$\times$ 480 GB SSD RAID 1 (OS), 8$\times$ 2.5'' Bays (Data), 2$\times$ 10Gb SFP+ NIC, 2$\times$ 550W PSU, 3Yr ProSupport. \textit{Est. 6.5M FCFA.}

\noindent\textbf{6. Security Implementation (ANTIC Compliance)}
\begin{itemize}
    \item \textbf{Perimeter:} FortiGate 60F: IPS, Application Control, Antivirus, Web Filter (Block Proxy/VPN, Adult, Malware), SSL Inspection (Deep Inspection), Geo-IP Block (High risk countries), Log to FortiAnalyzer Cloud / Syslog Server (VM3).
    \item \textbf{Endpoint:} GPO: Password Policy (12 chars, complexity, 90 days), Account Lockout (5/30/30), Restrict Local Admin, Disable USB Storage, Firewall ON, BitLocker on Laptops.
    \item \textbf{Network Access:} 802.1X (Wired + WiFi Staff) via NPS $\rightarrow$ Machine Auth (Domain Computer) + User Auth. MAB (MAC Auth Bypass) for VoIP Phones / Printers.
    \item \textbf{Backup (3-2-1):} Veeam on VM3 $\rightarrow$ Local NAS (Synology 4-bay, RAID 5, 16 TB) $\rightarrow$ Offsite: Encrypted Copy to Azure Blob / Wasabi / Local USB Rotated Weekly (taken home by Admin).
    \item \textbf{Audit:} AD Audit Policy (Success/Failure Logon, Object Access). Logs retained 1 year (Law 2010/012).
\end{itemize}

\noindent\textbf{7. Key Bill of Materials (Est. FCFA)}
\begin{center}
\begin{tabularx}{\linewidth}{|l|X|r|}
\hline
\rowcolor{popblueL}
\textbf{Item} & \textbf{Description} & \textbf{Est. Cost (FCFA)} \\
\hline
Server & Dell R350 (Xeon Silver, 128GB, 2x480SSD, 8-bay, 2x10G, 2xPSU, 3Yr) & 6,500,000 \\
Core Switch & HPE Aruba 6300M 48p PoE+ (JL659A) + 2x SFP+ DAC & 1,800,000 \\
Firewall & FortiGate 60F + 3Yr UTP (UTM) Licence & 1,500,000 \\
Access Points & 4x Aruba AP-535 (WiFi 6) + 3Yr Central Licence & 1,200,000 \\
NAS Backup & Synology DS923+ + 4x 4TB IronWolf (RAID 5 $\approx$ 12TB) & 800,000 \\
Cabling & Cat6a UTP (305m $\times$ 4 boxes), Fibre Patch (OM4), Patch Panels, Rack 22U, UPS 3kVA & 1,500,000 \\
Licences & Windows Server 2022 Std (16 Core) $\times$ 1, CALs User $\times$ 40, Veeam Std $\times$ 2 sockets & 1,200,000 \\
VoIP & Yealink T46U $\times$ 30 (PoE) & 900,000 \\
CCTV & Hikvision 8ch NVR + 8x IP Dome 4MP + 2TB HDD & 600,000 \\
Install/Config & Labour, Config, Documentation, Training (1 week) & 800,000 \\
\hline
\textbf{TOTAL} & & \textbf{16,800,000} \\
\hline
\end{tabularx}
\end{center}
\textbf{Note:} Slightly over 15M. \textbf{Optimisation:} Drop NAS to 2-bay (RAID 1, 8TB) -400k; Use Windows Server Backup (free) instead of Veeam -200k; Negotiate FortiGate licence -100k. \textbf{Revised Total $\approx$ 15.1M FCFA.} Acceptable contingency 5\%.

\noindent\textbf{8. Implementation Phases}
\begin{enumerate}
    \item \textbf{Week 1:} Server Rack Build, Hyper-V, VMs, AD, DNS, DHCP, NPS, File Shares, GPOs, Backup Jobs.
    \item \textbf{Week 2:} Core Switch Config (VLANs, SVIs, DHCP Relay, STP Root, Port Security, 802.1X). FortiGate Config (WAN, LAN, Policies, UTM, VPN Admin).
    \item \textbf{Week 3:} Cabling Termination (Admin + Lab), Patch Panels, Labelling (TIA-606). Connect PCs/Phones. Test Domain Join, GPO, Print, Internet.
    \item \textbf{Week 4:} AP Mounting, Aruba Central Onboarding, SSID Config, 802.1X Test, Guest Portal Branding. CCTV/VoIP VLANs + ACLs.
    \item \textbf{Week 5:} Documentation (As-Built Diagrams, IP Plan, Credentials Vault), Handover, Admin Training, Final Audit (ANTIC Checklist).
\end{enumerate}
\end{exoresolu}

\begin{methode}[Troubleshooting Common Architecture Misconfigurations (Practical/Oral Defence)]
\textbf{Objective:} Diagnose and fix typical design/implementation errors in a Client/Server deployment.

\begin{enumerate}
    \item \textbf{Symptom: ``Users can't log in / ``No logon servers available''.}
    \begin{itemize}
        \item \textbf{Check:} DNS on Client $\rightarrow$ Must point \textbf{ONLY} to Domain Controller(s) (DC IP). Not ISP DNS, Not Router DNS, Not 8.8.8.8.
        \item \textbf{Check:} SRV Records in DNS (\_ldap.\_tcp.dc.\_msdcs.domain.local). Run \texttt{ipconfig /flushdns}, \texttt{nltest /dsgetdc:domain.local}.
        \item \textbf{Check:} Time Sync (Max 5 min drift). \texttt{w32tm /query /status}. PDC Emulator syncs to external (time.windows.com / ntp.ubuntu.com); Clients sync to DC.
    \end{itemize}
    \item \textbf{Symptom: ``WiFi connects but no Internet / Can't reach Server''.}
    \begin{itemize}
        \item \textbf{Check:} VLAN assignment on AP / Switch Port. Is the port Trunk (tagged) or Access (untagged)?
        \item \textbf{Check:} DHCP Relay (IP Helper) on Core Switch SVI for that VLAN pointing to DHCP Server IP.
        \item \textbf{Check:} Firewall ACL on Core SVI or Edge Firewall blocking Inter-VLAN or Internet.
        \item \textbf{Check:} Captive Portal / 802.1X Radius timeout (NPS down? Shared Secret mismatch?).
    \end{itemize}
    \item \textbf{Symptom: ``File Server Slow / Disconnects''.}
    \begin{itemize}
        \item \textbf{Check:} NIC Teaming / LACP config mismatch (Switch: Active $\leftrightarrow$ Server: Active). Disable LACP if switch doesn't support; use Switch Independent Teaming.
        \item \textbf{Check:} SMB Signing / Encryption overhead. Disable signing if not required (GPO) for speed, but security trade-off.
        \item \textbf{Check:} Disk Queue Length (PerfMon: \texttt{PhysicalDisk\% Idle Time}, \texttt{Avg. Disk sec/Read}). RAID rebuild running? AV scanning live data? Exclude data volume from Real-time AV.
        \item \textbf{Check:} Opportunistic Locking (OpLocks) / SMB2/3 credits. Update NIC drivers / Firmware.
    \end{itemize}
    \item \textbf{Symptom: ``Print Jobs Stuck / Disappear''.}
\begin{itemize}
        \item \textbf{Check:} Print Spooler Service on Print Server (Restart). Clear \texttt{C:\textbackslash{}Windows\textbackslash{}System32\textbackslash{}spool\textbackslash{}PRINTERS}.
        \item \textbf{Check:} Driver Isolation (Print Management $\rightarrow$ Properties $\rightarrow$ Advanced $\rightarrow$ ``Print directly to printer'' vs Spool). Use \textbf{Type 4 (v4) Drivers} / Universal Print Driver.
        \item \textbf{Check:} Disk Space on Print Server (Spool folder).
    \end{itemize}
    \item \textbf{Symptom: ``Internet Slow / Intermittent''.}
    \begin{itemize}
        \item \textbf{Check:} Bandwidth Saturation (FortiGate / NetFlow / SNMP). Top talkers: Windows Updates (WSUS needed!), Video Streaming, Backups running daytime.
        \item \textbf{Check:} DNS Resolution Latency (Forwarders on DC $\rightarrow$ ISP / 1.1.1.1 / 8.8.8.8). Root Hints only = Slow.
        \item \textbf{Check:} MTU Mismatch (VPN / PPPoE). Ping with \texttt{-f -l 1472} (Don't Fragment).
    \end{itemize}
\end{enumerate}
\textbf{Oral Defence Tip:} Always structure answer: \textbf{Observation $\rightarrow$ Hypothesis $\rightarrow$ Verification Command $\rightarrow$ Fix $\rightarrow$ Prevention (Documentation/GPO/Monitoring)}.
\end{methode}

\begin{exoresolu}[Troubleshooting a ``Ghost'' Domain Join Failure]
\textbf{Statement:}
A technician at \textit{CamTech Solutions} (Buea) builds a new Windows 11 Pro workstation. Network settings: IP 192.168.10.50/24, GW 192.168.10.1, DNS 8.8.8.8. The Domain Controller is at 192.168.10.10 (DC01.camtech.local). The technician tries to join \texttt{camtech.local} and gets: ``The specified domain either does not exist or could not be contacted.'' Ping to 192.168.10.10 works. Ping to \texttt{dc01.camtech.local} fails. Identify the error, explain \textbf{why} it breaks domain join, and give the fix.

\tcblower
\textbf{Detailed Solution:}

\noindent\textbf{1. Identification of Error:}
The workstation is configured with \textbf{Public DNS (8.8.8.8)} as its primary (or only) DNS server instead of the \textbf{Domain Controller's IP (192.168.10.10)}.

\noindent\textbf{2. Technical Explanation (Why it breaks Domain Join):}
\begin{itemize}
    \item \textbf{Active Directory relies on DNS Service Locator (SRV) Records.} To find a Domain Controller for domain \texttt{camtech.local}, the client queries DNS for records like \texttt{\_ldap.\_tcp.dc.\_msdcs.camtech.local} and \texttt{\_kerberos.\_tcp.camtech.local}.
    \item \textbf{Public DNS Servers (Google 8.8.8.8, Cloudflare 1.1.1.1, ISP DNS) are \textbf{not authoritative} for the private zone \texttt{camtech.local}.} They do not host the internal AD-integrated DNS zone.
    \item \textbf{Result:} The public DNS server responds with \texttt{NXDOMAIN} (Non-Existent Domain) or \texttt{SERVFAIL} for the SRV records and the A record for \texttt{dc01.camtech.local}.
    \item \textbf{Client Logic:} The Windows NetLogon service / DNS Client receives ``Domain Not Found'' from the configured DNS. It does \textbf{not} fall back to broadcast or NetBIOS for domain location in modern versions (disabled by default for security). Hence the error: ``The specified domain either does not exist or could not be contacted.''
    \item \textbf{Ping 192.168.10.10 works} because it uses ICMP (Layer 3/4) directly to the IP, bypassing DNS name resolution entirely. \textbf{Ping dc01.camtech.local fails} because the name resolution step fails at the public DNS.
\end{itemize}

\noindent\textbf{3. The Fix (Two equivalent methods):}
\textbf{Method A (Static IP - Immediate):}
\begin{enumerate}
    \item Open \texttt{ncpa.cpl} $\rightarrow$ Right-click Ethernet $\rightarrow$ Properties $\rightarrow$ IPv4 $\rightarrow$ Properties.
    \item Set \textbf{Preferred DNS Server: 192.168.10.10} (DC IP).
    \item Set \textbf{Alternate DNS Server: 192.168.10.11} (Secondary DC) \textbf{OR} Leave blank. \textbf{NEVER put a public DNS here as Alternate.} (Client tries Preferred, if timeout $\rightarrow$ Alternate. If Alternate is Public, it gets NXDOMAIN for internal names $\rightarrow$ Fail).
    \item \texttt{ipconfig /flushdns} $\rightarrow$ Retry Domain Join.
\end{enumerate}

\textbf{Method B (DHCP - Correct Design):}
\begin{enumerate}
    \item Ensure the DHCP Scope (on DC or Router) for VLAN 10 (192.168.10.0/24) has \textbf{Option 006 (DNS Servers)} set to \textbf{192.168.10.10} (and 192.168.10.11).
    \item Ensure \textbf{Option 015 (DNS Domain Name)} is set to \texttt{camtech.local}.
    \item Set Workstation to ``Obtain DNS automatically''.
    \item \texttt{ipconfig /release \&\& ipconfig /renew}.
\end{enumerate}

\noindent\textbf{4. Prevention (Best Practice):}
\begin{itemize}
    \item \textbf{GPO:} ``Network\textbackslash{}DNS Client\textbackslash{}Dynamic Update'' / ``Primary DNS Suffix''.
    \item \textbf{Documentation:} Network Diagram shows ``DNS $\rightarrow$ DC IP'' on all internal segments.
    \item \textbf{Monitoring:} Script to alert if DHCP Option 006 $\neq$ DC IPs.
\end{itemize}
\end{exoresolu}

\begin{methode}[Evaluating Cloud / Hybrid Architecture vs On-Premises (Modern Context)]
\textbf{Objective:} Compare traditional On-Premises Client/Server with Cloud (Azure AD / Entra ID, Microsoft 365, AWS) and Hybrid models for a Cameroonian organisation. Increasingly relevant for GCE Modern ICT.

\begin{enumerate}
    \item \textbf{Define the Models:}
    \begin{itemize}
        \item \textbf{On-Prem (Traditional C/S):} Physical/Virtual Servers in local Server Room. AD DS, File Server, Print, Apps local. Full Control. High CAPEX/OPEX (Power, Cooling, Hardware Refresh, SysAdmin).
        \item \textbf{Cloud-Native (Azure AD / Entra ID + Intune + SharePoint/OneDrive + Azure Virtual Desktop):} No Domain Controller. Identity = Cloud (Entra ID). Management = MDM (Intune). Files = SharePoint/OneDrive. Apps = SaaS / AVD. Zero Trust. OpEx (Subscriptions). Requires \textbf{Reliable, High-Bandwidth Internet}.
        \item \textbf{Hybrid (AD Connect / Entra Connect Sync):} On-Prem AD DS $\leftrightarrow$ Sync $\rightarrow$ Entra ID. Devices: Hybrid Azure AD Joined (Domain Joined + Registered in Cloud). Best of both: Local auth speed (Kerberos/NTLM) + Cloud services (M365, Conditional Access, Intune co-management).
    \end{itemize}
    \item \textbf{Decision Matrix for Cameroon Context:}
    \begin{center}
    \begin{tabularx}{\linewidth}{|l|X|X|X|}
    \hline
    \rowcolor{popblueL}
    \textbf{Factor} & \textbf{On-Prem} & \textbf{Hybrid} & \textbf{Cloud-Native} \\
    \hline
    \textbf{Internet Dependency} & Low (Works offline) & Medium (Sync/Cloud apps need Net) & \textbf{Critical} (100\% uptime needed) \\
    \hline
    \textbf{Cameroon Internet Reality} & CAMTEL/MTN Fibre OK in Douala/Yaoundé. Unreliable/Expensive upcountry. & Feasible in Major Cities. Risky in Rural. & \textbf{High Risk} outside major hubs. \\
    \hline
    \textbf{Data Sovereignty / Law 2010/012} & Full Control. Data stays in Country. & Syncs Hashes to MS Cloud (EU/US Regions). Check ANTIC guidance. & Data in MS Cloud (South Africa / EU regions). Legal review needed. \\
    \hline
    \textbf{Legacy Apps (Custom ERP, Old FoxPro/Access)} & Native Support. & Supported (App Proxy / AVD). & Difficult (Refactor / AVD / RDWeb). \\
    \hline
    \textbf{Print / Large File (CAD/Video)} & Fast (LAN 1-10 Gbps). & Hybrid: Local Print Server / File Server + Cloud Sync. & Slow (Upload/Download over WAN). \\
    \hline
    \textbf{Skill Set Available} & MCSE/MCSA (Common in CM). & Hybrid Skills (Rare, High Value). & Cloud Admin (AZ-104, MS-102) - Emerging. \\
    \hline
    \textbf{Cost Model} & High CAPEX, Predictable OPEX. & Medium CAPEX + Subscriptions. & Pure OpEx (Per User/Month). FX Risk (USD/FCFA). \\
    \hline
    \end{tabularx}
    \end{center}
    \item \textbf{Recommendation Logic for Exam:}
    \begin{itemize}
        \item \textbf{Scenario: ``Rural School / Govt Dept, Unstable Internet, Legacy App, Low Budget''} $\rightarrow$ \textbf{On-Prem Client/Server.}
        \item \textbf{Scenario: ``Modern SME, Douala/Yaoundé, Good Fibre, Mobile Workforce, M365 Licences, No Legacy App''} $\rightarrow$ \textbf{Cloud-Native (Entra ID + Intune).}
        \item \textbf{Scenario: ``Large Org, Mixed Sites, Existing AD, Moving to M365, Need Single Sign-On''} $\rightarrow$ \textbf{Hybrid (AD Connect).}
    \end{itemize}
\end{enumerate}
\textbf{Key Terminology:} Entra ID (formerly Azure AD), Intune (MDM/MAM), AD Connect (Sync), Hybrid Join, Conditional Access, Zero Trust, Data Residency.
\end{methode}

\begin{exoresolu}[Architecture Choice for a Cameroonian NGO]
\textbf{Statement:}
\textit{Green Horizon Cameroon}, an environmental NGO, has headquarters in Yaoundé (25 staff, stable MTN Fibre 50 Mbps) and 4 field offices in Bertoua, Maroua, Bamenda, Buea (5 staff each, 4G MiFi / VSAT only, frequent outages). They use a custom Access/Excel based Project Tracking App (legacy, not web-based). They want to migrate to Microsoft 365 E3 for Email/Teams/SharePoint. They have 1 part-time IT consultant. Recommend the Identity \& Management Architecture (On-Prem, Hybrid, or Cloud-Native) with 3 justifications.

\tcblower
\textbf{Detailed Solution:}

\noindent\textbf{Recommendation: \textbf{Hybrid Architecture (On-Prem AD DS at HQ + Entra Connect Sync $\rightarrow$ Entra ID)}.}

\noindent\textbf{Justification 1: Field Office Connectivity Reality (The Decider).}
The 4 field offices rely on \textbf{4G MiFi / VSAT with frequent outages}. A Cloud-Native (Entra ID only) architecture requires \textbf{live Internet connectivity for every authentication} (Token acquisition, Conditional Access evaluation, Intune policy sync). During outages, staff cannot log on to laptops, access cached OneDrive files (requires token refresh), or use Teams. With Hybrid: Laptops are \textbf{Hybrid Azure AD Joined} and \textbf{Domain Joined}. They authenticate locally via \textbf{Cached Credentials (Kerberos/NTLM)} to the local domain profile when offline. They work seamlessly. When online, they sync with Entra ID for M365 services. This is the \textbf{only} architecture guaranteeing productivity in Cameroon's unstable upcountry connectivity.

\noindent\textbf{Justification 2: Legacy Application Support.}
The custom \textbf{Access/Excel Project Tracking App} requires local file locking (SMB/OpLocks) and low-latency LAN access to a shared backend (.accdb/.xlsx on File Server). It \textbf{cannot run over high-latency, lossy WAN/4G} to a cloud file share (SharePoint/OneDrive sync latency breaks Access locking).
\textbf{Hybrid Solution:} Keep the App \& File Server \textbf{On-Prem at HQ}. Field offices connect via \textbf{Site-to-Site VPN (IPsec / SD-WAN)} to HQ when online. For offline work: Use local cached copy + Sync script (PowerShell/Robocopy) on reconnect, or migrate App to Azure Virtual Desktop (AVD) hosted in Azure South Africa (Low latency from CM) accessed via RDP over VPN/4G. Hybrid allows this phased transition.

\noindent\textbf{Justification 3: Single IT Consultant Manageability \& Cost.}
\begin{itemize}
    \item \textbf{On-Prem Only:} Requires managing Exchange Server (complex, high skill), Patch Management (WSUS), Backup, AD Health \emph{alone}. High risk for 1 part-time person.
    \item \textbf{Cloud-Native:} Requires deep Intune/Entra ID/Graph/PowerShell skills. No local fallback. High subscription cost (E3 $\approx$ \$36/user/month $\times$ 45 users $\approx$ 120M FCFA/yr at 600 FCFA/\$ -- FX risk).
    \item \textbf{Hybrid (Sweet Spot):} On-Prem: 1 DC (VM) + 1 File/Print (VM) at HQ. Low maintenance (Core AD, DHCP, DNS, File). Cloud: Entra Connect Sync (Automated), M365 (Microsoft manages Exchange/SharePoint/Teams infra), Intune (Co-management: ConfigMgr/Intune for apps/policies). Consultant focuses on Sync Health, GPO/Intune Policies, VPN, Backup. Cost controlled (M365 Business Premium $\approx$ \$22/user cheaper than E3, or F3 for field staff).
\end{itemize}

\noindent\textbf{Implementation Sketch:}
\begin{itemize}
    \item \textbf{HQ (Yaoundé):} 1 Physical Host (Hyper-V) $\rightarrow$ VM1: DC01 (AD DS, DNS, DHCP, NPS, AD Connect), VM2: FILE01 (App Data, User Profiles, Print). FortiGate 60F (VPN Hub).
    \item \textbf{Field Offices:} FortiGate 40F (VPN Spoke, SD-WAN 4G+VSAT). Laptops: Domain Joined + Hybrid Entra Joined (GPO/Intune). Cached Logon enabled. OneDrive Known Folder Move (Desktop/Docs/Pics) $\rightarrow$ Syncs to SharePoint when VPN up.
    \item \textbf{Identity:} User created in On-Prem AD $\rightarrow$ Syncs to Entra ID (Password Hash Sync + Seamless SSO) $\rightarrow$ Licences assigned in Entra ID $\rightarrow$ M365 Apps deployed via Intune.
\end{itemize}
\end{exoresolu}

\begin{aretenir}[Summary: Master Checklist for Network Design Questions]
\begin{itemize}
    \item \textbf{Architecture ID:} Resource Location $\rightarrow$ Roles $\rightarrow$ P2P vs C/S.
    \item \textbf{Comparison Table:} 6 Criteria (Admin, Security, Scale, Cost, Perf, Reliability). Contrast Keywords.
    \item \textbf{Topology Diagram:} Hierarchical (Core-Dist-Access), Server Farm, VLANs, WLC, Firewall, Redundancy, IPs, Legend.
    \item \textbf{Server Specs:} Role $\rightarrow$ vCPU/vRAM $\rightarrow$ Storage (RAID 10 for Perf/DB, RAID 6 for Capacity) $\rightarrow$ NIC 10G $\rightarrow$ Dual PSU $\rightarrow$ Warranty.
    \item \textbf{IP Plan (VLSM):} Sort Desc $\rightarrow$ Calculate /n $\rightarrow$ Allocate Sequential $\rightarrow$ Document Table $\rightarrow$ DHCP Scope/Exclusions.
    \item \textbf{Proposal Structure:} Reqs $\rightarrow$ Arch Choice $\rightarrow$ Topology $\rightarrow$ IP $\rightarrow$ Servers $\rightarrow$ Security (ANTIC/Law) $\rightarrow$ Phases $\rightarrow$ BoM.
    \item \textbf{Troubleshooting:} DNS $\rightarrow$ Time $\rightarrow$ VLAN/DHCP Relay $\rightarrow$ Firewall ACL $\rightarrow$ NIC Teaming $\rightarrow$ SMB/Print Spooler $\rightarrow$ Bandwidth/MTU.
    \item \textbf{Modern Context:} On-Prem vs Hybrid vs Cloud-Native. \textbf{Cameroon Internet Reality} is the primary filter. Hybrid = Resilience for Field Offices.
\end{itemize}
\end{aretenir}

\newpage

\coursec{Exercises}

\courssub{I check my knowledge}

\begin{exercice}[Multiple choice questions]
For each question, choose the single correct answer.
\begin{enumerate}
\item In a peer-to-peer network, each computer:
\begin{enumerate}
\item can only receive resources from the server;
\item can act both as a client and as a server;
\item must be administered centrally;
\item cannot share its own files.
\end{enumerate}
\item The computer that provides services (files, printing, user accounts) to other computers in a client/server network is called:
\begin{enumerate}
\item a workstation;
\item a hub;
\item a server;
\item a repeater.
\end{enumerate}
\item Which of the following is a typical disadvantage of a peer-to-peer network?
\begin{enumerate}
\item It requires an expensive dedicated server;
\item Security and backups are difficult to centralise;
\item It cannot work without the Internet;
\item It needs a network administrator.
\end{enumerate}
\item In a client/server architecture, if the main server breaks down:
\begin{enumerate}
\item the workstations continue to access shared resources normally;
\item the whole network services may become unavailable;
\item the workstations automatically become servers;
\item only the printer stops working.
\end{enumerate}
\end{enumerate}
\end{exercice}

\begin{exercice}[True or false]
State whether each statement is \textbf{true} or \textbf{false}, and justify your answer in one or two sentences.
\begin{enumerate}
\item In a peer-to-peer network, there is no dedicated server.
\item A client/server network is always cheaper to set up than a peer-to-peer network.
\item In a client/server network, user accounts and passwords can be managed centrally on the server.
\item A peer-to-peer network is well suited to an organisation of 200 computers.
\item Each user in a peer-to-peer network is responsible for the security of his or her own computer.
\end{enumerate}
\end{exercice}

\begin{exercice}[Definitions]
Define each of the following terms in one or two clear sentences, as you would in the GCE examination.
\begin{enumerate}
\item Network architecture.
\item Peer-to-peer network.
\item Client/server network.
\item Workstation.
\item Dedicated server.
\end{enumerate}
\end{exercice}

\begin{exercice}[Quick questions]
Answer each question briefly.
\begin{enumerate}
\item State \textbf{two} advantages and \textbf{two} disadvantages of a peer-to-peer network.
\item State \textbf{two} advantages and \textbf{two} disadvantages of a client/server network.
\item Give one example of a situation in Cameroon where a peer-to-peer network is appropriate, and one example where a client/server network is appropriate.
\item A small office has 5 computers connected in peer-to-peer. If each computer stores its own users' files, how many computers must be checked to perform a complete backup? Justify your answer.
\end{enumerate}
\end{exercice}

\courssub{I practise}

\begin{exercice}[Completing a comparison table]
Copy and complete the following table by writing the correct entry for a peer-to-peer (P2P) network and a client/server (C/S) network for each criterion.

\begin{center}
\begin{tabularx}{\linewidth}{|l|X|X|}
\hline
\rowcolor{popblueL} \textbf{Criterion} & \textbf{Peer-to-peer} & \textbf{Client/server} \\
\hline
Cost of installation & & \\
\hline
Need for a network administrator & & \\
\hline
Centralised security & & \\
\hline
Centralised backup & & \\
\hline
Scalability (adding many computers) & & \\
\hline
Effect of one machine failing & & \\
\hline
\end{tabularx}
\end{center}

Then state which architecture you would choose for a school computer laboratory of 25 computers, and justify your choice with two reasons.
\end{exercice}

\begin{exercice}[Identifying the architecture]
For each of the following situations, state whether the network described is peer-to-peer or client/server, and justify your answer in one sentence.
\begin{enumerate}
\item Three friends in a quarter of Bamenda connect their laptops to share films and music; each laptop shares its own folders.
\item A bank in Douala stores all customer data on one powerful central machine; the cashiers' computers send requests to this machine.
\item A small printing press with 4 computers shares one printer; every computer can access the shared folders of the others.
\item A ministry uses a central computer to manage user accounts, emails and backups for all staff computers.
\end{enumerate}
\end{exercice}

\begin{exercice}[Drawing a client/server network]
A microfinance institution in Buea has one file server, one shared printer and five teller workstations, all connected through a central switch.
\begin{enumerate}
\item Draw and label the network diagram, showing the server, the switch, the printer and the five workstations.
\item On your diagram, use arrows to show the flow of a request when a teller prints a customer's statement.
\item Identify the role of each component in the network.
\end{enumerate}
\end{exercice}

\begin{exercice}[Choosing an architecture]
For each scenario below, recommend an architecture (peer-to-peer or client/server) and give \textbf{two} justified reasons for your choice.
\begin{enumerate}
\item A home network with 2 computers and one printer.
\item A hospital with 60 computers where patient records must be protected and backed up every night.
\item A small business centre with 6 computers and a very limited budget.
\item A secondary school where 300 students must each have a personal login account usable on any of the 40 computers in the laboratory.
\end{enumerate}
\end{exercice}

\courssub{I prepare for the GCE}

\begin{exercice}[Structured question: advising a cybercafé owner]
Mr. Tabi owns a cybercafé in Kumba. He has 8 computers, one printer, and a small budget. He wants all the computers to share the Internet connection, the printer and some files.
\begin{enumerate}
\item Define the term \emph{network architecture}. \hfill (2 marks)
\item Name the two main network architectures studied. \hfill (2 marks)
\item Advise Mr. Tabi on the architecture to use, giving \textbf{four} justified reasons linked to his situation. \hfill (4 marks)
\item State \textbf{two} limitations of the architecture you have recommended, and explain how Mr. Tabi can reduce each of them. \hfill (2 marks)
\end{enumerate}
\hfill \textbf{(Total: 10 marks)}
\end{exercice}

\begin{exercice}[Case study: connecting a ministry]
A ministry in Yaoundé wants to connect 50 computers spread over three departments. The management requires: centralised user accounts and passwords, automatic daily backups of all data, controlled access to files according to each employee's department, and the possibility of adding more computers in the future.
\begin{enumerate}
\item Recommend the most appropriate network architecture for this ministry, justifying your choice with reference to the four requirements above. \hfill (4 marks)
\item Describe the hardware equipment needed (server, workstations, switch, cabling or wireless access points) and the role of each element. \hfill (4 marks)
\item Describe the software roles needed on the server (user account management, file services, backup, security). \hfill (3 marks)
\item Explain what would happen to the network services if the single server failed, and propose one measure to reduce this risk. \hfill (2 marks)
\end{enumerate}
\hfill \textbf{(Total: 13 marks)}
\end{exercice}

\begin{exercice}[Essay: discussing a statement]
\emph{``For a large organisation, the client/server architecture is always the best choice.''}

Discuss this statement with reference to cost, security, scalability and administration. In your essay you should:
\begin{enumerate}
\item explain how a client/server network works and why it suits large organisations; \hfill (4 marks)
\item discuss the advantages of client/server networks concerning security, scalability and centralised administration; \hfill (4 marks)
\item discuss the limitations of client/server networks, particularly concerning cost and dependence on the server; \hfill (3 marks)
\item give a reasoned conclusion, mentioning situations where a peer-to-peer network may still be preferred, with at least one Cameroonian example. \hfill (2 marks)
\end{enumerate}
\hfill \textbf{(Total: 13 marks)}
\end{exercice}

\newpage

\coursec{Solutions to the exercises}

\coursec{Corrections and Worked Examples for Network Architecture}

\begin{corrige}[Exercise 1 — Multiple choice questions]
\begin{enumerate}
\item \textbf{Correct answer: (b) can act both as a client and as a server.}\par
In a \cle{peer-to-peer} network there is no dedicated server: every computer can request resources from the others (client role) and offer its own resources such as files or a printer (server role). Answer (a) is wrong because it describes a pure client; (c) is wrong because P2P administration is decentralised; (d) is wrong because sharing one's own files is precisely the principle of P2P.
\item \textbf{Correct answer: (c) a server.}\par
The \cle{server} is the powerful computer that provides services (file storage, printing, user accounts, authentication) to the \cle{client} workstations. A workstation consumes services, while a hub or a repeater only transmits signals.
\item \textbf{Correct answer: (b) Security and backups are difficult to centralise.}\par
Since each user manages his or her own machine, there is no central control of passwords, access rights or backups. Answers (a) and (d) are false because P2P needs neither a dedicated server nor an administrator; (c) is false because a P2P network works perfectly without the Internet.
\item \textbf{Correct answer: (b) the whole network services may become unavailable.}\par
The server is the single central point of the architecture: if it fails, authentication, shared files and centralised printing stop. This is the main weakness (\cle{single point of failure}) of the \cle{client/server} model.
\end{enumerate}
\end{corrige}

\begin{corrige}[Exercise 2 — True or false]
\begin{enumerate}
\item \textbf{True.} By definition, a \cle{peer-to-peer} network has no dedicated server: all computers are equal (\emph{peers}) and each one can both offer and request resources.
\item \textbf{False.} A \cle{client/server} network requires a powerful dedicated server, a server operating system (often licensed) and usually a network administrator, so it is generally \emph{more} expensive to set up than a peer-to-peer network.
\item \textbf{True.} Centralised management of user accounts and passwords (authentication) on the server is one of the main advantages of the \cle{client/server} architecture: a user can log on from any workstation.
\item \textbf{False.} Peer-to-peer is suitable only for small networks (roughly up to about ten computers). With 200 computers, security, backups and administration become unmanageable; a \cle{client/server} architecture is required.
\item \textbf{True.} In a P2P network there is no central administrator, so each user must configure the sharing rights, passwords and antivirus protection of his or her own computer.
\end{enumerate}
\end{corrige}

\begin{corrige}[Exercise 3 — Definitions]
\begin{enumerate}
\item \textbf{\cle{Network architecture}:} the way in which the computers of a network are organised and the roles (\cle{client}, \cle{server}, or both) that each machine plays in the sharing of resources and services.
\item \textbf{\cle{Peer-to-peer network}:} a network architecture in which all computers are equal; there is no dedicated server, and each computer can act both as a \cle{client} (requesting resources) and as a \cle{server} (offering its own resources).
\item \textbf{\cle{Client/server network}:} a network architecture in which one or more powerful \cle{dedicated servers} provide services (files, printing, user accounts, backups) to \cle{client workstations} that send requests to them.
\item \textbf{\cle{Workstation}:} a client computer used by an end user to send requests to the server and to use the shared resources of the network.
\item \textbf{\cle{Dedicated server}:} a powerful computer reserved entirely for providing network services (file storage, authentication, printing, backups); it is not used as an ordinary user machine.
\end{enumerate}
\end{corrige}

\begin{corrige}[Exercise 4 — Quick questions]
\begin{enumerate}
\item \textbf{Peer-to-peer network.}
\begin{itemize}
\item \emph{Advantages (any two):} cheap to set up (no dedicated server, no server licence); easy to install and configure; no need for a specialised network administrator.
\item \emph{Disadvantages (any two):} security and backups cannot be centralised; difficult to manage when the number of computers grows (poor \cle{scalability}); each machine must be switched on for its resources to be available.
\end{itemize}
\item \textbf{Client/server network.}
\begin{itemize}
\item \emph{Advantages (any two):} centralised security (accounts, passwords, access rights managed on the server); centralised and automatic backups; good scalability (easy to add many workstations); resources available even when user machines are off.
\item \emph{Disadvantages (any two):} high cost (server hardware, server operating system, administrator); the server is a \cle{single point of failure} — if it breaks down, network services stop; installation and administration are complex.
\end{itemize}
\item \textbf{Cameroonian examples.} P2P is appropriate for a small business centre or a home in Bamenda with 3--5 computers sharing a printer. Client/server is appropriate for a bank in Douala, a ministry in Yaoundé or a university, where accounts, data security and central backups are essential.
\item \textbf{All 5 computers must be checked.} In a peer-to-peer network there is no central storage: each computer keeps its own users' files locally. A complete backup therefore requires copying the data of each of the $5$ machines separately, i.e.\ $5$ backup operations. This illustrates precisely why backups are difficult in P2P.
\end{enumerate}
\end{corrige}

\begin{corrige}[Exercise 5 — Completing a comparison table]
\begin{center}
\begin{tabularx}{\linewidth}{|l|X|X|}
\hline
\rowcolor{popblueL} \textbf{Criterion} & \textbf{Peer-to-peer} & \textbf{Client/server} \\
\hline
Cost of installation & Low: ordinary computers, no dedicated server, no server licence & High: powerful server, server OS licences, administrator \\
\hline
Need for a network administrator & No: each user manages his/her own machine & Yes: a trained administrator manages the server \\
\hline
Centralised security & No: rights and passwords set on each machine & Yes: accounts, passwords and access rights managed on the server \\
\hline
Centralised backup & No: each computer must be backed up separately & Yes: automatic backup of all data on the server \\
\hline
Scalability (adding many computers) & Poor: becomes unmanageable beyond about 10 machines & Good: hundreds of workstations can be added \\
\hline
Effect of one machine failing & Only the resources of that machine become unavailable & If a workstation fails, little impact; if the server fails, all services stop \\
\hline
\end{tabularx}
\end{center}
\textbf{Choice for a school laboratory of 25 computers: client/server.} Justification (any two reasons):
\begin{itemize}
\item 25 computers exceed the practical limit of a peer-to-peer network; a server allows centralised management of all student accounts, so each student can log on from any machine.
\item Security and access rights (e.g.\ preventing students from deleting system files) are controlled centrally, and students' work can be backed up automatically on the server.
\end{itemize}
\end{corrige}

\begin{corrige}[Exercise 6 — Identifying the architecture]
\begin{enumerate}
\item \textbf{Peer-to-peer.} Each laptop shares its own folders and there is no dedicated server; all machines are equal.
\item \textbf{Client/server.} All customer data and services are concentrated on one powerful central machine which answers the requests of the cashiers' client workstations.
\item \textbf{Peer-to-peer.} The computers share a printer and their folders mutually, with no central server; each machine is both client and server.
\item \textbf{Client/server.} A central computer manages user accounts, emails and backups for all staff machines, which is the defining role of a dedicated server.
\end{enumerate}
\end{corrige}

\begin{corrige}[Exercise 7 — Drawing a client/server network]
\textbf{(a) and (b) Diagram with the flow of a print request:}

\begin{popfigure}
\begin{tikzpicture}[scale=0.95, every node/.style={font=\small}]
% switch
\node[draw=popblue, fill=popblueL, rounded corners, minimum width=2.6cm, minimum height=0.9cm] (sw) at (0,0) {\textbf{Switch}};
% server
\node[draw=poporange, fill=popgold!25, rounded corners, minimum width=2.6cm, minimum height=0.9cm] (srv) at (-4.5,2.2) {\textbf{File server}};
% printer
\node[draw=popgreen, fill=popgreenL, rounded corners, minimum width=2.4cm, minimum height=0.9cm] (pr) at (4.5,2.2) {\textbf{Shared printer}};
% workstations
\node[draw=popdark, fill=white, rounded corners, minimum width=1.5cm, minimum height=0.7cm] (w1) at (-5.2,-2.4) {WS1};
\node[draw=popdark, fill=white, rounded corners, minimum width=1.5cm, minimum height=0.7cm] (w2) at (-2.6,-2.4) {WS2};
\node[draw=popdark, fill=white, rounded corners, minimum width=1.5cm, minimum height=0.7cm] (w3) at (0,-2.4) {WS3};
\node[draw=popdark, fill=white, rounded corners, minimum width=1.5cm, minimum height=0.7cm] (w4) at (2.6,-2.4) {WS4};
\node[draw=popdark, fill=white, rounded corners, minimum width=1.5cm, minimum height=0.7cm] (w5) at (5.2,-2.4) {WS5};
% links
\draw[thick, popline] (srv) -- (sw);
\draw[thick, popline] (pr) -- (sw);
\draw[thick, popline] (w1) -- (sw);
\draw[thick, popline] (w2) -- (sw);
\draw[thick, popline] (w3) -- (sw);
\draw[thick, popline] (w4) -- (sw);
\draw[thick, popline] (w5) -- (sw);
% print request arrows
\draw[->, very thick, poppink] (w2.north) .. controls (-2.6, -1.0) and (-1.0, -0.5) .. (sw.south) node[midway, left, font=\scriptsize, text=poppink, align=center]{1. request};
\draw[->, very thick, poppink] (sw.east) .. controls (3.0, 0.0) and (4.0, 1.0) .. (pr.south) node[midway, below right, font=\scriptsize, text=poppink, align=center]{2. print job};
\end{tikzpicture}
\legende{Star topology of the microfinance network; arrows show a print request from workstation WS2 to the shared printer via the switch.}
\end{popfigure}

\textbf{(c) Role of each component:}
\begin{itemize}
\item \textbf{File server:} stores the customers' data and statements, manages teller accounts and access rights, and may manage the print queue.
\item \textbf{Switch:} central device that receives frames from one machine and forwards them only to the correct destination (server, printer or workstation).
\item \textbf{Shared printer:} receives print jobs from the workstations (through the network) and prints the customers' statements.
\item \textbf{Teller workstations (WS1--WS5):} client computers used by the tellers to send requests (open a file, print a statement) to the server and the printer.
\end{itemize}
\end{corrige}

\begin{corrige}[Exercise 8 — Choosing an architecture]
\begin{enumerate}
\item \textbf{Home network (2 computers, 1 printer): peer-to-peer.} Reasons: (i) it is very cheap — no dedicated server or licence is needed; (ii) with only 2 machines, administration is trivial and centralised security is unnecessary.
\item \textbf{Hospital (60 computers, sensitive patient records): client/server.} Reasons: (i) patient data must be protected by centralised accounts, passwords and access rights on a server; (ii) automatic nightly backups can be performed centrally on the server, and 60 machines are far beyond the capacity of a P2P network.
\item \textbf{Small business centre (6 computers, very limited budget): peer-to-peer.} Reasons: (i) no expensive server or administrator is required, which fits the limited budget; (ii) 6 computers are few enough to be managed machine by machine.
\item \textbf{School laboratory (300 students, 40 computers, personal logins): client/server.} Reasons: (i) only a server can store 300 personal accounts centrally so that each student logs on from \emph{any} of the 40 computers; (ii) security, disk quotas and backups of students' work are managed centrally, which is impossible in P2P.
\end{enumerate}
\end{corrige}

\begin{corrige}[Exercise 9 — Structured question: advising a cybercafé owner]
\textbf{Indicative mark scheme (10 marks).}
\begin{enumerate}
\item \textbf{Definition of network architecture (2 marks).} The network architecture is the way the computers of a network are organised and the roles (\cle{client}, \cle{server}, or both) that each machine plays in sharing resources and services. \emph{(1 mark for ``organisation/roles of machines'', 1 mark for mention of resource/service sharing.)}
\item \textbf{The two architectures (2 marks).} \cle{Peer-to-peer} architecture \emph{(1 mark)} and \cle{client/server} architecture \emph{(1 mark)}.
\item \textbf{Advice: peer-to-peer (4 marks — 1 mark per justified reason).}
\begin{itemize}
\item His budget is small: P2P needs no expensive dedicated server nor server licence.
\item With only 8 computers, the network stays below the practical P2P limit (about 10 machines), so it remains manageable.
\item No specialised network administrator is needed: Mr. Tabi can configure sharing himself, reducing running costs.
\item Sharing the printer and files is easy to set up directly between the machines (and the Internet connection can be shared through the router/modem), which covers all his stated needs.
\end{itemize}
\item \textbf{Two limitations and remedies (2 marks — 1 mark each).}
\begin{itemize}
\item \emph{No centralised security:} each computer must be protected separately. Remedy: install an antivirus and configure passwords and sharing rights carefully on every machine.
\item \emph{Backups are not centralised:} data on each PC must be saved separately. Remedy: schedule regular manual backups of each computer onto an external hard disk.
\end{itemize}
(Also acceptable: a machine must be switched on for its resources to be available — remedy: keep the machine sharing the printer/files permanently on.)
\end{enumerate}
\emph{Examiner's expectations:} precise definitions, reasons explicitly linked to Mr. Tabi's situation (budget, 8 computers), and realistic remedies — not just a list of disadvantages copied from the course.
\end{corrige}

\begin{corrige}[Exercise 10 — Case study: connecting a ministry]
\textbf{Indicative mark scheme (13 marks).}
\begin{enumerate}
\item \textbf{Recommended architecture: client/server (4 marks — 1 mark per requirement linked).}
\begin{itemize}
\item \emph{Centralised accounts and passwords:} the server manages authentication, so each employee logs on from any workstation.
\item \emph{Automatic daily backups:} all data are stored on the server, where backup software runs automatically every day.
\item \emph{Controlled access per department:} the administrator defines access rights and groups (e.g.\ Finance, HR, Registry) on the server.
\item \emph{Future growth:} client/server is scalable — new workstations are simply connected and given accounts, unlike P2P which becomes unmanageable with 50+ machines.
\end{itemize}
\item \textbf{Hardware (4 marks — 1 mark per element with its role).}
\begin{itemize}
\item \emph{Server:} powerful dedicated machine storing data and providing services to the 50 clients.
\item \emph{Workstations:} the 50 employees' client computers sending requests to the server.
\item \emph{Switch(es):} connect all machines and forward frames to the correct destination; one switch per department, interconnected, is appropriate.
\item \emph{Cabling (e.g.\ twisted pair / fibre) or wireless access points:} the transmission medium linking workstations to the switches and the server.
\end{itemize}
\item \textbf{Software roles on the server (3 marks — 1 mark each, any three).}
\begin{itemize}
\item \emph{User account management:} creation of accounts, passwords and department groups (authentication/authorisation service).
\item \emph{File services:} shared folders with access rights per department.
\item \emph{Backup software:} scheduled automatic daily backups of all data.
\item \emph{Security software:} server antivirus, firewall and update management.
\end{itemize}
\item \textbf{Server failure (2 marks).} If the single server fails, all centralised services stop: employees can no longer log on, open shared files or access backups — the server is a \cle{single point of failure} \emph{(1 mark)}. \emph{Measure (1 mark, any one):} install a second (backup/redundant) server that takes over; perform regular backups to an external disk or another site; use a UPS to protect against power cuts.
\end{enumerate}
\end{corrige}

\begin{corrige}[Exercise 11 — Essay: discussing a statement]
\textbf{Indicative mark scheme (13 marks) and model outline.}
\begin{enumerate}
\item \textbf{How client/server works and why it suits large organisations (4 marks).} In a \cle{client/server} network, powerful \cle{dedicated servers} provide services (file storage, user accounts, printing, backups) to \cle{client workstations} that send requests to them \emph{(2 marks)}. Large organisations have many users and machines, sensitive data and a need for permanent availability: only a centralised architecture can manage hundreds of accounts, protect data and keep services running continuously \emph{(2 marks)}.
\item \textbf{Advantages (4 marks).}
\begin{itemize}
\item \emph{Security:} accounts, passwords and access rights are managed centrally on the server; access to files can be controlled per user or per department.
\item \emph{Scalability:} new workstations and users are added easily without reorganising the whole network.
\item \emph{Centralised administration:} one administrator manages accounts, software updates and automatic backups from the server, saving time and reducing errors.
\end{itemize}
\item \textbf{Limitations (3 marks).}
\begin{itemize}
\item \emph{Cost:} powerful server hardware, server operating system licences and a paid administrator make the initial investment high.
\item \emph{Dependence on the server:} the server is a \cle{single point of failure} — if it breaks down, all network services stop (hence the need for a backup server, a UPS, etc.).
\item \emph{Complexity:} installation and administration require specialised skills.
\end{itemize}
\item \textbf{Reasoned conclusion (2 marks).} The statement is largely true for large organisations, because security, scalability and centralised administration outweigh the cost. However, ``always'' is too strong: for very small structures a \cle{peer-to-peer} network remains preferable because it is cheap and sufficient — for example, a small business centre in Kumba or a home in Bamenda with a few computers sharing a printer does not need a server. The best choice therefore depends on the \emph{size, budget and security needs} of the organisation.
\end{enumerate}
\emph{Examiner's expectations:} a structured essay (introduction, developed arguments, conclusion), correct technical vocabulary (\cle{server}, \cle{workstation}, \cle{single point of failure}, \cle{scalability}), a balanced discussion of both sides, and a conclusion supported by a concrete Cameroonian example.
\end{corrige}

\newpage

\coursec{Summary sheet}

\begin{popfigure}
\begin{center}\tcbox[colback=popdark,colframe=popdark,coltext=white,arc=9pt,boxsep=4pt,left=6pt,right=6pt]{\bfseries\small Network Architecture}\end{center}
\vspace{-2pt}
\begin{tcbraster}[raster columns=3,raster equal height=rows,raster column skip=6pt,raster row skip=6pt,enhanced,blankest]
\begin{tcolorbox}[enhanced,colback=white,colframe=popblue,boxrule=0.9pt,arc=3pt,title={Definition \& Goals},fonttitle=\bfseries\scriptsize,coltitle=white,colbacktitle=popblue,left=4pt,right=4pt,top=2pt,bottom=2pt,boxsep=1pt,toptitle=1.5pt,bottomtitle=1.5pt,halign=left]
\begin{itemize}[nosep,leftmargin=*,itemsep=1pt,font=\scriptsize]
\item Sharing resources
\item Communication
\end{itemize}
\end{tcolorbox}
\begin{tcolorbox}[enhanced,colback=white,colframe=popgreen,boxrule=0.9pt,arc=3pt,title={Peer-to-Peer (P2P)},fonttitle=\bfseries\scriptsize,coltitle=white,colbacktitle=popgreen,left=4pt,right=4pt,top=2pt,bottom=2pt,boxsep=1pt,toptitle=1.5pt,bottomtitle=1.5pt,halign=left]
\begin{itemize}[nosep,leftmargin=*,itemsep=1pt,font=\scriptsize]
\item All PCs equal
\item No dedicated server
\end{itemize}
\end{tcolorbox}
\begin{tcolorbox}[enhanced,colback=white,colframe=poporange,boxrule=0.9pt,arc=3pt,title={Client/Server},fonttitle=\bfseries\scriptsize,coltitle=white,colbacktitle=poporange,left=4pt,right=4pt,top=2pt,bottom=2pt,boxsep=1pt,toptitle=1.5pt,bottomtitle=1.5pt,halign=left]
\begin{itemize}[nosep,leftmargin=*,itemsep=1pt,font=\scriptsize]
\item Central server
\item Clients request
\end{itemize}
\end{tcolorbox}
\begin{tcolorbox}[enhanced,colback=white,colframe=poppurple,boxrule=0.9pt,arc=3pt,title={Comparison Criteria},fonttitle=\bfseries\scriptsize,coltitle=white,colbacktitle=poppurple,left=4pt,right=4pt,top=2pt,bottom=2pt,boxsep=1pt,toptitle=1.5pt,bottomtitle=1.5pt,halign=left]
\begin{itemize}[nosep,leftmargin=*,itemsep=1pt,font=\scriptsize]
\item Cost, security
\item Size, admin
\end{itemize}
\end{tcolorbox}
\begin{tcolorbox}[enhanced,colback=white,colframe=poppink,boxrule=0.9pt,arc=3pt,title={Choosing a Strategy},fonttitle=\bfseries\scriptsize,coltitle=white,colbacktitle=poppink,left=4pt,right=4pt,top=2pt,bottom=2pt,boxsep=1pt,toptitle=1.5pt,bottomtitle=1.5pt,halign=left]
\begin{itemize}[nosep,leftmargin=*,itemsep=1pt,font=\scriptsize]
\item Needs analysis
\item Scalability
\end{itemize}
\end{tcolorbox}
\begin{tcolorbox}[enhanced,colback=white,colframe=popgold,boxrule=0.9pt,arc=3pt,title={Cameroon Context},fonttitle=\bfseries\scriptsize,coltitle=white,colbacktitle=popgold,left=4pt,right=4pt,top=2pt,bottom=2pt,boxsep=1pt,toptitle=1.5pt,bottomtitle=1.5pt,halign=left]
\begin{itemize}[nosep,leftmargin=*,itemsep=1pt,font=\scriptsize]
\item Schools, cybercaf\'es
\item Banks, e-gov
\end{itemize}
\end{tcolorbox}
\end{tcbraster}

\legende{Mind map of the chapter: Network Design Strategies (network architecture).}
\end{popfigure}

\begin{bilan}
\textbf{Key definitions}
\begin{itemize}
\item \cle{Network architecture}: the overall design or structural model that defines how computers in a network are organised, how they communicate, and how resources (files, printers, Internet) are shared and managed.
\item \cle{Peer-to-peer (P2P) network}: an architecture in which \textbf{all computers are equal} (peers); each machine can act both as a client (requesting resources) and as a server (providing resources). There is \textbf{no dedicated central server}.
\item \cle{Client/server network}: an architecture built around one or more \textbf{dedicated, powerful servers} that provide services (files, printing, email, authentication, databases) to \textbf{client} machines that request them.
\item \cle{Server}: a powerful computer (hardware + server software, e.g.\ Windows Server, Linux) that stores shared data centrally and manages access to it.
\item \cle{Client}: a workstation that sends requests to the server and uses the services offered.
\item \cle{Network administrator}: the person who manages user accounts, security, backups and resources on a client/server network.
\end{itemize}

\textbf{Peer-to-peer architecture — essentials}
\begin{itemize}
\item Typical size: small networks (about \textbf{2 to 10 computers}); easy to set up with ordinary operating systems (Windows Home editions).
\item \textbf{Advantages}: low cost (no server, no specialist administrator); simple installation and configuration; no single point of failure for the whole network.
\item \textbf{Disadvantages}: weak security (each user manages his own machine); no centralised backup; difficult to manage as the network grows; a peer that is switched off makes its shared resources unavailable; performance drops when a machine serves many requests.
\item Typical uses: home networks, small cybercaf\'es, a few computers in a small office in Bamenda or Bafoussam sharing one printer.
\end{itemize}

\textbf{Client/server architecture — essentials}
\begin{itemize}
\item Roles are \textbf{specialised}: servers provide services, clients consume them. Centralised storage, user accounts and security policies (e.g.\ a domain controller).
\item \textbf{Advantages}: strong centralised security and access control; centralised backup and data management; easy administration of many users; \textbf{scalable} (grows well to hundreds of machines); better performance under heavy load.
\item \textbf{Disadvantages}: high cost (server hardware, server licences, trained administrator); more complex to install and maintain; the server can be a \textbf{single point of failure} (reduced by backups and redundant servers).
\item Typical uses: schools and universities, banks, ministries and e-government services, MTN/Orange service platforms, hospital information systems.
\end{itemize}

\textbf{Comparison table (learn it by heart)}
\begin{center}
\begin{tabularx}{\linewidth}{|l|X|X|}
\hline
\rowcolor{popblueL} \textbf{Criterion} & \textbf{Peer-to-peer} & \textbf{Client/Server} \\
\hline
Cost & Low & High \\
\hline
Security & Weak, decentralised & Strong, centralised \\
\hline
Administration & Each user, no specialist & Central administrator \\
\hline
Recommended size & Small ($\leq 10$ PCs) & Medium to large \\
\hline
Data backup & On each machine & Centralised on server \\
\hline
Failure impact & Limited to one peer & Server failure stops services \\
\hline
\end{tabularx}
\end{center}

\textbf{Method: choosing a design strategy}
\begin{enumerate}
\item Analyse the \textbf{needs}: number of users, types of resources shared, expected growth.
\item Evaluate the \textbf{security requirements}: sensitive data (marks, money, medical records) demand client/server.
\item Evaluate the \textbf{budget} and the availability of a qualified administrator.
\item Check \textbf{scalability}: will the network grow beyond about ten machines?
\item Decide: small scale + low budget + low sensitivity $\Rightarrow$ \textbf{peer-to-peer}; otherwise $\Rightarrow$ \textbf{client/server}.
\end{enumerate}

\textbf{Common mistakes to avoid}
\begin{itemize}
\item Confusing \emph{architecture} (peer-to-peer vs client/server: \textbf{roles} of machines) with \emph{topology} (bus, star, ring: \textbf{physical layout}). They are independent: a star topology can host either architecture.
\item Saying ``P2P has no server at all'': each peer \emph{can} act as a server; there is simply no \emph{dedicated} server.
\item Believing client/server removes all risks: without backups and redundancy, server failure paralyses the network.
\item Recommending client/server for a 3-computer home network: unjustified cost and complexity.
\end{itemize}

\textbf{GCE tips}
\begin{itemize}
\item A very frequent question: ``\emph{State two advantages and two disadvantages of each architecture}'' — prepare exactly two of each, clearly worded.
\item In scenario questions (e.g.\ ``a school with 200 students''), always \textbf{justify} your choice with at least two criteria (size, security, cost, administration).
\item Define terms before comparing: examiners reward precise definitions of \cle{peer}, \cle{client}, \cle{server}.
\item Use a labelled diagram (clients around a central server) to earn presentation marks.
\end{itemize}
\end{bilan}

\findecours

\end{document}
